% Les révolutions scientifiques et techniques depuis la fin de la Deuxième Guerre mondiale Chapitre 6: La marche vers l’indépendance — Histoire Tle Tech — Cours signé OnBuch+
% Programme officiel MINESEC. Compiler avec : tectonic N08-revolutions-scientifiques-techniques-depuis-fin-deuxiem.tex
\newcommand{\DOCMATIERE}{Histoire}
\newcommand{\DOCNIVEAU}{Tle Tech}
\newcommand{\DOCMODULE}{Histoire – classe de Terminale technique}
\newcommand{\DOCLECON}{N08}
\newcommand{\DOCTITRE}{Les révolutions scientifiques et techniques depuis la fin de la Deuxième Guerre mondiale Chapitre 6: La marche vers l’indépendance}
% OnBuch+ — préambule « cours signés » ENRICHI (compilé avec tectonic / XeLaTeX)
% Système visuel « Pop » d'OnBuch+ : fond crème (jamais blanc), bordures encre
% épaisses, ombres « dures » décalées, titres Archivo Black en capitales.
% Version MATHÉMATIQUES : graphes (pgfplots/tikz), boîtes pédagogiques
% (définition, propriété, méthode, attention, le savais-tu, exercice résolu,
% exercice, TP, bilan), page de garde et signature OnBuch+ ancrée en bas.
%
% Macros attendues dans main.tex AVANT \input{preamble} :
%   \DOCMATIERE  \DOCNIVEAU  \DOCMODULE  \DOCLECON (numéro)  \DOCTITRE
\documentclass[11pt]{article}

\usepackage{fontspec}
\usepackage[french]{babel}
\usepackage[a4paper,margin=2cm,top=3.4cm,bottom=2.8cm,headheight=1.4cm,footskip=1.3cm]{geometry}
\usepackage{amsmath,amssymb,amsfonts}
\usepackage{mathtools} % \xmapsto, \coloneqq... souvent utilisés par les modèles
\usepackage{cancel} % simplifications barrées (bilans de forces)
% \abs{z} (module, valeur absolue) et \norm{u} (norme), utilisés par les modèles
\providecommand{\abs}{}\let\abs\relax\DeclarePairedDelimiter{\abs}{\lvert}{\rvert}
\providecommand{\norm}{}\let\norm\relax\DeclarePairedDelimiter{\norm}{\lVert}{\rVert}
\usepackage[table]{xcolor} % [table] : fournit aussi \rowcolors (alternance de lignes)
\usepackage{pagecolor}
\usepackage{tikz}
\usetikzlibrary{arrows.meta,positioning,calc,shapes.geometric,shapes.misc,shapes.arrows,decorations.pathmorphing,decorations.pathreplacing,decorations.markings,patterns,shadows,mindmap,trees,fit,backgrounds,matrix,intersections,angles,quotes}
\usepackage{pgfplots}
\usepackage[version=4]{mhchem} % équations nucléaires \ce{^{235}_{92}U}
\usepackage[european,siunitx]{circuitikz} % schémas électriques
\pgfplotsset{compat=1.18}
\usepgfplotslibrary{fillbetween,groupplots} % aires entre courbes (calcul intégral)
\usepackage{enumitem}
\usepackage{fancyhdr}
% [most] charge toutes les bibliothèques tcolorbox (listings, minted, raster,
% documentation...) et gonfle inutilement la mémoire TeX sur un document long
% (~300 boîtes) : on ne charge que ce qui est réellement utilisé.
\usepackage[skins,breakable]{tcolorbox}
\usepackage{graphicx}
\usepackage{colortbl}
\usepackage{booktabs}
\usepackage{tabularx}
\usepackage{multirow}
\usepackage{diagbox} % case d'en-tête diagonale des tableaux à double entrée
% Colonnes extensibles centrée / gauche / droite (C, L, R), souvent utilisées par les modèles
\newcolumntype{C}{>{\centering\arraybackslash}X}
\newcolumntype{L}{>{\raggedright\arraybackslash}X}
\newcolumntype{R}{>{\raggedleft\arraybackslash}X}
\usepackage{array}
\usepackage{multicol}
\usepackage{siunitx}
% parse-numbers=false : accepte aussi \SI{2,5 \times 10^{-3}}{...}
\sisetup{locale=FR,output-decimal-marker={,},group-minimum-digits=4,parse-numbers=false}
\DeclareSIUnit\ohm{\ensuremath{\symup{\Omega}}} % U+2126 absent de la police
\DeclareSIUnit\division{div} % réglages d'oscilloscope (ms/div, V/div)
\DeclareSIUnit\tr{tr} % tours (vitesse de rotation, ex. \SI{1500}{\tr\per\minute})
\DeclareSIUnit\molar{M} % concentration molaire (ex. \SI{3}{\molar}, \SI{1}{\micro\molar})
\DeclareSIUnit\an{an} % année (le modèle confond parfois avec \year, un compteur TeX) — ex. \SI{5}{\kilo\meter\per\an}
\DeclareSIUnit\FCFA{FCFA} % franc CFA (ex. \SI{500}{\FCFA\per\kilo\gram})
\DeclareSIUnit\degreeBrix{\textdegree Brix} % teneur en sucre d'un sirop/jus (réfractométrie)
\DeclareSIUnit\month{mois} % le modèle utilise parfois \month, un compteur TeX, comme unité de durée
\DeclareSIUnit\microliter{\micro\liter} % le modèle écrit parfois \microliter en un seul token
\DeclareSIUnit\million{millions} % dénombrement (ex. \SI{15}{\million\per\mL} spermatozoïdes)
\usepackage{qrcode}
% \degree : commande fréquemment utilisée par les modèles pour un angle ou une
% température en dehors de \SI{}{} (non fournie par les paquets chargés).
\providecommand{\degree}{^{\circ}}
% \qed (fin de démonstration), utilisé par les modèles sans amsthm
\providecommand{\qed}{\hfill$\square$}
% \barre : double barre des valeurs interdites dans un tableau de variations (tkz-tab)
\providecommand{\barre}{\ensuremath{\|}}
% environnement proof (amsthm non chargé) utilisé par les modèles
\ifdefined\proof\else\newenvironment{proof}[1][Démonstration]{\par\noindent\textit{#1.}\ }{\qed\par}\fi
\usepackage{lastpage}
\usepackage{hyperref}
\hypersetup{hidelinks}

% ============================================================
% Palette « Pop » OnBuch+ — mêmes hex que la classe P (plus_ui.dart)
% ============================================================
\definecolor{poporange}{HTML}{F59321}
\definecolor{popdark}{HTML}{E07A0C}
\definecolor{popcream}{HTML}{FFF6E8}
\definecolor{popink}{HTML}{1C1714}
\definecolor{poppurple}{HTML}{7A5AE0}
\definecolor{popgreen}{HTML}{2E9E6A}
\definecolor{poppink}{HTML}{E0457B}
\definecolor{popblue}{HTML}{2F7DE1}
\definecolor{popgold}{HTML}{C98A12}
\definecolor{popmuted}{HTML}{8A7F76}
\definecolor{popline}{HTML}{EADFD2}
\definecolor{popgray}{HTML}{8A7F76}
\colorlet{popgrey}{popgray}
% Alias tolérants (le modèle nomme parfois les couleurs différemment)
\colorlet{popwhite}{white}
\colorlet{popblack}{popink}
\colorlet{popred}{poppink}
\colorlet{popyellow}{popgold}
% Teintes claires (fonds de tableaux/encadrés) — le modèle nomme parfois
% les couleurs avec un suffixe L (ex. popblueL) sans les avoir définies.
\colorlet{poporangeL}{poporange!20}
\colorlet{popdarkL}{popdark!20}
\colorlet{poppurpleL}{poppurple!20}
\colorlet{popgreenL}{popgreen!20}
\colorlet{poppinkL}{poppink!20}
\colorlet{popblueL}{popblue!20}
\colorlet{popgoldL}{popgold!20}
\colorlet{popredL}{popred!20}
\colorlet{popgrayL}{popgray!20}
% Teintes claires pour fonds de tableaux / zones
\colorlet{popblueL}{popblue!10!white}
\colorlet{poporangeL}{poporange!14!white}
\colorlet{popgreenL}{popgreen!12!white}
\colorlet{poppurpleL}{poppurple!12!white}
\colorlet{poppinkL}{poppink!12!white}
\colorlet{popgoldL}{popgold!14!white}

\pagecolor{popcream}

% ============================================================
% Typographie
% ============================================================
\defaultfontfeatures{Ligatures=TeX}
\setmainfont{PlusJakartaSans-Medium.ttf}[Path=../../../fonts/,BoldFont=PlusJakartaSans-Bold.ttf]
% Maths en sans-serif (Fira Math) avec chiffres et lettres droites de Plus
% Jakarta, pour que les formules se fondent dans le texte.
\usepackage[math-style=ISO,bold-style=ISO]{unicode-math}
\setmathfont{FiraMath-Regular.otf}[Path=../../../fonts/]
\setmathfont{PlusJakartaSans-Medium.ttf}[Path=../../../fonts/,range={up/num,up/{Latin,latin},\mathup}]
\usepackage{icomma} % virgule décimale sans espace en mode math
% Caractères mathématiques Unicode absents des polices de texte (Plus Jakarta,
% Archivo Black) : rendus par leur équivalent mathématique, en texte comme en
% formule — sinon ils s'affichent en carré vide (ex. « Arithmétique dans ℤ »).
\usepackage{newunicodechar}
\newunicodechar{ℝ}{\ensuremath{\mathbb{R}}}
\newunicodechar{ℤ}{\ensuremath{\mathbb{Z}}}
\newunicodechar{ℂ}{\ensuremath{\mathbb{C}}}
\newunicodechar{ℕ}{\ensuremath{\mathbb{N}}}
\newunicodechar{ℚ}{\ensuremath{\mathbb{Q}}}
\newunicodechar{𝔻}{\ensuremath{\mathbb{D}}}
\newunicodechar{α}{\ensuremath{\alpha}}
\newunicodechar{β}{\ensuremath{\beta}}
\newunicodechar{γ}{\ensuremath{\gamma}}
\newunicodechar{δ}{\ensuremath{\delta}}
\newunicodechar{ε}{\ensuremath{\varepsilon}}
\newunicodechar{θ}{\ensuremath{\theta}}
\newunicodechar{λ}{\ensuremath{\lambda}}
\newunicodechar{μ}{\ensuremath{\mu}}
\newunicodechar{σ}{\ensuremath{\sigma}}
\newunicodechar{τ}{\ensuremath{\tau}}
\newunicodechar{φ}{\ensuremath{\varphi}}
\newunicodechar{ψ}{\ensuremath{\psi}}
\newunicodechar{ω}{\ensuremath{\omega}}
\newunicodechar{Σ}{\ensuremath{\Sigma}}
% majuscules produites par \MakeUppercase dans les titres (α-aminés -> Α)
\newunicodechar{Α}{\ensuremath{\alpha}}
\newunicodechar{Β}{\ensuremath{\beta}}
\newunicodechar{Γ}{\ensuremath{\Gamma}}
\newunicodechar{Δ}{\ensuremath{\Delta}}
\newunicodechar{Ε}{\ensuremath{\varepsilon}}
\newunicodechar{Θ}{\ensuremath{\Theta}}
\newunicodechar{Λ}{\ensuremath{\Lambda}}
\newunicodechar{Μ}{\ensuremath{\mu}}
\newunicodechar{Τ}{\ensuremath{\tau}}
\newunicodechar{Φ}{\ensuremath{\Phi}}
\newunicodechar{Ψ}{\ensuremath{\Psi}}
\newunicodechar{Ω}{\ensuremath{\Omega}}
\newunicodechar{↦}{\ensuremath{\mapsto}}
\newunicodechar{∈}{\ensuremath{\in}}
\newunicodechar{∉}{\ensuremath{\notin}}
\newunicodechar{∧}{\ensuremath{\wedge}}
\newunicodechar{⇔}{\ensuremath{\Leftrightarrow}}
\newunicodechar{⟺}{\ensuremath{\Longleftrightarrow}}
\newunicodechar{⇒}{\ensuremath{\Rightarrow}}
\newunicodechar{⟹}{\ensuremath{\Longrightarrow}}
\newunicodechar{∘}{\ensuremath{\circ}}
\newunicodechar{∪}{\ensuremath{\cup}}
\newunicodechar{∩}{\ensuremath{\cap}}
\newunicodechar{≡}{\ensuremath{\equiv}}
\newunicodechar{⊂}{\ensuremath{\subset}}
\newunicodechar{⊥}{\ensuremath{\perp}}
\newunicodechar{∃}{\ensuremath{\exists}}
\newunicodechar{∀}{\ensuremath{\forall}}
\newunicodechar{∛}{\ensuremath{\sqrt[3]{\,}}}
\newunicodechar{∜}{\ensuremath{\sqrt[4]{\,}}}
\newunicodechar{⃗}{\ensuremath{{}^{\to}}}
\newunicodechar{ⁿ}{\ifmmode^{n}\else\textsuperscript{\ensuremath{n}}\fi}
\newunicodechar{ˣ}{\ifmmode^{x}\else\textsuperscript{\ensuremath{x}}\fi}
\newunicodechar{ᵇ}{\ifmmode^{b}\else\textsuperscript{\ensuremath{b}}\fi}
\newunicodechar{ʳ}{\ifmmode^{r}\else\textsuperscript{\ensuremath{r}}\fi}
\newunicodechar{ᵘ}{\ifmmode^{u}\else\textsuperscript{\ensuremath{u}}\fi}
\newunicodechar{ʸ}{\ifmmode^{y}\else\textsuperscript{\ensuremath{y}}\fi}
\newunicodechar{ᵃ}{\ifmmode^{a}\else\textsuperscript{\ensuremath{a}}\fi}
\newunicodechar{ᵉ}{\ifmmode^{e}\else\textsuperscript{\ensuremath{e}}\fi}
\newunicodechar{ᵏ}{\ifmmode^{k}\else\textsuperscript{\ensuremath{k}}\fi}
\newunicodechar{ᵖ}{\ifmmode^{p}\else\textsuperscript{\ensuremath{p}}\fi}
\newunicodechar{ᴺ}{\ifmmode^{N}\else\textsuperscript{\ensuremath{N}}\fi}
\newunicodechar{ᶜ}{\ifmmode^{c}\else\textsuperscript{\ensuremath{c}}\fi}
\newunicodechar{ʰ}{\ifmmode^{h}\else\textsuperscript{\ensuremath{h}}\fi}
\newunicodechar{ᵠ}{\ifmmode^{q}\else\textsuperscript{\ensuremath{q}}\fi}
\newunicodechar{ₙ}{\ifmmode_{n}\else\textsubscript{\ensuremath{n}}\fi}
\newunicodechar{ₐ}{\ifmmode_{a}\else\textsubscript{\ensuremath{a}}\fi}
\newunicodechar{ᵢ}{\ifmmode_{i}\else\textsubscript{\ensuremath{i}}\fi}
\newunicodechar{ᵣ}{\ifmmode_{r}\else\textsubscript{\ensuremath{r}}\fi}
\newunicodechar{ₖ}{\ifmmode_{k}\else\textsubscript{\ensuremath{k}}\fi}
\newunicodechar{ᵦ}{\ifmmode_{\beta}\else\textsubscript{\ensuremath{\beta}}\fi}
\newunicodechar{ₓ}{\ifmmode_{x}\else\textsubscript{\ensuremath{x}}\fi}
\newfontfamily\popdisplay{ArchivoBlack-Regular.ttf}[Path=../../../fonts/]
\newfontfamily\poplogo{SpaceGrotesk-Bold.ttf}[Path=../../../fonts/]
\newfontfamily\popsemi{PlusJakartaSans-SemiBold.ttf}[Path=../../../fonts/]
\newfontfamily\popxbold{PlusJakartaSans-ExtraBold.ttf}[Path=../../../fonts/]
\newfontfamily\popxboldls{PlusJakartaSans-ExtraBold.ttf}[Path=../../../fonts/,LetterSpace=3.0]

\setlist[enumerate]{leftmargin=1.7em,itemsep=5pt,topsep=4pt}
\setlist[itemize]{leftmargin=1.7em,itemsep=4pt,topsep=4pt,label={\color{poporange}\textbullet}}
\setlength{\parindent}{0pt}
\setlength{\parskip}{5pt}
\renewcommand{\arraystretch}{1.25}

% pgfplots : style Pop par défaut
\pgfplotsset{
  every axis/.append style={axis line style={popink,line width=1pt},tick style={popink},
    grid style={popline,line width=0.5pt},label style={font=\small},tick label style={font=\footnotesize}},
  popcurve/.style={poporange,line width=1.8pt,no markers,smooth},
}
% Même style disponible dans un \draw TikZ simple (hors pgfplots)
\tikzset{popcurve/.style={poporange,line width=1.8pt}}
\tikzset{popgrid/.style={popline,very thin}}
\colorlet{popcurve}{poporange} % le nom du style est parfois utilisé comme couleur

% ============================================================
% Pill / squiggle
% ============================================================
\newcommand{\poppill}[3]{%
  \tikz[baseline=-0.55ex]\node[draw=popink,line width=1.2pt,rounded corners=2.6pt,%
    fill=#2,inner xsep=6.5pt,inner ysep=3pt]{%
    \popxboldls\fontsize{7}{7}\selectfont\color{#3}\MakeUppercase{#1}};%
}
\newcommand{\popsquiggle}[1]{%
  \tikz[baseline=-0.4ex]{
    \draw[#1,line width=2.6pt,line cap=round]
      plot [smooth] coordinates {(0,0.09) (0.34,0.01) (0.68,0.21) (1.02,0.01) (1.36,0.10)};
  }%
}
% Mot-clé surligné (marqueur orange sous le texte)
\newcommand{\cle}[1]{{\popxbold\color{popdark}#1}}

% ============================================================
% En-tête / pied
% ============================================================
\pagestyle{fancy}
\fancyhf{}
\renewcommand{\headrulewidth}{0pt}
\renewcommand{\footrulewidth}{0pt}
\lhead{\raisebox{-0.15ex}{\poplogo\fontsize{13}{13}\selectfont\color{popink}OnBuch\color{poporange}+}}
\rhead{\poppill{\DOCMATIERE\ \textbullet\ \DOCNIVEAU\ \textbullet\ Leçon \DOCLECON}{white}{popink}}
\lfoot{\popxbold\fontsize{7.6}{7.6}\selectfont\color{popmuted}\MakeUppercase{Cours signé OnBuch+}\ \textbullet\ {\popsemi\DOCTITRE}}
\rfoot{\tikz[baseline=-0.6ex]\node[rounded corners=8.5pt,fill=popink,minimum height=17pt,minimum width=17pt,inner xsep=5pt,inner sep=0]{\popxbold\fontsize{8}{8}\selectfont\color{popcream}\thepage\,/\,\pageref*{LastPage}};}

% ============================================================
% Titres
% ============================================================
\newcommand{\chaptitre}[2]{%
  \begin{tcolorbox}[enhanced,colback=poporange,colframe=popink,boxrule=2.6pt,arc=11pt,
      drop shadow=popink,left=15pt,right=15pt,top=12pt,bottom=11pt,before skip=3pt,after skip=18pt]
    \poppill{#2}{popink}{popcream}\par\vspace{9pt}
    {\raggedright\hyphenpenalty=10000\exhyphenpenalty=10000\popdisplay\fontsize{22}{24}\selectfont\color{popcream}\MakeUppercase{#1}\par}
  \end{tcolorbox}
}
% Section numérotée : pastille orange avec le numéro + titre Archivo
\newcounter{popsec}
\newcommand{\coursec}[1]{%
  \stepcounter{popsec}\setcounter{popsub}{0}%
  \par\vspace{16pt}\noindent
  \begin{minipage}{\linewidth}
  \tikz[baseline=-0.7ex]\node[draw=popink,line width=1.6pt,fill=poporange,rounded corners=4pt,minimum size=22pt,inner sep=0]{\popdisplay\fontsize{12}{12}\selectfont\color{popcream}\thepopsec};%
  \hspace{8pt}{\popdisplay\fontsize{15}{17}\selectfont\color{popink}\MakeUppercase{#1}}\par
  \vspace{4pt}\noindent{\color{poporange}\rule{36pt}{3.6pt}}
  \end{minipage}\par\nopagebreak\vspace{8pt}
}
\newcounter{popsub}
\newcommand{\courssub}[1]{%
  \stepcounter{popsub}%
  \par\vspace{9pt}\noindent{\popxbold\fontsize{11.5}{13}\selectfont\color{popink}{\color{poporange}\thepopsec.\thepopsub}\hspace{6pt}#1}\par\nopagebreak\vspace{4pt}
}

% ============================================================
% Boîtes pédagogiques — même recette : fond blanc, bordure épaisse colorée,
% ombre dure encre, badge plein. Titre optionnel [#1].
% ============================================================
\tcbset{popbase/.style={breakable,enhanced,colback=white,boxrule=2.3pt,arc=9pt,
  shadow={3pt}{-3pt}{0pt}{popink},left=14pt,right=14pt,top=11pt,bottom=11pt,before skip=11pt,after skip=15pt,
  fontupper=\popsemi\fontsize{10.6}{14.5}\selectfont\color{popink},
  fontlower=\popsemi\fontsize{10.6}{14.5}\selectfont\color{popink}}}
% Coche dessinée (le glyphe ✓ n'existe pas dans les polices du document)
\newcommand{\popcheck}[1]{\tikz[baseline=-0.5ex]\draw[#1,line width=1.6pt,line cap=round,line join=round](-0.13,0.0)--(-0.04,-0.09)--(0.13,0.1);}
\providecommand{\coche}{\popcheck{popgreen}}
\providecommand{\resultat}[1]{\par\smallskip\noindent\fcolorbox{popgreen}{popgreenL}{\parbox{\dimexpr\linewidth-2\fboxsep-2\fboxrule\relax}{\textbf{Résultat.} #1}}\par\smallskip}
\newcommand{\popboxhead}[3]{\poppill{#1}{#2}{white}\ifx\relax#3\relax\else\hspace{6pt}{\popxbold\fontsize{10.6}{12}\selectfont\color{popink}#3}\fi\par\vspace{7pt}}

\newtcolorbox{aretenir}[1][]{popbase,colframe=popgreen,before upper={\popboxhead{À retenir}{popgreen}{#1}}}
\newtcolorbox{exemplebox}[1][]{popbase,colframe=poporange,before upper={\popboxhead{Exemple}{poporange}{#1}}}
\newtcolorbox{definition}[1][]{popbase,colframe=popblue,before upper={\popboxhead{Définition}{popblue}{#1}}}
\newtcolorbox{propriete}[1][]{popbase,colframe=poppurple,before upper={\popboxhead{Propriété}{poppurple}{#1}}}
\newtcolorbox{methode}[1][]{popbase,colframe=popgold,before upper={\popboxhead{Méthode}{popgold}{#1}}}
\newtcolorbox{attention}[1][]{popbase,colframe=poppink,before upper={\popboxhead{Attention}{poppink}{#1}}}
\newtcolorbox{experience}[1][]{popbase,colframe=popblue,colback=popblueL,before upper={\popboxhead{Expérience}{popblue}{#1}}}
% « Le savais-tu ? » : carte violette pleine (contexte, culture, Cameroun)
\newtcolorbox{savaistu}[1][]{popbase,colback=poppurple,colframe=popink,
  fontupper=\popsemi\fontsize{10.4}{14.2}\selectfont\color{white},
  before upper={\poppill{Le savais-tu\,?}{white}{poppurple}\ifx\relax#1\relax\else\hspace{6pt}{\popxbold\color{white}#1}\fi\par\vspace{7pt}}}
% Exercice résolu : énoncé en haut, \tcblower puis la solution sur fond crème
\newcounter{popexo}\newcounter{popexr}
\newtcolorbox{exoresolu}[1][]{popbase,colframe=poporange,colbacklower=popcream,
  segmentation style={popink,line width=1.2pt,dashed},
  before upper={\stepcounter{popexr}\popboxhead{Exercice résolu \thepopexr}{poporange}{#1}},
  before lower={\poppill{Solution}{popink}{popcream}\par\vspace{6pt}}}
% Exercice (non résolu en place) — la correction est dans la partie Corrigés
\newtcolorbox{exercice}[1][]{popbase,colframe=popink,
  before upper={\stepcounter{popexo}\popboxhead{Exercice \thepopexo}{popink}{#1}}}
% Corrigé
\newtcolorbox{corrige}[1][]{popbase,colframe=popgreen,colback=popgreenL,
  before upper={\popboxhead{Corrigé}{popgreen}{#1}}}
% Objectifs / prérequis / bilan
\newtcolorbox{objectifs}{popbase,colframe=popink,colback=poporangeL,
  before upper={\popboxhead{Objectifs de la leçon}{popink}{}}}
\newtcolorbox{prerequis}{popbase,colframe=popink,colback=popblueL,
  before upper={\popboxhead{Prérequis}{popblue}{}}}
\newtcolorbox{bilan}{popbase,colframe=popink,colback=white,boxrule=2.8pt,
  before upper={\popboxhead{Fiche bilan}{poporange}{L'essentiel à connaître pour l'examen}}}
% Figure encadrée avec légende
\newtcolorbox{popfigure}[1][]{enhanced,breakable=false,colback=white,colframe=popline,boxrule=1.4pt,arc=8pt,
  left=10pt,right=10pt,top=10pt,bottom=8pt,before skip=10pt,after skip=12pt,halign=center,
  lower separated=false,halign lower=center,
  fontlower=\popsemi\fontsize{9}{11.5}\selectfont\color{popmuted},
  #1}
\newcommand{\legende}[1]{\par\vspace{6pt}{\popsemi\fontsize{9}{11.5}\selectfont\color{popmuted}#1}}

% ============================================================
% Signature OnBuch+
% ============================================================
\newcommand{\poplogoinline}[1]{{\poplogo\fontsize{#1}{#1}\selectfont\color{popink}OnBuch\color{poporange}+}}
\newcommand{\signatureonbuch}{%
  \begin{tcolorbox}[enhanced,colback=popink,colframe=popink,boxrule=2.2pt,arc=10pt,
      drop shadow=poporange,left=14pt,right=14pt,top=10pt,bottom=10pt,before skip=0pt,after skip=0pt]
    \tikz[baseline=-0.7ex]\node[circle,fill=popgreen,draw=popcream,line width=1.3pt,minimum size=22pt,inner sep=0]{\popcheck{white}};%
    \hspace{9pt}\begin{minipage}[c]{0.62\linewidth}
      {\popxbold\fontsize{12}{13}\selectfont\color{popcream}Signé\ \textbullet\ {\poplogo OnBuch\color{poporange}+}}\par\vspace{2pt}
      {\raggedright\popsemi\fontsize{8}{10.5}\selectfont\color{popline}\MakeUppercase{Équipe pédagogique OnBuch+}\\\MakeUppercase{Conforme au programme officiel MINESEC}\par}
    \end{minipage}\hfill
    {\poplogo\fontsize{22}{22}\selectfont\color{popcream}OnBuch\color{poporange}+}
  \end{tcolorbox}%
}

% ============================================================
% Page de garde
% #1 titre · #2 matière · #3 niveau · #4 module · #5 n° leçon · #6 durée ·
% #7 description / objectifs (texte ou itemize)
% ============================================================
\newcommand{\pagegarde}[7]{%
  \thispagestyle{empty}
  \newgeometry{a4paper,margin=1.7cm,top=1.6cm,bottom=1.4cm}%
  \begin{tikzpicture}[remember picture,overlay]
    \fill[poporange!18] ([xshift=-3.2cm,yshift=-3.6cm]current page.north east) circle (4.6cm);
    \fill[poppurple!14] ([xshift=2.4cm,yshift=7.6cm]current page.south west) circle (3.8cm);
  \end{tikzpicture}%
  {\poplogo\fontsize{30}{30}\selectfont\color{popink}OnBuch\color{poporange}+}\hfill
  \raisebox{0.6ex}{\poppill{Cours signé \textbullet\ Premium}{popink}{popcream}}\par
  \vspace{1pt}\popsquiggle{poppurple}\par
  \vspace{8pt}
  \begin{tcolorbox}[enhanced,colback=poporange,colframe=popink,boxrule=2.8pt,arc=13pt,
      drop shadow=popink,left=18pt,right=18pt,top=12pt,bottom=13pt,after skip=10pt]
    \poppill{#2 \textbullet\ #3}{popink}{popcream}\hspace{5pt}\poppill{#4}{white}{popink}\par\vspace{8pt}
    {\popdisplay\fontsize{14}{14}\selectfont\color{popink}LEÇON #5}\par\vspace{4pt}
    {\raggedright\hyphenpenalty=10000\exhyphenpenalty=10000\popdisplay\fontsize{25}{27}\selectfont\color{popcream}\MakeUppercase{#1}\par}
    \vspace{8pt}
    {\popxbold\fontsize{9.5}{9.5}\selectfont\color{popink}\MakeUppercase{Durée conseillée : #6}}
  \end{tcolorbox}
  \begin{tcolorbox}[enhanced,colback=white,colframe=popink,boxrule=2.2pt,arc=10pt,
      drop shadow=popink,left=16pt,right=16pt,top=10pt,bottom=10pt,after skip=8pt]
    \poppill{Dans cette leçon}{poppurple}{white}\par\vspace{5pt}
    \popsemi\fontsize{9.3}{12.6}\selectfont\color{popink}#7
  \end{tcolorbox}
  \noindent\begin{minipage}[t]{0.487\linewidth}
    \begin{tcolorbox}[enhanced,colback=popgreen,colframe=popink,boxrule=2.2pt,arc=10pt,
        drop shadow=popink,left=11pt,right=11pt,top=10pt,bottom=10pt,height=3.1cm,valign=center]
      \tcbox[colback=white,colframe=popink,boxrule=1.3pt,arc=5pt,boxsep=0pt,nobeforeafter,
        left=3pt,right=3pt,top=3pt,bottom=3pt,valign=center]{\qrcode[height=1.9cm]{https://play.google.com/store/apps/details?id=com.luvvix.onbuch&hl=fr}}%
      \hspace{8pt}\begin{minipage}[c]{0.5\linewidth}
        {\popdisplay\fontsize{11}{12.5}\selectfont\color{white}\MakeUppercase{Télécharge OnBuch}}\par\vspace{4pt}
        {\raggedright\popsemi\fontsize{8.4}{11}\selectfont\color{white}Gratuit \textbullet\ Google Play\\Tuteur IA, annales, concours, résultats\par}
      \end{minipage}
    \end{tcolorbox}
  \end{minipage}\hfill
  \begin{minipage}[t]{0.487\linewidth}
    \begin{tcolorbox}[enhanced,colback=poppurple,colframe=popink,boxrule=2.2pt,arc=10pt,
        drop shadow=popink,left=11pt,right=11pt,top=10pt,bottom=10pt,height=3.1cm,valign=center]
      \tikz[baseline=-0.7ex]\node[circle,fill=white,draw=popink,line width=1.3pt,minimum size=1.6cm,inner sep=0]{\popdisplay\fontsize{24}{24}\selectfont\color{poppurple}+};%
      \hspace{8pt}\begin{minipage}[c]{0.6\linewidth}
        {\popdisplay\fontsize{11}{12.5}\selectfont\color{white}\MakeUppercase{Passe à OnBuch+}}\par\vspace{4pt}
        {\raggedright\popsemi\fontsize{8.4}{11}\selectfont\color{white}Tous les cours signés, Tuteur IA illimité, fiches PDF — onglet OnBuch+ de l'app\par}
      \end{minipage}
    \end{tcolorbox}
  \end{minipage}
  \vspace{8pt}
  \popcontenu
  \vfill
  \signatureonbuch
  \restoregeometry
  \newpage
}

% Rangée « Ce cours contient » (page de garde)
\newcommand{\poptile}[3]{% #1 couleur #2 gros texte #3 légende
  \begin{tcolorbox}[enhanced,colback=#1,colframe=popink,boxrule=2pt,arc=9pt,shadow={2.5pt}{-2.5pt}{0pt}{popink},
      left=6pt,right=6pt,top=9pt,bottom=9pt,halign=center,valign=center,height=1.75cm,width=\linewidth,nobeforeafter]
    {\popdisplay\fontsize{12.5}{13}\selectfont\color{white}\MakeUppercase{#2}}\par\vspace{3pt}
    {\centering\popsemi\fontsize{7.8}{9.5}\selectfont\color{white}#3\par}
  \end{tcolorbox}}
\newcommand{\popcontenu}{%
  \par\noindent{\popxboldls\fontsize{7.5}{7.5}\selectfont\color{popmuted}CE COURS SIGNÉ CONTIENT}\par\vspace{7pt}
  \noindent\begin{minipage}[t]{0.235\linewidth}\poptile{popblue}{Cours}{complet, illustré, pas à pas}\end{minipage}\hfill
  \begin{minipage}[t]{0.235\linewidth}\poptile{poporange}{Méthodes}{et exercices résolus}\end{minipage}\hfill
  \begin{minipage}[t]{0.235\linewidth}\poptile{poppink}{Exercices}{type examen + corrigés}\end{minipage}\hfill
  \begin{minipage}[t]{0.235\linewidth}\poptile{popgreen}{Fiche bilan}{l'essentiel à retenir}\end{minipage}\par}

% Sommaire
\newtcolorbox{sommaire}{enhanced,colback=white,colframe=popink,boxrule=2.2pt,arc=10pt,
  drop shadow=popink,left=18pt,right=18pt,top=13pt,bottom=13pt,before skip=6pt,after skip=20pt,
  before upper={\poppill{Sommaire}{poppurple}{white}\par\vspace{9pt}\popsemi\fontsize{10.6}{14.5}\selectfont\color{popink}}}

% Signature de fin de cours — poussée tout en bas de la dernière page
\newcommand{\findecours}{%
  \par\vfill\nopagebreak
  \begin{center}\popsquiggle{poporange}\end{center}\vspace{4pt}
  \signatureonbuch
}
\AtBeginDocument{\def\llbracket{\mathopen{[\mkern-3.5mu[}}\def\rrbracket{\mathclose{]\mkern-3.5mu]}}} % intervalles d'entiers (glyphe ⟦ absent de la police)
% Glyphes absents de Fira Math : redéfinis à partir de glyphes disponibles
\AtBeginDocument{%
  \def\setminus{\mathbin{\backslash}}\let\smallsetminus\setminus
  \def\vdots{\mathord{\vbox{\baselineskip3.2pt\lineskiplimit0pt\kern4pt\hbox{.}\hbox{.}\hbox{.}}}}%
  \def\ddots{\mathinner{\mkern1mu\raise6.5pt\hbox{.}\mkern2mu\raise3.6pt\hbox{.}\mkern2mu\raise0.7pt\hbox{.}\mkern1mu}}%
  \def\triangle{\mathord{\scalebox{0.9}{$\symup{\Delta}$}}}%
  \def\nabla{\mathord{\raisebox{\depth}{\scalebox{1}[-1]{$\symup{\Delta}$}}}}%
  \def\checkmark{\popcheck{popdark}}%
  \def\star{\mathbin{\ast}}\def\bigstar{\mathord{\ast}}%
  \def\diamond{\mathbin{\scalebox{0.6}{$\Diamond$}}}%
  \def\bigcup{\mathop{\mathchoice{\raisebox{-0.3em}{\scalebox{1.7}{$\cup$}}}{\scalebox{1.25}{$\cup$}}{\cup}{\cup}}\displaylimits}%
  \def\bigcap{\mathop{\mathchoice{\raisebox{-0.3em}{\scalebox{1.7}{$\cap$}}}{\scalebox{1.25}{$\cap$}}{\cap}{\cap}}\displaylimits}%
  \def\lesseqqgtr{\mathrel{\lessgtr}}\def\bigoplus{\mathop{\oplus}\displaylimits}%
}
\providecommand{\ssi}{si et seulement si} % abréviation fréquente
\AtBeginDocument{\providecommand{\wideparen}{\overparen}} % arc de cercle

%% FIGFIX-BEGIN v5 — correctifs globaux des figures (docs/onbuch-generation/tools/figfix)
\usepackage{adjustbox}\usepackage{etoolbox}\tcbuselibrary{raster}
% toute figure TikZ/pgfplots plus large que la ligne est réduite à la largeur disponible
\BeforeBeginEnvironment{tikzpicture}{\begin{adjustbox}{max width=\linewidth}}
\AfterEndEnvironment{tikzpicture}{\end{adjustbox}}
% légendes pgfplots lisibles par-dessus les courbes
\pgfplotsset{every axis/.append style={legend style={fill=white,fill opacity=0.92,text opacity=1,draw=black!15,inner sep=3pt}}}
% étiquettes d'arêtes (syntaxe edge["..."]) avec fond blanc
\tikzset{every edge quotes/.append style={fill=white,inner sep=1.2pt}}
%% FIGFIX-END
\begin{document}

\pagegarde{Les révolutions scientifiques et techniques depuis la fin de la Deuxième Guerre mondiale Chapitre 6: La marche vers l’indépendance}{Histoire}{Tle Tech}{Histoire – classe de Terminale technique}{N08}{3 séances (fiche de progression harmonisée)}{Cette leçon analyse la transition du Cameroun de la tutelle onusienne à la souveraineté nationale (1945-1961), en mettant en lumière le rôle décisif des mouvements nationalistes et de leurs leaders historiques, conformément au programme officiel d'Histoire de Terminale Technique.
\vspace{4pt}
\begin{multicols}{2}\small
\begin{itemize}
\item À la fin de cette leçon, je sais expliquer le système de tutelle de l'ONU au Cameroun et le distinguer du mandat de la SDN.
\item Je sais analyser les causes profondes (économiques, sociales, politiques) de l'émergence du nationalisme camerounais après 1945.
\item Je sais identifier les grandes figures du nationalisme (Um Nyobé, Moumié, Ahidjo, Foncha, etc.) et comparer leurs stratégies politiques via le Dossier 3.
\item Je sais construire une frise chronologique commentée des étapes clés menant aux indépendances du 1er janvier 1960 et du 1er octobre 1961.
\item Je sais rédiger une réponse argumentée justifiant le choix de la réunification par les populations du Southern Cameroons en 1961.
\item Je sais appliquer la notion de souveraineté à une situation concrète actuelle (ex: gestion des ressources, siège à l'ONU).
\end{itemize}\end{multicols}}

\begin{sommaire}
\begin{multicols}{2}
\begin{enumerate}
\item 1. Le Cameroun sous la tutelle de l'ONU (1945-1960) : un cadre international contraignant
\item 2. Genèse et essor du nationalisme camerounais (1945-1955) : des syndicaux aux partis politiques
\item 3. Dossier 3 : Les grandes figures du nationalisme camerounais - Portraits croisés et stratégies
\item 4. Le cheminement institutionnel vers les indépendances et la réunification (1956-1961)
\item 5. Héritages, mémoires et enjeux historiographiques (Complément Bac / Éducation à la citoyenneté)
\item Activité d'intégration
\item Méthodes et exercices résolus
\item Exercices
\item Corrigés des exercices
\item Fiche bilan
\end{enumerate}
\end{multicols}
\end{sommaire}

\begin{objectifs}
\begin{itemize}
\item À la fin de cette leçon, je sais expliquer le système de tutelle de l'ONU au Cameroun et le distinguer du mandat de la SDN.
\item Je sais analyser les causes profondes (économiques, sociales, politiques) de l'émergence du nationalisme camerounais après 1945.
\item Je sais identifier les grandes figures du nationalisme (Um Nyobé, Moumié, Ahidjo, Foncha, etc.) et comparer leurs stratégies politiques via le Dossier 3.
\item Je sais construire une frise chronologique commentée des étapes clés menant aux indépendances du 1er janvier 1960 et du 1er octobre 1961.
\item Je sais rédiger une réponse argumentée justifiant le choix de la réunification par les populations du Southern Cameroons en 1961.
\item Je sais appliquer la notion de souveraineté à une situation concrète actuelle (ex: gestion des ressources, siège à l'ONU).
\item Je sais réaliser un croquis historique légendé « Le Cameroun sous tutelle : espaces et acteurs du nationalisme ».
\end{itemize}
\end{objectifs}

\begin{prerequis}
\begin{itemize}
\item Notion de mandat de la SDN et traité de Versailles (1919) - Programme de 3ème / 1ère.
\item Contexte de la Seconde Guerre mondiale et création de l'ONU (1945) - Programme de 1ère.
\item Géographie physique et administrative du Cameroun (régions, capitale, frontières).
\item Vocabulaire de base : colonie, protectorat, tutelle, nationalisme, parti politique, syndicat.
\item Méthode d'analyse de document historique (nature, auteur, date, portée).
\end{itemize}
\end{prerequis}

\newpage

\coursec{Le Cameroun sous la tutelle de l'ONU (1945-1960) : un cadre international contraignant}

\courssub{Accroche et problématique}

En 1956, à Nkongsamba, cœur de la zone de culture du café robusta, des milliers de cultivateurs refusent de livrer leur récolte aux acheteurs agréés par l'administration coloniale française. Ils brandissent le drapeau rouge de l'Union des Populations du Cameroun (UPC) et scandent : « Indépendance immédiate ! ». Ce geste de résistance paysanne, réprimé dans le sang, n'est pas un événement isolé. Il est l'expression locale d'une tension mondiale : le Cameroun, placé sous tutelle de l'Organisation des Nations Unies (ONU) depuis 1946, est censé être préparé à l'autonomie. Pourtant, l'administration française maintient le Code de l'indigénat jusqu'en 1946, le travail forcé sous couvert de « prestations » et verrouille la vie politique. Comment ce territoire, unique en Afrique par sa double tutelle franco-britannique, passe-t-il du statut de mandat de la Société des Nations (SDN) à celui de territoire sous tutelle de l'ONU, et pourquoi ce cadre international, théoriquement protecteur, devient-il le théâtre d'un affrontement violent entre nationalistes et puissances administratrices ?

\courssub{Du mandat de la SDN à la tutelle de l'ONU : rupture et continuité (1919-1946)}

\begin{definition}[Tutelle de l'ONU]
Régime juridique international défini par le Chapitre XII de la Charte des Nations Unies (1945), confiant l'administration d'un territoire à une puissance administrante (France ou Royaume-Uni) pour le préparer à l'autonomie ou à l'indépendance, sous le contrôle de l'Assemblée générale et du Conseil de tutelle. Contrairement au mandat de la SDN, la tutelle pose explicitement le principe de non-annexion et l'obligation de développement politique, économique, social et éducatif avec une finalité d'accession à la souveraineté.
\end{definition}

\begin{attention}[Mandat vs Tutelle : ne pas confondre]
\begin{itemize}
    \item \textbf{Mandat (SDN, 1919-1945) :} Article 22 du Pacte de la SDN. But : « bien-être et développement » des populations. Pas d'échéance fixe d'indépendance. Contrôle par la Commission des mandats (faible pouvoir de sanction).
    \item \textbf{Tutelle (ONU, 1946-1960) :} Chapitre XII de la Charte. But explicite : « évolution progressive vers l'autonomie ou l'indépendance » (Art. 76 b). Principe de non-annexion (Art. 77). Contrôle renforcé : rapports annuels obligatoires, pétitions recevables, missions de visite sur place.
\end{itemize}
La transition n'est pas une rupture totale : les puissances mandatées (France, UK) deviennent puissances administratrices. Les structures coloniales (administration, économie de plantation, Code de l'indigénat) persistent initialement.
\end{attention}

\begin{savaistu}[Le cas unique du Cameroun]
Le Cameroun est le \textbf{seul territoire d'Afrique} à avoir été administré simultanément par deux puissances tutélaires différentes (France et Royaume-Uni) sous l'égide de l'ONU. Cette double tutelle, héritée du partage de 1916 (ligne Picot / Milner-Simon), crée deux espaces juridiques, politiques et économiques distincts qui compliqueront la réunification.
\end{savaistu}

\paragraph{Les accords de tutelle du 13 décembre 1946.}
Signés à New York, ils lient la France et le Royaume-Uni à l'ONU.
\begin{itemize}
    \item \textbf{Cameroun français (Cameroun sous tutelle française) :} La France s'engage à appliquer l'article 76 de la Charte. Elle institue une Assemblée Représentative du Cameroun (ARTCAM, 1946) puis l'Assemblée Législative du Cameroun (ALCAM, 1952) avec des pouvoirs consultatifs puis délibératifs limités. Le Haut-Commissaire (représentant de la France) conserve l'essentiel de l'exécutif.
    \item \textbf{British Cameroons (Cameroun sous tutelle britannique) :} Le Royaume-Uni choisit l'administration indirecte via le Nigeria. Le territoire est divisé en Northern Cameroons (intégré à la région du Nord du Nigeria) et Southern Cameroons (région à part entière de la Fédération du Nigeria depuis 1954). Aucune assemblée territoriale propre n'existe avant 1954 (House of Assembly à Buéa pour le Sud).
\end{itemize}

\begin{popfigure}
\begin{tikzpicture}[scale=0.85]
% Cadre général
\draw[thick, popdark] (0,0) rectangle (14,7);
% Ligne de séparation (frontière Picot/Milner-Simon approximative)
\draw[thick, dashed, popmuted] (5.5,0) -- (5.5,7);
% Zone Française (Gauche)
\fill[popblue!15] (0,0) rectangle (5.5,7);
\node[align=center, font=\bfseries\large, text=popblue] at (2.75, 6.5) {CAMEROUN FRANÇAIS};
\node[align=center, font=\small] at (2.75, 5.8) {Tutelle : France (Haut-Commissaire)};
\node[align=center, font=\small] at (2.75, 5.2) {Capital : \textbf{Yaoundé}};
\node[align=left, font=\small] at (2.75, 4.2) {
    \textbf{Institutions :}\\
    ARTCAM (1946) $\to$ ALCAM (1952)\\
    Suffrage : Censitaire puis universel (1956)\\
    \textbf{Partis majeurs :}\\
    UPC (1948), PDC (1951), UC (1958)
};
\node[align=center, font=\tiny, text=popmuted] at (2.75, 1.5) {Économie : Plantations (Banane, Café, Cacao, Palmier)\newline Monnaie : Franc CFA (Zone Franc)};

% Zone Britannique (Droite)
\fill[popred!15] (5.5,0) rectangle (14,7);
\node[align=center, font=\bfseries\large, text=popred] at (9.75, 6.5) {BRITISH CAMEROONS};
\node[align=center, font=\small] at (9.75, 5.8) {Tutelle : Royaume-Uni (Gouverneur du Nigeria)};
\node[align=center, font=\small] at (9.75, 5.2) {Capital Sud : \textbf{Buéa} | Nord : \textbf{Bamenda}};
\node[align=left, font=\small] at (9.75, 4.2) {
    \textbf{Intégration administrative au Nigeria}\\
    Northern Cameroons $\to$ Région Nord Nigeria\\
    Southern Cameroons $\to$ Région Fédérée (1954)\\
    \textbf{Institutions (Sud) :}\\
    House of Assembly (Buéa, 1954)\\
    \textbf{Partis majeurs :}\\
    KNC, KPP, KNDP (J. N. Foncha)
};
\node[align=center, font=\tiny, text=popmuted] at (9.75, 1.5) {Économie : Intégrée au Nigeria (Marketing Boards)\newline Monnaie : Livre sterling / Livre nigériane};

% Flèche Nord
\begin{scope}[shift={(12.5, 6.0)}, rotate=0]
\draw[thick, ->, popdark] (0,0) -- (0,0.8) node[above] {N};
\end{scope}

% Légende frontière
\draw[<->, thick, poporange] (3, 0.5) -- (8, 0.5);
\node[below, font=\small, text=poporange] at (5.5, 0.2) {Frontière de tutelle (ligne Picot / Milner-Simon, 1916)};
\end{tikzpicture}
\legende{Figure 1 : Le Cameroun sous double tutelle (1946-1960) : structures institutionnelles et partisanes contrastées.}
\end{popfigure}

\courssub{Le contrôle international : l'ONU face aux réalités coloniales (1946-1960)}

L'ONU dispose de trois leviers principaux pour contrôler l'action des puissances administratrices : l'examen des rapports annuels, le droit de pétition et les missions de visite.

\paragraph{1. Les rapports annuels et le Conseil de tutelle.}
Chaque année, la France et le Royaume-Uni soumettent un rapport sur l'évolution politique, économique, sociale et éducative. Le Conseil de tutelle (organe principal de l'ONU jusqu'en 1994) l'examine, pose des questions écrites et orales, et adopte des résolutions.
\begin{itemize}
    \item \textbf{Points de friction récurrents :} Lenteur de la préparation politique (peu de cadres africains dans l'administration), maintien du travail forcé (prestations, corvées), répression des mouvements nationalistes (interdiction UPC en 1955), inégalité des budgets alloués à l'éducation/santé.
    \item \textbf{Exemple :} Le rapport français de 1952 vante les « progrès matériels » (routes, hôpital de Yaoundé) mais occulte l'absence de libertés syndicales et politiques réelles.
\end{itemize}

\paragraph{2. Le droit de pétition : la voix des populations.}
L'article 87 de la Charte autorise les habitants des territoires sous tutelle à adresser des pétitions à l'ONU. C'est une innovation majeure par rapport à la SDN.
\begin{exemplebox}[Les pétitions de l'UPC (1952-1955)]
L'Union des Populations du Cameroun (UPC), dirigée par Ruben Um Nyobé, utilise massivement ce droit.
\begin{itemize}
    \item \textbf{Pétition de 1952 (Um Nyobé à New York) :} Dénonce le « régime de terreur », l'absence de libertés, l'exploitation économique, exige l'indépendance immédiate et la réunification.
    \item \textbf{Pétition de 1954 :} Réclame la tenue d'élections libres sous contrôle de l'ONU.
    \item \textbf{Impact :} Ces pétitions obligent le Conseil de tutelle à inscrire « la question du Cameroun » à son ordre du jour, internationalisant le conflit. Elles irritent profondément la France qui tente d'en limiter la recevabilité (argument de « recevabilité » : l'UPC n'est pas un « organe représentatif » officiel).
\end{itemize}
\end{exemplebox}

\paragraph{3. Les missions de visite de l'ONU (1952, 1955, 1958).}
Des délégations du Conseil de tutelle se rendent sur place pour vérifier les dires des rapports et entendre les populations.
\begin{itemize}
    \item \textbf{Mission 1952 (présidée par le Libanais Charles Malik) :} Constate le fossé entre le discours officiel et la réalité (travail forcé, bas salaires, analphabétisme). Recommande l'abolition du Code de l'indigénat (déjà fait en 1946 mais pratiques persistantes) et l'accélération de la formation de cadres.
    \item \textbf{Mission 1955 (post-insurrection UPC) :} Contexte de guerre (mai-décembre 1955). La mission est encadrée de près par l'armée française. Son rapport est plus modéré, notant les « troubles » mais recommandant des réformes politiques.
    \item \textbf{Mission 1958 (pré-indépendance) :} Constate l'avancée de l'autonomie interne (Gouvernement André Marie Mbida puis Ahmadou Ahidjo) mais note l'exclusion de l'UPC (interdite, maquis) et le problème non résolu des British Cameroons.
\end{itemize}

\begin{methode}[Analyser un rapport de tutelle ou une mission de l'ONU]
Pour exploiter un document d'archive (rapport administrant, compte-rendu de mission, pétition) :
\begin{enumerate}
    \item \textbf{Identifier la source :} Qui écrit ? (Puissance administrante, Secrétariat ONU, Mission, Pétitionnaires). Quelle date ? Quel public visé ?
    \item \textbf{Repérer les indicateurs de « progrès » :} Chiffres (écoles, lits d'hôpitaux, km de routes, budget), création d'institutions (assemblées, gouvernements).
    \item \textbf{Repérer les indicateurs de « libertés » :} Existence de partis légaux, syndicats libres, presse non censurée, suffrage universel, absence de prisonniers politiques.
    \item \textbf{Confronter les discours :} Comparer le rapport de la France (auto-satisfaction) avec les pétitions de l'UPC (dénonciation) et les observations de la mission ONU (vérité terrain).
    \item \textbf{Conclure sur l'effectivité de la tutelle :} L'ONU a-t-elle obtenu des résultats concrets ? (Ex: abolition travail forcé 1946, suffrage universel 1956, mais répression 1955-1960).
\end{enumerate}
\end{methode}

\courssub{Les limites structurelles de la tutelle : un cadre contraignant mais insuffisant}

Malgré le cadre juridique protecteur de la Charte, la tutelle bute sur la souveraineté réelle des puissances administratrices.

\begin{propriete}[Les quatre limites majeures de la tutelle au Cameroun]
\begin{enumerate}
    \item \textbf{Maintien de l'ordre colonial :} Le Code de l'indigénat est aboli légalement en 1946 (Loi Lamine Guèye), mais les « prestations » (travail forcé pour routes, plantations) persistent jusqu'en 1950-1952. La justice indigène reste discriminatoire.
    \item \textbf{Absence de réelle préparation des cadres :} En 1955, sur 5000 fonctionnaires de catégorie supérieure, moins de 50 sont Camerounais. L'enseignement secondaire et supérieur est quasi inexistant sur place (création de l'Ecole Normale Supérieure de Yaoundé en 1961 seulement). La France forme des auxiliaires, pas des dirigeants.
    \item \textbf{Ingérence dans le processus électoral :} La loi-cadre Defferre (1956) instaure le suffrage universel. Mais les élections de 1956 (ALCAM) et 1957 (Assemblée territoriale) sont marquées par l'exclusion de l'UPC (interdite juillet 1955), le contrôle des listes par l'administration, et l'appui aux partis « modérés » (PDC, UC) par le Haut-Commissaire.
    \item \textbf{Division du territoire :} La tutelle britannique, en intégrant les British Cameroons au Nigeria, brise l'unité territoriale, économique (douane, monnaie, chemins de fer) et politique, rendant la réunification future complexe.
\end{enumerate}
\end{propriete}

\begin{exoresolu}[Exercice type Bac : Commentaire de document]
\textbf{Énoncé :} « Le système de tutelle, tel qu'il a été appliqué au Cameroun, a davantage servi les intérêts stratégiques des puissances administratrices que l'accession réelle des populations à la souveraineté. » À partir de vos connaissances et du document ci-après (extrait du rapport de la Mission de visite de l'ONU, 1958), discutez cette affirmation.

\textbf{Document (extrait simulé) :} « La Mission note avec satisfaction l'instauration du suffrage universel et la formation d'un gouvernement africain responsable devant l'Assemblée législative au Cameroun français. Elle note cependant que le principal parti nationaliste (UPC) reste interdit et que ses dirigeants sont en exil ou dans le maquis. Au British Cameroons, l'intégration au Nigeria se poursuit, rendant incertaine l'expression de la volonté des populations sur leur avenir distinct. »

\textbf{Correction détaillée :}
\begin{enumerate}
    \item \textbf{Analyse de l'affirmation :} Elle pose la thèse d'un détournement de la tutelle au profit des puissances coloniales.
    \item \textbf{Arguments POUR (Validation) :}
        \begin{itemize}
            \item \textit{Intérêt stratégique France :} Maintien du Cameroun dans la zone Franc, accords de coopération (défense, matières premières) signés avant l'indépendance (1960). Soutien aux partis modérés (Ahidjo) contre l'UPC (indépendantiste radicale).
            \item \textit{Intérêt stratégique UK :} Intégration économique au Nigeria (pétrole, main d'œuvre). Volonté de ne pas créer un micro-État non viable au Nord, ou de favoriser le rattachement au Nigeria pour le Sud.
            \item \textit{Preuve document :} Interdiction UPC (exclusion politique), poursuite intégration Nigeria (négation droit à l'autodétermination distincte).
        \end{itemize}
    \item \textbf{Arguments CONTRE (Nuance) :}
        \begin{itemize}
            \item \textit{Cadre contraignant :} L'ONU a imposé le suffrage universel (1956), l'abolition du travail forcé, la création d'institutions représentatives (ALCAM, Gouvernement).
            \item \textit{Preuve document :} « Satisfaction » pour suffrage universel et gouvernement responsable. Ce sont des acquis de la tutelle.
            \item \textit{Dynamique propre :} Le nationalisme camerounais (UPC, puis Ahidjo) a forcé la main aux puissances tutélaires. La tutelle a fourni une tribune internationale (pétitions, missions).
        \end{itemize}
    \item \textbf{Synthèse :} La tutelle a été un \textbf{cadre contraignant} (obligation de rendre des comptes, droit de pétition, missions) qui a accéléré les échéances (indépendance 1960 au lieu de 1970-80 hypothétique). Mais elle a été \textbf{détournée} par la France pour organiser une « indépendance surveillée » (élection d'Ahidjo, accords de coopération) et par le UK pour noyer le choix dans le plébiscite nigérian. L'affirmation est \textbf{largement fondée mais incomplète} : la tutelle a été un champ de bataille, pas seulement un outil de domination.
\end{enumerate}
\end{exoresolu}

\courssub{Focus Bac : La résolution 1476 (XIV) de 1959 et le sort des British Cameroons}

Alors que le Cameroun français accède à l'indépendance le 1er janvier 1960, le sort des British Cameroons reste en suspens. L'ONU doit organiser la consultation promise.

\begin{definition}[Résolution 1476 (XIV) - 12 décembre 1959]
Adoptée par l'Assemblée générale de l'ONU, elle fixe les modalités du plébiscite dans les British Cameroons (Nord et Sud) pour février 1961. Elle pose deux questions alternatives (pas d'option « indépendance pure ») :
\begin{enumerate}
    \item \textbf{Intégration au Nigeria indépendant (prévu pour octobre 1960).}
    \item \textbf{Réunification avec la République du Cameroun (indépendante depuis le 1er janv. 1960).}
\end{enumerate}
La résolution prévoit un scrutin à bulletin secret, sous supervision d'officiers de l'ONU, mais l'administration électorale reste britannique/nigériane.
\end{definition}

\begin{attention}[Piège historiographique : L'absence de l'option « Indépendance séparée »]
La résolution 1476 (XIV) \textbf{ne propose pas} l'indépendance comme État souverain distinct pour les British Cameroons. L'ONU (sur pression UK/USA) juge le territoire « non viable » économiquement (petite taille, pas de port en eau profonde au Sud, enclavement au Nord). C'est une entorse à l'article 76 b de la Charte (« autonomie \textbf{ou} indépendance »). Cela fausse le choix : les populations doivent choisir entre deux « tuteurs » (Nigeria ou Cameroun), pas la liberté pleine.
\end{attention}

\begin{popfigure}
\begin{tikzpicture}[scale=0.9]
% Titre
\node[font=\bfseries\large, text=popdark] at (7, 5.5) {Le Plébiscite de Février 1961 (Résolution 1476 XIV)};

% Carte schématique très simplifiée
% Nigeria (Gauche)
\fill[popgreen!10] (0,0) rectangle (5,4);
\node[align=center, font=\small, text=popgreen!80!black] at (2.5, 3.5) {NIGERIA\\ (Indépendance oct. 1960)};
\node[align=center, font=\tiny] at (2.5, 2.5) {Option A : Intégration};

% Cameroun République (Droite)
\fill[popblue!10] (9,0) rectangle (14,4);
\node[align=center, font=\small, text=popblue] at (11.5, 3.5) {RÉPUBLIQUE DU CAMEROUN\\ (Indépendance 1er janv 1960)};
\node[align=center, font=\tiny] at (11.5, 2.5) {Option B : Réunification};

% British Cameroons Nord (Haut milieu)
\fill[poporange!20] (5,2.5) rectangle (9,4);
\draw[thick, poporange] (5,2.5) rectangle (9,4);
\node[align=center, font=\small\bfseries, text=poporange] at (7, 3.5) {NORTHERN CAMEROONS};
\node[align=center, font=\scriptsize] at (7, 3.0) {Population : ~600 000};
\node[align=center, font=\scriptsize] at (7, 2.7) {Résultat : \textbf{60\% Nigeria} / 40\% Cameroun};
\draw[->, thick, popgreen] (7, 2.5) -- (3, 2.5) node[midway, below, font=\tiny] {Intégration (1er juin 1961)};

% British Cameroons Sud (Bas milieu)
\fill[poppurple!20] (5,0) rectangle (9,2.5);
\draw[thick, poppurple] (5,0) rectangle (9,2.5);
\node[align=center, font=\small\bfseries, text=poppurple] at (7, 1.75) {SOUTHERN CAMEROONS};
\node[align=center, font=\scriptsize] at (7, 1.25) {Population : ~300 000};
\node[align=center, font=\scriptsize] at (7, 0.95) {Résultat : 30\% Nigeria / \textbf{70\% Cameroun}};
\draw[->, thick, popblue] (7, 0) -- (11, 0) node[midway, below, font=\tiny] {Réunification (1er oct. 1961)};

% Flèches Nord
\begin{scope}[shift={(13, 4.5)}]
\draw[thick, ->, popdark] (0,0) -- (0,0.6) node[above] {N};
\end{scope}
\end{tikzpicture}
\legende{Figure 2 : Les choix du plébiscite de 1961 dans les British Cameroons. Le Nord choisit le Nigeria, le Sud choisit la réunification avec la République du Cameroun.}
\end{popfigure}

\courssub{Synthèse de la section : Un cadre contraignant, un tremplin contesté}

La tutelle de l'ONU (1946-1960) a constitué un \textbf{cadre international contraignant} pour la France et le Royaume-Uni : obligation de rendre des comptes, droit de pétition pour les populations, missions de contrôle. Elle a légitimé la revendication nationaliste (référence à la Charte, article 76) et imposé des échéances (suffrage universel 1956, indépendance 1960).

Pourtant, elle n'a pas empêché la \textbf{guerre coloniale} (1955-1960, 10 000 à 100 000 morts selon les estimations) ni l'\textbf{ingérence électorale}. Elle a géré la double tutelle sans imposer l'unité préalable, aboutissant à une réunification par plébiscite (1961) tronquée (perte du Nord). Le Cameroun accède à la souveraineté dans un cadre négocié par les élites modérées (Ahidjo) sous l'œil de l'ONU, mais au prix de l'exclusion de la frange radicale (UPC) et de l'amputation territoriale du Northern Cameroons.

\begin{aretenir}[Essentiel pour le Bac - Section 1]
\begin{itemize}
    \item \textbf{1946 :} Accords de tutelle France/UK-ONU. Double particularité : double tutelle + double héritage colonial.
    \item \textbf{Institutions :} ALCAM (Fr) vs Intégration Nigeria (UK). Suffrage universel 1956 (Loi-cadre Defferre).
    \item \textbf{Contrôle ONU :} Rapports annuels, Pétitions (UPC/Um Nyobé), Missions de visite (1952, 1955, 1958).
    \item \textbf{Limites :} Travail forcé persistant, peu de cadres formés, répression UPC (interdite 1955), ingérence électorale.
    \item \textbf{1959-1961 :} Indépendance Cameroun français (1 janv 1960). Résolution 1476 (XIV) $\to$ Plébiscite Fév 1961. Nord $\to$ Nigeria, Sud $\to$ Réunification (1er oct 1961).
\end{itemize}
\end{aretenir}

\coursec{Genèse et essor du nationalisme camerounais (1945-1955) : des syndicats aux partis politiques}

Nous avons vu dans la section précédente que la tutelle de l'ONU, loin d'être une simple continuation du mandat de la SDN, imposait aux puissances administrantes (France et Grande-Bretagne) une obligation de résultat : conduire les populations du Cameroun vers l'autonomie, voire l'indépendance. Ce cadre international contraignant ouvre un espace juridique et politique nouveau. Encore faut-il que des hommes et des femmes s'en saisissent. C'est précisément l'objet de cette section : comprendre comment, entre 1945 et 1955, le mécontentement social hérité de la colonisation se transforme en un \cle{nationalisme} organisé, passant des revendications syndicales (salaires, travail forcé) à l'exigence politique de l'indépendance et de la réunification.

\courssub{Les facteurs endogènes : un terreau de mécontentements}

Le nationalisme ne naît pas d'une simple idée abstraite : il plonge ses racines dans les conditions concrètes de vie imposées par la colonisation. Quatre griefs majeurs structurent la colère des Camerounais après 1945.

\begin{itemize}
\item \textbf{L'exploitation économique.} Les planteurs camerounais vendent leur cacao et leur café à des \cle{bas prix} fixés unilatéralement par l'administration et les grandes compagnies commerciales européennes, qui détiennent des \cle{monopoles} d'achat et d'exportation. Le producteur africain est ainsi spolié de la plus-value de son travail, tandis que les bénéfices s'envolent vers la métropole.
\item \textbf{Le travail forcé.} Malgré sa condamnation internationale, le système des \cle{prestations} (journées de travail gratuites imposées pour les routes, les ports, les chantiers publics) perdure dans les faits après 1945, surtout dans les zones rurales. Il est vécu comme un vestige humiliant de l'esclavage.
\item \textbf{La discrimination raciale.} Les villes comme Douala et Yaoundé sont spatialement séparées : quartiers européens modernes (eau courante, électricité) d'un côté, quartiers indigènes surpeuplés et sans infrastructures de l'autre. L'accès aux écoles, aux hôpitaux et aux emplois qualifiés reste quasi fermé aux Africains.
\item \textbf{La fiscalité lourde.} L'\cle{impôt de capitation}, payé par chaque adulte, étrangle les paysans qui doivent vendre leurs récoltes à bas prix pour obtenir la monnaie exigée par l'administration.
\end{itemize}

\begin{attention}[Piège fréquent]
Ne pas confondre \emph{facteurs endogènes} (internes : exploitation, travail forcé, discriminations, impôts) et \emph{facteurs exogènes} (externes : ONU, Brazzaville, Bandung). Au Baccalauréat technique, une dissertation qui mélange les deux catégories sans les distinguer perd des points d'organisation. Annoncez toujours clairement le plan : causes internes, puis causes externes.
\end{attention}

\courssub{Les facteurs exogènes : le souffle du monde}

Le mécontentement intérieur rencontre un contexte international devenu favorable à l'émancipation des peuples colonisés.

\begin{itemize}
\item \textbf{La Conférence de Brazzaville (1944).} Réunie par le général de Gaulle, elle promet des réformes (suppression du travail forcé, représentation politique) mais exclut explicitement toute perspective d'\emph{indépendance}. C'est une immense \cle{déception} pour les élites africaines : la France veut réformer l'Empire, pas le dissoudre.
\item \textbf{La Charte de l'ONU (1945).} Son article premier consacre le \cle{droit des peuples à disposer d'eux-mêmes}. Pour les territoires sous tutelle comme le Cameroun, ce principe devient une arme juridique : les nationalistes pourront désormais pétitionner directement auprès de l'ONU.
\item \textbf{L'indépendance de l'Inde (1947).} La libération du « joyau de la couronne britannique » prouve qu'un empire colonial peut être vaincu par la mobilisation politique. Elle électrise les élites africaines.
\item \textbf{La Conférence de Bandung (1955).} Vingt-neuf pays d'Asie et d'Afrique y affirment l'anticolonialisme et le non-alignement. Le Cameroun n'y participe pas (il n'est pas encore indépendant), mais l'événement conforte les radicaux dans l'idée que l'émancipation est inéluctable.
\item \textbf{Le rôle des anciens combattants.} Des milliers de \cle{tirailleurs} camerounais ont combattu pour la liberté de la France entre 1939 et 1945. De retour au pays, ils ne comprennent pas qu'on leur refuse à eux-mêmes cette liberté. Ils apportent au mouvement nationaliste une expérience du monde, une conscience politique et parfois un savoir-faire militaire.
\end{itemize}

\begin{definition}[Nationalisme camerounais]
Le \cle{nationalisme camerounais} est le mouvement idéologique et politique qui, après la Deuxième Guerre mondiale, vise la libération du joug colonial, la \cle{réunification} des deux Cameroun (français et britannique) et l'accession du pays à la souveraineté internationale.
\end{definition}

\courssub{Du syndicalisme à la politique : l'USCC et les grèves de 1945-1946}

Avant les partis, ce sont les \cle{syndicats} qui servent de creuset à la contestation. En 1945 est créée l'\textbf{USCC} (Union des Syndicats Confédérés du Cameroun), liée à la CGT française. Son programme initial est socio-économique : hausse des salaires, suppression du travail forcé, égalité de traitement entre travailleurs africains et européens.

Les grandes grèves de 1945-1946 à \textbf{Douala} (le port, cœur économique du pays) et à \textbf{Edéa} (site industriel et hydroélectrique) marquent un tournant. Les grévistes constatent que l'administration coloniale protège systématiquement les intérêts des compagnies européennes. La leçon est vite tirée : \emph{on ne peut pas obtenir la justice sociale sans conquérir d'abord le pouvoir politique}. Le revendicatif socio-économique glisse ainsi vers le politique : c'est la naissance du nationalisme organisé.

\begin{popfigure}
\begin{tikzpicture}[scale=0.75]
% Axe principal
\draw[->, thick, popdark] (0,0) -- (12.5,0) node[right] {\small Temps};
% Graduations et dates
\foreach \x/\lab in {0/1945, 2/1948, 4/1952, 6/1955, 8/1958, 10/1960}
  \draw[popdark] (\x,0.1) -- (\x,-0.1) node[below] {\small \lab};
% Événements (alignés sur les graduations ou intercalés)
\node[align=center, fill=popgoldL, rounded corners, draw=popgold, font=\scriptsize, inner sep=2pt] at (0.5,1.8) {Fin 2\up{de} GM\\Retour tirailleurs};
\node[align=center, fill=poporangeL, rounded corners, draw=poporange, font=\scriptsize, inner sep=2pt] at (2,3.0) {Création UPC\\(10 avril 1948)};
\node[align=center, fill=popgreenL, rounded corners, draw=popgreen, font=\scriptsize, inner sep=2pt] at (4,1.8) {Multiplication partis\\(PDC 1951, UC 1952)};
\node[align=center, fill=popink!80, rounded corners, text=white, font=\scriptsize, inner sep=2pt] at (6,3.0) {Insurrection UPC\\Interdiction (juil. 1955)};
\node[align=center, fill=popblue!70, rounded corners, text=white, font=\scriptsize, inner sep=2pt] at (6.7,1.8) {Loi-cadre Defferre\\(juin 1956)};
\node[align=center, fill=poppurple!70, rounded corners, text=white, font=\scriptsize, inner sep=2pt] at (8,3.0) {Autonomie interne\\(1958)};
\node[align=center, fill=popred!70, rounded corners, text=white, font=\scriptsize, inner sep=2pt] at (10,1.8) {Indépendance\\1$^{\text{er}}$ janv. 1960};
% Lignes de liaison
\foreach \x in {0.5, 2, 4, 6, 6.7, 8, 10}
  \draw[thin, popmuted] (\x,0.2) -- (\x,1.4);
\end{tikzpicture}
\legende{Frise chronologique : de la fin de la Deuxième Guerre mondiale à l'indépendance du Cameroun (1945-1960). On observe l'accélération du processus après la création de l'UPC en 1948 et le cadre légal posé par la loi-cadre Defferre (1956).}
\end{popfigure}

\courssub{La création des premiers partis politiques}

\textbf{L'UPC : l'option radicale.} Le 10 avril 1948 naît à Douala l'\textbf{Union des Populations du Cameroun} (UPC), dont la figure emblématique est \cle{Ruben Um Nyobé}, secrétaire général. Son programme est le plus avancé de tous : \emph{indépendance immédiate} et \emph{réunification} des deux Cameroun. L'UPC multiplie les pétitions à l'ONU, s'implante massivement chez les paysans du Sanaga-Maritime et du pays Bamiléké, et refuse toute compromission avec l'administration française.

\textbf{Les partis modérés.} En face, d'autres formations préfèrent la voie de la négociation :
\begin{itemize}
\item le \textbf{PDC} (Parti des Démocrates Camerounais, 1951), dirigé par Jean Mouliom, se réclame de l'amitié franco-camerounaise et d'une évolution progressive ;
\item l'\textbf{UC} (Union Camerounaise, 1952), créée par \cle{Ahmadou Ahidjo}, défend une évolution négociée avec la France et une organisation de type fédéral.
\end{itemize}

\textbf{Dans le Cameroun britannique}, le nationalisme s'organise autour du \textbf{KNC} (Kamerun National Congress) d'Emmanuel Endeley, puis du \textbf{KNDP} (Kamerun National Democratic Party) de John Ngu Foncha, qui fera le choix de la réunification avec le Cameroun francophone.

\begin{propriete}[Le clivage structurel du nationalisme camerounais]
Le nationalisme camerounais est traversé par un clivage durable entre deux stratégies :
\begin{itemize}
\item les \textbf{Radicaux / Réunificationnistes} (UPC) : indépendance immédiate, réunification, lutte de masse voire armée, hostilité à la tutelle, sensibilité panafricaniste ;
\item les \textbf{Modérés / Évolutionnistes} (UC, PDC, KNC/KNDP) : indépendance négociée par étapes, coopération avec la France ou la Grande-Bretagne, légalité électorale.
\end{itemize}
Ce clivage porte à la fois sur la \emph{méthode} (lutte armée contre négociation) et sur l'\emph{orientation internationale} (panafricanisme contre francophonie ou commonwealth).
\end{propriete}

\begin{methode}[Comparer deux programmes de parti]
Pour comparer rigoureusement deux partis nationalistes (par exemple UPC et UC) dans une copie d'examen :
\begin{enumerate}
\item Construire un tableau à double entrée avec quatre critères : objectif final, moyens d'action, rapport à la puissance tutrice, base sociale.
\item Remplir chaque case avec un élément factuel daté (pas d'impression vague).
\item Conclure en nommant le clivage (radical / modéré) et en expliquant ses conséquences (interdiction de l'UPC en 1955, accès des modérés au pouvoir).
\end{enumerate}
\end{methode}

\begin{exemplebox}[L'ancrage populaire de l'UPC en chiffres]
En 1955, l'UPC revendique environ \num{300000} adhérents (chiffre du parti, à manier avec prudence), contre environ \num{50000} pour l'UC selon les estimations administratives. Même en divisant le chiffre upéciste par deux, l'écart reste considérable : il témoigne de l'enracinement initial de l'UPC dans les campagnes du Sanaga-Maritime et de l'Ouest (pays Bamiléké), là où l'exploitation économique et le travail forcé étaient les plus durement ressentis.
\end{exemplebox}

\begin{center}
\begin{tabularx}{\linewidth}{|l|X|X|}\hline
\rowcolor{popblueL} \textbf{Critère} & \textbf{UPC (Um Nyobé, 1948)} & \textbf{UC (Ahidjo, 1952)} \\ \hline
Objectif final & Indépendance immédiate et réunification & Indépendance négociée par étapes \\ \hline
Moyens & Pétitions à l'ONU, lutte de masse puis lutte armée & Légalité, élections, institutions \\ \hline
Rapport à la France & Rupture, dénonciation de la tutelle & Coopération, francophonie \\ \hline
Base sociale & Paysans du Sanaga-Maritime et Bamiléké, syndicalistes & Élites administratives, Nord et notables \\ \hline
\end{tabularx}
\end{center}

\courssub{L'apport des femmes au nationalisme}

Le nationalisme camerounais n'est pas une affaire d'hommes seuls. En 1952 est créée l'\textbf{Union des Femmes du Cameroun} (UFC), proche de l'UPC. Des figures comme \cle{Marthe Moumié} et \cle{Julienne Ngoumou} jouent un rôle décisif : mobilisation des femmes des marchés et des quartiers, rédaction et transport des pétitions, intendance (ravitaillement, logistique) des maquis après l'interdiction de l'UPC. Les femmes assurent la continuité du mouvement lorsque les dirigeants masculins sont emprisonnés ou exilés.

\begin{savaistu}[Le saviez-vous ?]
Ruben Um Nyobé, surnommé le \emph{Mpodol} (« celui qui porte la parole » en langue bassa), s'est rendu trois fois à New York (1952, 1953, 1954) pour défendre personnellement la cause camerounaise à la tribune de l'ONU. C'est l'un des premiers nationalistes africains à utiliser ainsi la scène internationale. Assassiné par l'armée française le 13 septembre 1958, il reste une figure fondatrice de la mémoire nationale camerounaise.
\end{savaistu}

\begin{exoresolu}[Exercice d'application : analyser une évolution]
\textbf{Énoncé.} En 1945, l'USCC réclame la hausse des salaires et la fin des prestations. En 1948, l'UPC exige l'indépendance immédiate et la réunification. En trois ans, quelle transformation du mouvement de contestation observe-t-on ? Expliquez-en deux causes.
\tcblower
\textbf{Solution détaillée.} On observe le passage d'un mouvement \emph{syndical} (revendications socio-économiques : salaires, travail forcé) à un mouvement \emph{politique et nationaliste} (conquête de la souveraineté : indépendance, réunification). Deux causes expliquent cette radicalisation : (1) l'échec des grèves de 1945-1946, qui démontre que l'administration coloniale protège les intérêts économiques européens et que seule la conquête du pouvoir politique permettrait de changer les rapports de force ; (2) le contexte international favorable (Charte de l'ONU et droit des peuples à disposer d'eux-mêmes, indépendance de l'Inde en 1947), qui fournit aux militants camerounais un cadre juridique et un exemple à imiter.
\end{exoresolu}

\begin{aretenir}[L'essentiel de la section]
\begin{itemize}
\item Le nationalisme camerounais naît de la rencontre entre des griefs internes (bas prix, prestations, discriminations, capitation) et un contexte externe porteur (ONU, Brazzaville, Inde, Bandung, tirailleurs).
\item L'USCC (1945) et les grèves de Douala et Edéa (1945-1946) font le pont entre le social et le politique.
\item L'UPC (1948, Um Nyobé) incarne l'option radicale ; le PDC (1951) et l'UC (1952, Ahidjo) l'option modérée ; le KNC puis le KNDP structurent le nationalisme anglophone.
\item Les femmes, à travers l'UFC (1952), sont des actrices à part entière de la lutte.
\end{itemize}
\end{aretenir}

\begin{exercice}[Pour s'entraîner]
À l'aide d'un tableau à quatre critères (objectif, moyens, rapport à la puissance tutrice, base sociale), comparez le KNDP de John Ngu Foncha et l'UPC de Ruben Um Nyobé, puis montrez que la question de la réunification les rapproche malgré leurs différences de méthode.
\end{exercice}

\coursec{Dossier 3 : Les grandes figures du nationalisme camerounais — Portraits croisés et stratégies}

Nous venons d'analyser, dans la section précédente, comment le nationalisme camerounais est né après 1945 : d'abord dans les syndicats et les associations, puis dans les partis politiques (UPC en 1948, puis UC, KNDP, PDC...). Mais derrière ces sigles se cachent des \cle{hommes et des femmes} : des instituteurs, des médecins, des syndicalistes, des écrivains, des fonctionnaires. Comprendre la marche vers l'indépendance, c'est donc comprendre leurs \cle{choix stratégiques} : fallait-il négocier avec la France ou lutter les armes à la main ? Fallait-il accepter une indépendance « dans les liens » avec l'ancien tuteur, ou exiger une indépendance totale et immédiate ? Ce dossier dresse les portraits croisés des grandes figures du nationalisme camerounais et compare leurs stratégies dans le contexte de la \cle{Guerre froide}.

\begin{methode}[Comment construire la fiche d'une figure du nationalisme (Dossier 3)]
Pour chaque personnalité étudiée, la fiche doit suivre cinq rubriques :
\begin{enumerate}
\item \textbf{Origine sociale et formation} (ethnie, région, métier, études).
\item \textbf{Engagement politique} (parti, rôle, date d'adhésion ou de fondation).
\item \textbf{Action majeure} (discours à l'ONU, négociation, maquis, plébiscite).
\item \textbf{Fin tragique ou accession au pouvoir} (assassinat, exécution, présidence).
\item \textbf{Postérité et mémoire} (héros national, réhabilitation, oubli volontaire, controverse).
\end{enumerate}
Cette grille permet de comparer rigoureusement les parcours et d'éviter le simple récit biographique. Elle est exigible au Probatoire / Baccalauréat technique pour toute question de synthèse sur les acteurs de la décolonisation.
\end{methode}

\begin{popfigure}
\begin{center}
\begin{tikzpicture}[scale=0.85, every node/.style={font=\small}]
% Frise chronologique simplifiée 1945-1971
\draw[popdark, thick] (0,0) -- (16,0);
\foreach \x/\year in {0/1945, 2/1948, 4/1952, 6/1955, 8/1957, 10/1958, 12/1960, 14/1961, 16/1971} {
  \draw (\x,-0.15) -- (\x,0.15) node[above=2pt] {\year};
}
% Événements clés
\node[align=center, anchor=north] at (1,-0.5) {Fin 2e GM\\Création ONU};
\node[align=center, anchor=north, fill=popblueL, rounded corners, inner sep=2pt] at (2.5,-0.5) {Fondation\\UPC (avril 1948)};
\node[align=center, anchor=north] at (4,-0.5) {1er discours\\Um Nyobé à l'ONU};
\node[align=center, anchor=north, fill=poporangeL, rounded corners, inner sep=2pt] at (6.5,-0.5) {Loi-cadre Defferre\\Interdiction UPC\\Maquis débute};
\node[align=center, anchor=north, fill=popgreenL, rounded corners, inner sep=2pt] at (8.5,-0.5) {Autonomie interne\\Ahidjo P.M.};
\node[align=center, anchor=north, fill=poppinkL, rounded corners, inner sep=2pt] at (10.5,-0.5) {Mort Um Nyobé\\Indépendance 1er janv 1960};
\node[align=center, anchor=north, fill=popgoldL, rounded corners, inner sep=2pt] at (12.5,-0.5) {Plébiscite 11 fév\\Réunification 1er oct};
\node[align=center, anchor=north] at (15,-0.5) {Fin maquis\\Exécution Ouandié};
% Acteurs
\node[align=center, anchor=south, text width=2.5cm] at (1,1.2) {\textbf{Um Nyobé}\newline 1913-1958};
\node[align=center, anchor=south, text width=2.5cm] at (5,1.2) {\textbf{Moumié}\newline 1925-1960};
\node[align=center, anchor=south, text width=2.5cm] at (9,1.2) {\textbf{Ahidjo}\newline 1924-1989};
\node[align=center, anchor=south, text width=2.5cm] at (13,1.2) {\textbf{Foncha}\newline 1916-1999};
\end{tikzpicture}
\end{center}
\legende{Frise chronologique des grandes figures et des étapes clés (1945-1971). En couleur : fondations (bleu), rupture 1955 (orange), accession au pouvoir modérés (vert), indépendance/réunification (rose/or), fin du maquis (or).}
\end{popfigure}

\courssub{Ruben Um Nyobé (1913-1958) : le « Mpodol », théoricien de la lutte de masse}

\textbf{Origine et formation.} Né en 1913 à Song Mpeck, dans le département de la Sanaga-Maritime (région du Littoral), Ruben Um Nyobé est issu du peuple \cle{bassa}. Comptable de formation, il travaille comme commis-comptable à la société forestière \emph{SFID} (Société forestière et industrielle du Douala) à Édéa, où il devient syndicaliste à l'\cle{USCC} (Union des syndicats confédérés du Cameroun), affiliée à la CGT française. Il y défend les travailleurs camerounais contre les abus du travail forcé (\cle{prestation}) et des compagnies coloniales.

\textbf{Engagement politique.} En avril 1948, il participe à la fondation de l'\cle{UPC} (Union des populations du Cameroun) lors du congrès constitutif à la salle des fêtes de Douala, et en devient le \cle{secrétaire général}. Surnommé le \emph{Mpodol} (le « porte-parole » en langue bassa), il fait de l'UPC un parti de masse structuré en comités de base, qui revendique deux objectifs inséparables : l'\cle{indépendance totale} du Cameroun et la \cle{réunification} des deux territoires (Cameroun français sous tutelle ONU et Cameroun britannique sous administration nigériane).

\textbf{Action majeure : le pétitionnaire de l'ONU.} Um Nyobé est un pétitionnaire infatigable. Il se rend trois fois à New York pour plaider la cause camerounaise devant la \cle{Quatrième Commission} (décolonisation) de l'Assemblée générale des Nations unies : en \textbf{1952}, \textbf{1954} et \textbf{1958}. Sa stratégie est originale : utiliser le cadre juridique international (la tutelle de l'ONU, article 76 de la Charte) contre la puissance tutélaire elle-même, tout en organisant sur le terrain la mobilisation de masse (grèves, meetings, pétitions).

\begin{savaistu}[Le discours du 17 décembre 1952]
Ruben Um Nyobé prononce son célèbre discours à la tribune de l'ONU le 17 décembre 1952, s'exprimant en langue \emph{bassa} (traduit pour l'assemblée), geste hautement symbolique d'affirmation culturelle et de refus de l'assimilation linguistique. Il y déclare : « Nous voulons l'indépendance du Cameroun, non pas dans dix ans, non pas dans cinq ans, mais \emph{maintenant}. » Ce discours fait de lui le visage du nationalisme camerounais dans le monde et contraint la France à rendre des comptes sur son administration de la tutelle.
\end{savaistu}

\textbf{Fin tragique.} Après l'interdiction de l'UPC par le Haut-commissaire Pierre Messmer (décret du 13 juillet 1955), Um Nyobé entre en clandestinité et organise la lutte armée dans le maquis de la Sanaga-Maritime. Il est tué au \cle{maquis} le \textbf{13 septembre 1958}, près de Boumnyebel, par une patrouille des forces de sécurité françaises (groupement de commandos parachutistes). Son corps est exposé publiquement à Éseka pour décourager les maquisards.

\textbf{Postérité.} Longtemps effacé de la mémoire officielle (son nom était tabou sous Ahidjo), Um Nyobé est réhabilité à partir des années 1990. Il est aujourd'hui honoré comme le symbole de l'indépendance « vraie » — celle qui aurait dû couper tous les liens de dépendance néocoloniale — et de la réunification. Il figure au Panthéon national camerounais (stèle au Mémorial de Yaoundé).

\begin{definition}[Le maquis camerounais]
Le \cle{maquis} désigne la guérilla nationaliste menée par l'UPC dans les zones forestières du Cameroun (Sanaga-Maritime, pays bamiléké à l'Ouest, pays bamoun au Nord-Ouest, puis région de l'Est) de \textbf{1955 à 1971}. Objectifs : obtenir l'indépendance totale et la réunification du Cameroun. Le maquis est violemment réprimé par l'armée française (opérations \emph{Bulldozer}, \emph{Epervier}...), puis par l'armée camerounaise après 1960, avec un lourd bilan humain (estimé entre 10 000 et 100 000 morts selon les sources), notamment dans les populations civiles soumises à la \cle{regroupement} (villages de regroupement forcés).
\end{definition}

\courssub{Félix-Roland Moumié (1925-1960) : le médecin qui internationalisa la lutte}

\textbf{Origine et formation.} Né en 1925 à Foumban (région de l'Ouest, pays bamoun), Félix-Roland Moumié est \cle{médecin}, formé à l'École de médecine de Dakar (William Ponty). Il représente la première génération d'intellectuels camerounais formés dans les institutions coloniales d'excellence.

\textbf{Engagement et action majeure.} Dirigeant de l'UPC (vice-président), il succède \emph{de facto} à Um Nyobé après la mort de celui-ci en 1958. Exilé (Conakry, Accra, Alger, Le Caire, Pékin), il mène une stratégie d'\cle{internationalisation de la lutte} : il transforme la cause camerounaise en cause panafricaine et tiers-mondiste. Il installe à Conakry (Guinée de Sékou Touré) un « Comité révolutionnaire » de l'UPC, obtient le soutien du Ghana (Nkrumah), de l'Algérie (FLN, GPRA), de la Chine (Zhou Enlai) et de l'URSS.

\begin{exemplebox}[Un soutien international chiffré]
Félix Moumié obtient le soutien de \textbf{17 chefs d'État africains et asiatiques} pour la cause de l'UPC à l'ONU, notamment grâce à la Conférence d'Accra (décembre 1958) et au mouvement des non-alignés (Conférence de Belgrade, 1961). Cette diplomatie parallèle montre que la question camerounaise est devenue un enjeu de la \cle{Guerre froide} en Afrique : l'UPC est perçue par l'Ouest comme le « Viet Cong africain », par l'Est comme le fer de lance de la révolution anti-impérialiste.
\end{exemplebox}

\textbf{Fin tragique.} Le \textbf{3 novembre 1960}, Félix-Roland Moumié meurt à Genève (Suisse), empoisonné au \cle{thallium} (un poison indétectable à l'époque, interdit par la Convention de Genève de 1925) par un agent des services secrets français (le SDECE, Service de documentation extérieure et de contre-espionnage), William Bechtel. Son assassinat, reconnu officiellement par la justice suisse des décennies plus tard (arrêt du Tribunal fédéral suisse, 2004), illustre la violence de la répression menée contre l'UPC, y compris hors du territoire camerounais et après l'indépendance formelle.

\textbf{Postérité.} Moumié incarne le nationalisme panafricain et radical, éliminé physiquement parce qu'il menaçait les intérêts stratégiques français (uranium, bases militaires, influence) au Cameroun indépendant. Son corps repose au cimetière de Foumban.

\courssub{Ahmadou Ahidjo (1924-1989) : le stratège du « fédéralisme pragmatique »}

\textbf{Origine et formation.} Né en 1924 à Garoua, dans le Nord-Cameroun (région du Nord), Ahmadou Ahidjo est le fils d'un chef peul (lamido). Employé des Postes et Télécommunications (radio-opérateur), il entre en politique en 1947 comme conseiller à l'Assemblée représentative du Cameroun (ARCAM), puis à l'Assemblée territoriale (ATCAM) en 1952.

\textbf{Engagement et stratégie.} Contrairement à l'UPC, Ahidjo choisit la voie de la \cle{négociation} avec la France. Sa stratégie, que les historiens qualifient de « fédéralisme pragmatique », consiste à obtenir l'indépendance par étapes, sans rupture brutale avec l'ancienne puissance tutélaire, tout en préservant l'unité territoriale :
\begin{itemize}
\item \textbf{1957 (mai)} : il négocie, comme vice-président de l'ATCAM, le statut d'\cle{État sous tutelle autonome} (loi-cadre Defferre) ;
\item \textbf{février 1958} : il devient \cle{Premier ministre} du gouvernement autonome, après l'éviction d'André-Marie Mbida par le Haut-commissaire Messmer ;
\item \textbf{1\textsuperscript{er} janvier 1960} : il proclame l'indépendance du Cameroun oriental (ex-Cameroun français) ;
\item \textbf{mai 1960} : il est élu premier Président de la République du Cameroun (suffrage universel indirect) ;
\item \textbf{1\textsuperscript{er} octobre 1961} : il devient Président de la \cle{République fédérale du Cameroun} après la réunification.
\end{itemize}

\textbf{Politique d'union nationale.} Au pouvoir, Ahidjo mène une politique d'« union nationale » qui aboutit en \textbf{1966} à la création du \cle{parti unique} (UNC, Union nationale camerounaise), fusionnant son parti (l'UC, Union camerounaise) avec le KNDP de Foncha, le PDC (restes) et d'autres formations. Il justifie ce choix par la nécessité de l'unité nationale face aux divisions ethniques et régionales, mais il instaure aussi un régime autoritaire qui poursuit la répression du maquis UPC (loi du 28 juin 1962 sur la sûreté de l'État).

\begin{attention}[Le piège de la « collaboration »]
Ne pas confondre \emph{stratégie modérée} et \emph{collaboration}. Ahidjo a été syndicaliste (USCC), a siégé dans les institutions coloniales pour les transformer de l'intérieur, et a obtenu l'indépendance \emph{réelle} (souveraineté internationale, siège à l'ONU) le 1\textsuperscript{er} janvier 1960. Son choix tactique — négocier plutôt que mourir au maquis — a permis l'accession à la souveraineté, mais au prix d'une continuité administrative, militaire et économique avec la France (accords de coopération de 1960). L'historien doit analyser ce \cle{compromis historique}, non le juger a priori.
\end{attention}

\courssub{John Ngu Foncha (1916-1999) : l'artisan de la réunification}

\textbf{Origine et formation.} Né en 1916 à Nkwen (Bamenda, Grassfield, Nord-Ouest), John Ngu Foncha est \cle{instituteur} de l'enseignement missionnaire (Basel Mission). Il représente le \cle{Cameroun britannique} (Southern Cameroons), trop souvent oublié dans les récits centrés sur le Cameroun français.

\textbf{Engagement et action majeure.} Foncha fonde le \cle{KNDP} (Kamerun National Democratic Party) en février 1955, après avoir quitté le NCNC (National Council of Nigeria and the Cameroons) de Nnamdi Azikiwe, qu'il juge trop nigérianocentrique. Son combat est spécifique : arracher le Southern Cameroons à la tutelle britannique et à la tentation du rattachement au Nigeria, pour le \cle{réunifier} avec le Cameroun francophone. Il plaide cette cause à l'ONU (pétitions 1957, 1958, 1959) et mobilise l'électorat anglophone.

\textbf{Victoire électorale.} Il remporte le \cle{plébiscite du 11 février 1961}, organisé par l'ONU : sur 194 586 votants, \textbf{133 586 (68,6 \%)} choisissent la réunification avec la République du Cameroun, contre 61 001 (31,4 \%) pour le rattachement au Nigeria. (Au Northern Cameroons, le résultat est inverse : 60 \% pour le Nigeria).

\textbf{Accession au pouvoir.} Le 1\textsuperscript{er} octobre 1961, la République fédérale du Cameroun est proclamée : Ahidjo en est le Président, Foncha le \cle{Vice-Président} et Premier ministre du Cameroun occidental (West Cameroon). Il est ainsi l'artisan du \cle{fédéralisme à deux États}, qui garantit provisoirement la spécificité anglophone (common law, système éducatif anglais, administration locale) — jusqu'à l'abolition du fédéralisme par référendum le 20 mai 1972.

\courssub{André-Marie Mbida (1917-1980) : le premier Chef de gouvernement évincé}

\textbf{Origine et formation.} Né en 1917 à Edinding (région du Centre, pays eton), André-Marie Mbida est un catholique pratiquant, avocat de formation (faculté de droit de Paris). Il fonde le \cle{PDC} (Parti des démocrates camerounais) en 1955, formation modérée, pro-occidentale, soutenue par l'administration et l'Église catholique.

\textbf{Action et éviction.} En \textbf{mai 1957}, à la faveur de la loi-cadre Defferre et des premières institutions autonomes, Mbida devient le \cle{premier Chef de gouvernement} du Cameroun (Vice-président du Conseil de gouvernement). Mais il refuse d'engager le Cameroun dans la \emph{Communauté française} proposée par le général de Gaulle (référendum de septembre 1958) et réclame l'indépendance immédiate, ce que la France juge être un « nationalisme excessif » et un risque de basculement dans le bloc de l'Est. Sous la pression de l'administration française (Haut-commissaire Messmer) et de ses alliés à l'ATCAM (députés du Nord, Ahidjo en tête), il est renversé par un vote de défiance en \textbf{février 1958} et remplacé par Ahmadou Ahidjo, jugé plus « fiable » et plus malléable. Le cas Mbida montre que même les nationalistes « modérés » pouvaient être écartés dès qu'ils contrariaient les intérêts stratégiques français.

\courssub{Autres figures du nationalisme (Complément Bac / Approfondissement)}

\begin{itemize}
\item \textbf{Ernest Ouandié} (1924-1971) : dernier chef historique de l'UPC au maquis (région de l'Ouest, puis Est), il poursuit la lutte armée après 1960 contre le régime d'Ahidjo, qu'il considère comme néocolonial. Trahi, arrêté en 1970, il est condamné à mort par un tribunal militaire et \emph{exécuté} publiquement par peloton d'exécution à Bafoussam le \textbf{15 janvier 1971}. Sa mort marque la fin militaire du maquis UPC.
\item \textbf{Osendé Afana} (1930-1966) : économiste brillant (doctorat à Paris, thèse sur l'économie camerounaise), il abandonne une carrière prometteuse à l'ONU pour rejoindre le maquis UPC dans l'Est du pays (région de Moloundou). Il est tué au combat en mars 1966. Il incarne le sacrifice de l'élite intellectuelle engagée.
\item \textbf{Mongo Beti} (1932-2001, Alexandre Biyidi) : écrivain et intellectuel engagé, auteur notamment du \emph{Pauvre Christ de Bomba} (1956, prix Sainte-Beuve). Exilé en France, il dénonce dans ses romans et essais (\emph{Main basse sur le Cameroun}, 1972, censuré par le gouvernement français à la demande d'Ahidjo) la colonisation puis la « néo-colonisation » du Cameroun indépendant.
\item \textbf{Marthe Moumié} (née en 1937) : épouse de Félix-Roland Moumié, résistante et diplomate de l'UPC en exil (Accra, Conakry, Alger). Témoin direct de l'assassinat de son mari à Genève, elle a consacré sa vie à la mémoire de la lutte nationaliste et à la réhabilitation des « maquisards ».
\item \textbf{Charles Assalé} (1911-1999) : figure du Nord, premier Président de l'Assemblée législative (1959), artisan de l'alliance Nord-Sud (Ahidjo-Foncha).
\item \textbf{Solomon Tandeng Muna} (1912-2002) : Premier ministre du Cameroun occidental (1965-1968), Vice-Président de la République fédérale (1970-1972), premier Président de l'Assemblée nationale unifiée (1973).
\end{itemize}

\courssub{Synthèse : deux stratégies face à la décolonisation}

\begin{popfigure}
\begin{center}
\begin{tikzpicture}[scale=0.88, every node/.style={font=\small}]
% Trois blocs principaux
\node[draw, rectangle, fill=popblueL, rounded corners, align=center, text width=3.8cm, minimum height=2.2cm] (UPC) at (0,0) {\textbf{UPC — VOIE RADICALE}\\[3pt] \textbf{Chefs :} Um Nyobé, Moumié, Ouandié, Afana\\[3pt] \textbf{Base :} Paysans, ouvriers, jeunesse, intellectuels de gauche\\[3pt] \textbf{Zone forte :} Sanaga-Maritime, Ouest, Est};
\node[draw, rectangle, fill=popgreenL, rounded corners, align=center, text width=3.8cm, minimum height=2.2cm] (UC) at (7,0) {\textbf{UC / PDC — MODÉRÉS (EST)}\\[3pt] \textbf{Chefs :} Ahidjo, Mbida (évincé), Assalé\\[3pt] \textbf{Base :} Évolués, chefs traditionnels, Nord, fonctionnaires\\[3pt] \textbf{Zone forte :} Nord, Centre, Littoral (villes)};
\node[draw, rectangle, fill=poporangeL, rounded corners, align=center, text width=3.8cm, minimum height=2.2cm] (KNDP) at (14,0) {\textbf{KNDP — MODÉRÉS (OUEST)}\\[3pt] \textbf{Chefs :} Foncha, Muna\\[3pt] \textbf{Base :} Enseignants, élites anglophones, Grassfield\\[3pt] \textbf{Zone forte :} Southern Cameroons (Nord-Ouest, Sud-Ouest)};
% Flèches relations
\draw[<->, very thick, popink, bend left=15] (UPC.east) to node[midway, above, align=center, font=\tiny] {Confrontation armée\\ Répression (1955-1971)} (UC.west);
\draw[<->, very thick, popdark, bend left=15] (UC.east) to node[midway, above, align=center, font=\tiny] {Alliance pour la\\ réunification (1961)\newline Parti unique UNC (1966)} (KNDP.west);
\draw[->, thick, poppurple, dashed, bend right=20] (UPC.south east) to node[midway, below, align=center, font=\tiny] {Tentative d'alliance\\ avortée (1959-60)} (KNDP.south west);
% Synthèse box
\node[align=center, fill=popcream, rounded corners, draw=popline, text width=13cm, minimum height=2.5cm] at (7,-3.2) {\textbf{Synthèse des stratégies nationalistes (1948-1961)}\\[4pt]
\textbf{UPC (Radicale)} : Lutte de masse + pétitions ONU (droit international) + maquis armé (1955-71) + diplomatie panafricaine/tiers-mondiste (Accra, Conakry, Pékin, Alger). Soutien : bloc de l'Est, Non-alignés. \emph{Issue :} Élimination physique des chefs, interdiction, clandestinité.\\[4pt]
\textbf{UC / KNDP (Modérée)} : Négociation avec le tuteur + participation aux institutions (ATCAM, gouvernement autonome) + élections + plébiscite ONU (1961). Soutien : France, Royaume-Uni, bloc occidental. \emph{Issue :} Accession au pouvoir (Ahidjo, Foncha), indépendance 1960, réunification 1961, parti unique 1966.};
\end{tikzpicture}
\end{center}
\legende{Schéma des trois grandes forces nationalistes camerounaises (1948-1961) et de leurs relations. L'UPC radicale est combattue par les modérés soutenus par la France ; l'UC (Est) et le KNDP (Ouest) s'allient pour réaliser la réunification de 1961. La flèche pointillée rappelle la tentative avortée de front unique UPC-KNDP en 1959-1960.}
\end{popfigure}

Ce schéma résume l'essentiel : le nationalisme camerounais n'est pas un bloc unique. Il est traversé par une \cle{fracture stratégique} qui oppose deux voies vers l'indépendance, dans le contexte de la Guerre froide.

\begin{center}
\begin{tabularx}{\linewidth}{|l|X|X|}
\hline
\rowcolor{popblueL} \textbf{Critère de comparaison} & \textbf{Voie radicale (UPC)} & \textbf{Voie modérée (UC, PDC, KNDP)} \\
\hline
Objectif proclamé & Indépendance \emph{totale et immédiate} + réunification & Indépendance \emph{par étapes}, dans la coopération (Communauté/accords) \\
\hline
Moyen d'action principal & Pétitions ONU, mobilisation de masse, grève générale, maquis armé & Négociations bilatérales, participation aux institutions, élections, plébiscite ONU \\
\hline
Soutien extérieur & Ghana, Guinée, Algérie, Chine, URSS, pays non-alignés (bloc de l'Est / Tiers-Monde) & France, puis Grande-Bretagne, USA, Allemagne de l'Ouest (bloc occidental) \\
\hline
Rapport à la France & Rupture : puissance tutélaire illégitime, ennemie principale & Partenariat : puissance tutélaire dont on obtient le départ négocié \\
\hline
Issue politique (1960-61) & Répression, interdiction (1955), chefs assassinés (Moumié 1960) ou exécutés (Ouandié 1971) & Accession au pouvoir : Ahidjo Président (1960), Foncha Vice-Président (1961) \\
\hline
Héritage institutionnel & Mémoire de la « vraie » indépendance, référence pour l'opposition ultérieure & État camerounais actuel (institutions, frontières, administration, armée) \\
\hline
\end{tabularx}
\end{center}

\begin{attention}[Ni héros parfaits, ni traîtres absolus — L'exigence historienne]
Il faut éviter deux écueils symétriques : l'\cle{hagiographie} (faire des nationalistes des saints sans défauts, gommer les divisions ethniques, les erreurs stratégiques, les violences du maquis) et la \cle{diabolisation} (traiter les modérés de simples « collaborateurs » ou « valets du néocolonialisme »). Il faut analyser les choix stratégiques \emph{dans leur contexte} : celui de la \cle{Guerre froide} (1947-1991). L'UPC, soutenue par le bloc soviétique et chinois, apparaît à la France comme une menace communiste au cœur de l'Afrique équatoriale ; l'UC et le KNDP, pro-occidentaux, bénéficient en retour du soutien français et britannique. Chaque camp fait avec les rapports de force de son époque, ses marges de manœuvre et ses contraintes. L'historien juge les stratégies et leurs conséquences, il ne distribue pas les bons points moraux.
\end{attention}

\begin{exoresolu}[Exercice de synthèse : comparer deux figures du nationalisme — Type Bac]
\textbf{Énoncé.} À l'aide de la grille du Dossier 3 (5 rubriques), comparez les parcours de Ruben Um Nyobé et d'Ahmadou Ahidjo, puis dégagez en cinq lignes maximum ce qui oppose fondamentalement leurs conceptions de l'indépendance.
\tcblower
\textbf{Corrigé détaillé.}
\begin{enumerate}
\item \emph{Origine et formation} : Um Nyobé, comptable et syndicaliste bassa (Sanaga-Maritime, zone côtière/forestière) ; Ahidjo, radio-opérateur PTT, fils de chef peul (Nord-Cameroun, zone sahélienne). Tous deux sont des « évolués » formés dans le système colonial, mais aucun n'a fait d'études supérieures en métropole.
\item \emph{Engagement} : Um Nyobé fonde et dirige l'UPC (avril 1948), parti de masse hors institutions ; Ahidjo gravit les échelons institutionnels (ARCAM 1947, ATCAM 1952) et dirige l'UC (1958).
\item \emph{Action majeure} : Um Nyobé plaide trois fois à l'ONU (1952, 1954, 1958) puis organise le maquis (1955-58) ; Ahidjo négocie l'autonomie interne (1957), devient Premier ministre (fév. 1958), proclame l'indépendance (1\textsuperscript{er} janv. 1960).
\item \emph{Fin} : Um Nyobé est tué au maquis le 13 septembre 1958 (à 45 ans) ; Ahidjo accède à la présidence (mai 1960), instaure le parti unique (1966), démissionne en 1982, meurt en exil à Dakar en 1989.
\item \emph{Postérité} : Um Nyobé, longtemps occulté, devient le symbole de l'indépendance « vraie » et de la réunification ; Ahidjo reste le « père de l'indépendance » officielle, mais contesté pour son autoritarisme et la continuité néocoloniale.
\end{enumerate}
\emph{Conclusion (5 lignes max) :} Um Nyobé conçoit l'indépendance comme une \textbf{rupture totale et immédiate} avec la puissance tutélaire, obtenue par la mobilisation souveraine du peuple et, au besoin, les armes ; Ahidjo la conçoit comme un \textbf{processus négocié, par étapes}, préservant la coopération économique, militaire et technique avec la France. Le premier est éliminé physiquement par la puissance coloniale, le second accède au pouvoir avec son appui : l'indépendance camerounaise de 1960 est donc celle de la voie modérée, ce qui nourrit jusqu'à aujourd'hui le débat sur son caractère « incomplet » ou « confisquée ».
\end{exoresolu}

\begin{exercice}[Entraînement Probatoire — Analyse de documents croisés]
\textbf{Document 1 :} Extrait du discours de Ruben Um Nyobé à l'ONU, 17 décembre 1952 : « Nous voulons l'indépendance du Cameroun, non pas dans dix ans, non pas dans cinq ans, mais maintenant. »
\textbf{Document 2 :} Déclaration d'Ahmadou Ahidjo, 1\textsuperscript{er} janvier 1960 : « Le Cameroun accède à la souveraineté internationale dans la paix et l'amitié avec la France. »
\textbf{Questions :}
\begin{enumerate}
\item Identifiez la nature et la date de chaque document.
\item Quel objectif chaque homme assigne-t-il à l'indépendance ?
\item Expliquez pourquoi la France a soutenu l'un et combattu l'autre.
\item En quoi ces deux extraits illustrent-ils la « fracture stratégique » du nationalisme camerounais ?
\end{enumerate}
\end{exercice}

\begin{corrige}[Exercice n°1 — Éléments de corrigé]
\begin{enumerate}
\item Doc 1 : Discours à la tribune de l'ONU (Quatrième Commission), New York, 17 déc. 1952. Doc 2 : Discours de proclamation d'indépendance, Yaoundé, 1\textsuperscript{er} janv. 1960.
\item Um Nyobé : indépendance \emph{immédiate, totale, inconditionnelle} (rupture). Ahidjo : indépendance \emph{négociée, pacifique, dans la continuité} (coopération).
\item La France combat Um Nyobé (UPC) car il menace ses intérêts stratégiques (ressources, bases, influence) et s'allie au bloc de l'Est ; elle soutient Ahidjo (UC) car il garantit la préservation de ces intérêts via les accords de coopération.
\item Ils illustrent les deux pôles : rupture vs continuité, masse vs élites, armes vs urnes, bloc de l'Est vs bloc occidental. La voie modérée l'emporte institutionnellement en 1960, mais la revendication radicale persiste en mémoire.
\end{enumerate}
\end{corrige}

\begin{aretenir}[L'essentiel du Dossier 3 — Fiche de révision Bac]
\begin{itemize}
\item \textbf{Ruben Um Nyobé} (UPC) : le « Mpodol », pétitionnaire à l'ONU (1952, 1954, 1958), organisateur du maquis, tué le 13 septembre 1958 près de Boumnyebel.
\item \textbf{Félix-Roland Moumié} (UPC) : médecin, internationalise la lutte (Accra, Conakry, Alger, Pékin), empoisonné au thallium à Genève le 3 novembre 1960 par le SDECE.
\item \textbf{Ahmadou Ahidjo} (UC) : négociateur, Premier ministre (fév. 1958), proclame l'indépendance le 1\textsuperscript{er} janvier 1960, premier Président, parti unique UNC (1966).
\item \textbf{John Ngu Foncha} (KNDP) : instituteur anglophone, vainqueur du plébiscite du 11 février 1961 (68,6 \% pour la réunification), Vice-Président de la République fédérale.
\item \textbf{André-Marie Mbida} (PDC) : premier Chef de gouvernement (mai 1957), évincé par la France en fév. 1958 pour « nationalisme excessif » (refus de la Communauté).
\item \textbf{Ernest Ouandié} : dernier chef maquis UPC, exécuté le 15 janvier 1971 à Bafoussam (fin du maquis).
\item La fracture centrale oppose la voie \cle{radicale} (lutte armée, indépendance totale, panafricanisme) à la voie \cle{modérée} (négociation, indépendance par étapes, pro-occidental), dans le contexte de la Guerre froide.
\end{itemize}
\end{aretenir}

\begin{savaistu}[Débat historiographique actuel : « L'indépendance confisquée » ?]
Depuis les années 1990 (conférence nationale avortée, multipartisme), le débat public camerounais oppose deux lectures : (1) L'indépendance de 1960 est une \emph{indépendance confisquée} par la France et ses alliés modérés (Ahidjo), qui ont éliminé les vrais patriotes (Um Nyobé, Moumié, Ouandié) pour installer un néocolonialisme durable. (2) L'indépendance de 1960 est une \emph{réussite pragmatique} : Ahidjo a obtenu la souveraineté internationale, l'unité territoriale (réunification 1961) et la stabilité, évitant le chaos de guerres civiles type Congo (1960) ou Angola (1975). L'historien ne tranche pas : il montre comment ce débat structure encore la vie politique camerounaise (revendications anglophones, mémoire des maquisards, relations France-Cameroun).
\end{savaistu}

Ces portraits croisés préparent la section suivante : c'est bien l'affrontement entre ces stratégies — et la victoire institutionnelle de la voie modérée — qui détermine le \cle{cheminement institutionnel} vers les indépendances de 1960 et la réunification de 1961, que nous allons maintenant étudier en détail.

\coursec{Le cheminement institutionnel vers les indépendances et la réunification (1956-1961)}

Après avoir analysé les figures emblématiques du nationalisme camerounais — de l'action syndicale d'Uma Ndoumbe à la radicalité de l'UPC en passant par le pragmatisme d'Ahidjo et de Foncha — nous devons désormais comprendre comment ces acteurs, dans un cadre international contraignant (tutelle ONU) et un contexte de guerre froide, ont négocié, imposé ou subi les étapes juridiques et institutionnelles qui ont conduit à la naissance de l'État camerounais moderne. Ce cheminement n'est pas une ligne droite : il oscille entre transfert de compétences progressif (loi-cadre Defferre), répression violente (guerre du maquis), souveraineté formelle (1$^{\text{er}}$ janvier 1960) et construction fédérale inédite (1$^{\text{er}}$ octobre 1961).

\courssub{La loi-cadre Defferre (23 juin 1956) : le tournant de l'autonomie interne}

La loi-cadre Defferre, du nom du ministre de la France d'Outre-mer Gaston Defferre, constitue la \cle{charte fondatrice} de l'autonomie interne des territoires d'outre-mer français. Pour le Cameroun sous tutelle française, elle marque la fin du régime de l'indigénat (déjà aboli en 1946 mais dont les séquelles persistent) et surtout la fin du \cle{dualisme électoral} (collèges distincts pour citoyens français et « sujets » africains), remplacé par le \cle{suffrage universel direct}.

\begin{definition}[Loi-cadre Defferre (1956)]
Ensemble de lois organiques votées le 23 juin 1956 qui accordent l'autonomie interne aux territoires d'outre-mer français : création d'assemblées territoriales élues au suffrage universel, formation d'un conseil de gouvernement présidé par le Haut-commissaire mais responsable devant l'assemblée, et transfert de compétences (éducation, santé, économie locale) aux élus locaux. La défense, la police, les finances et les affaires étrangères restent sous contrôle français.
\end{definition}

\begin{popfigure}
\begin{tikzpicture}[scale=0.7, every node/.style={font=\small}]
\node[draw, ellipse, fill=red!20, minimum width=3cm, minimum height=1.2cm, align=center] (Start) at (0,0){Loi-cadre Defferre\\23 juin 1956};
\node[draw, rectangle, fill=yellow!20, rounded corners, minimum width=3.5cm, minimum height=1.5cm, align=center] (E1) at (5,2){Élections législatives\\déc. 1956 / mars 1957\\\textbf{Victoire UC + alliés}\\\textbf{Ahidjo PM (1958)}};
\node[draw, rectangle, fill=yellow!20, rounded corners, minimum width=3.5cm, minimum height=1.5cm, align=center] (E2) at (5,-2){Interdiction UPC\\13 juil. 1955 (Messmer)\\\textbf{Guerre du maquis}\\\textbf{Répression militaire}};
\node[draw, ellipse, fill=green!20, minimum width=3cm, minimum height=1.2cm, align=center] (Indep) at (10,0){Indépendance\\1$^{\text{er}}$ janv. 1960};
\node[draw, diamond, fill=cyan!20, aspect=2, minimum width=3.5cm, align=center] (Pleb) at (14,2){Plébiscite ONU\\11 fév. 1961\\\textbf{Sud : Réunification (70\%)}\\\textbf{Nord : Nigeria (60\%)}};
\node[draw, ellipse, fill=blue!20, minimum width=3.5cm, minimum height=1.5cm, align=center] (Fed) at (18,0){République Fédérale\\1$^{\text{er}}$ oct. 1961\\\textbf{Conf. Foumban (juil. 1961)}};
\draw[->, thick, popblue] (Start) -- node[above, sloped, pos=0.5] {Cadre légal} (E1);
\draw[->, thick, poporange] (Start) -- node[below, sloped, pos=0.5] {Contexte répressif} (E2);
\draw[->, thick, popblue] (E1) -- node[above, sloped, pos=0.5] {Transfert compétences} (Indep);
\draw[->, thick, poporange] (E2) -- node[below, sloped, pos=0.5] {Sécurité -> Ahidjo (1959)} (Indep);
\draw[->, thick, popgreen] (Indep) -- node[above, sloped, pos=0.5] {Adhésion ONU sept. 1960} (Pleb);
\draw[->, thick, poppurple] (Pleb) -- node[above, sloped, pos=0.5] {Résultat Sud} (Fed);
\end{tikzpicture}
\legende{Chronologie synthétique du cheminement institutionnel (1956-1961) : du cadre légal français à la fédération, en passant par l'alternative répressive (UPC) et le choix plébiscitaire du British Southern Cameroons.}
\end{popfigure}

\courssub{Les élections de 1957 et l'accession au pouvoir d'Ahmadou Ahidjo}

L'application de la loi-cadre entraîne les premières élections à l'Assemblée législative du Cameroun (ALCM) les 23 décembre 1956 et 4 mars 1957 (second tour). L'\cle{Union Camerounaise (UC)}, fondée en 1958 par fusion de plusieurs groupes modérés (dont le Bloc démocratique camerounais de Louis-Paul Aujoulat et des éléments de l'UC d'Ahidjo), remporte une majorité relative. André-Marie Mbida, chef de file des démocrates-chrétiens, devient le premier Premier ministre (février 1957). Mais son refus de réclamer l'indépendance immédiate et sa gestion autoritaire provoquent une motion de censure. Le 17 mai 1958, \cle{Ahmadou Ahidjo} (vice-Premier ministre, ministre de l'Intérieur) prend la tête du gouvernement. Il obtient de Paris le transfert progressif des compétences de sécurité (police, gendarmerie) en 1959, crucial pour mater la rébellion de l'UPC.

\begin{exoresolu}[Analyse du basculement Mbida -> Ahidjo (1958)]
\textbf{Document :} Extrait du discours d'Ahmadou Ahidjo à l'Assemblée législative, mai 1958 : « \emph{Nous voulons l'indépendance, mais nous la voulons dans l'ordre, la paix et la fraternité. Nous ne transigerons pas avec ceux qui, sous couvert de nationalisme, sèment la mort et la désolation.} »

\textbf{Question :} Expliquer pourquoi Ahidjo succède à Mbida en 1958 et quelle stratégie il met en œuvre.

\textbf{Solution détaillée :}
\begin{enumerate}
\item \textbf{Contexte :} Mbida (Premier ministre depuis fév. 1957) refuse de fixer une date d'indépendance, craignant le chaos. Il s'aliène l'UC et la France.
\item \textbf{Déclencheur :} Motion de censure votée le 10 mai 1958 par la majorité UC + alliés (soutenus par le Haut-commissaire Pierre Messmer), reprochant à Mbida son immobilisme face à la rébellion UPC et son autoritarisme.
\item \textbf{Profil d'Ahidjo :} Nordiste (Garoua), musulman, fonctionnaire, perçu comme un modéré capable de rassurer la France (sécurité) et de fédérer au-delà des clivages ethniques/régionaux (Sud/Bamiléké vs Nord).
\item \textbf{Stratégie « Paix et Ordre » :} 
\begin{itemize}
\item Négociation avec la France : transfert de la Sûreté nationale (1959) -> contrôle de l'appareil répressif.
\item Marginalisation de l'UPC : interdite depuis 1955, ses leaders (Um Nyobé tué en 1958, Ouandié, Kingué) en maquis ou exil.
\item Revendication de l'indépendance « négociée » pour le 1$^{\text{er}}$ janv. 1960, acceptée par de Gaulle (Communauté française).
\end{itemize}
\end{enumerate}
\end{exoresolu}

\courssub{L'interdiction de l'UPC (1955) et la guerre du maquis : le coût de la répression}

Parallèlement à la voie institutionnelle, la répression de l'Union des Populations du Cameroun (UPC) structure la période. Interdite par arrêté du Haut-commissaire Pierre Messmer le 13 juillet 1955 (après des émeutes en mai 1955), l'UPC bascule dans la clandestinité et la lutte armée (\cle{maquis}) dès 1956 sous l'impulsion de Ruben Um Nyobé, puis d'Ernest Ouandié et Abel Kingué après l'assassinat d'Um Nyobé par l'armée française (13 septembre 1958).

\begin{attention}[Piège chronologique : Interdiction vs Indépendance]
L'UPC est interdite \textbf{avant} la loi-cadre Defferre (juillet 1955 vs juin 1956) et \textbf{bien avant} l'indépendance (1960). La « guerre du maquis » (1955-1964, avec un pic 1958-1960) se déroule donc \textbf{sous administration française} (Haut-commissaire Messmer puis Bolotte) puis sous le gouvernement Ahidjo (dès 1958 pour la police, 1960 pour l'armée). Ne pas écrire : « L'UPC s'est rebellée après l'indépendance ». Elle a été interdite puis a pris le maquis \emph{pendant} la période d'autonomie interne.
\end{attention}

Cette répression, menée par l'armée française (opérations « Mouette », « Épervier ») puis par les forces camerounaises encadrées par des officiers français, fait des dizaines de milliers de morts (estimations : 30\,000 à 100\,000), principalement dans le Sanaga-Maritime, le Wouri et le Moungo (pays Bassa, Bamiléké). Elle laisse un traumatisme durable et une mémoire conflictuelle (voir Section 5).

\courssub{L'indépendance du Cameroun français (1$^{\text{er}}$ janvier 1960) : souveraineté et choix de la Communauté}

Le 1$^{\text{er}}$ janvier 1960, le Cameroun français accède à la \cle{souveraineté internationale}. Ahmadou Ahidjo proclame l'indépendance à Yaoundé. Le 20 septembre 1960, le Cameroun devient membre de l'ONU (résolution 1476). Le choix politique majeur est l'adhésion à la \cle{Communauté française} (instituée par la Constitution de la V$^{\text{e}}$ République, 1958), contrairement à la Guinée de Sékou Touré (1958) qui a choisi l'indépendance immédiate hors Communauté.

\begin{propriete}[Originalité de l'indépendance camerounaise (1960)]
Le Cameroun est le \textbf{premier territoire d'Afrique subsaharienne française} à accéder à l'indépendance (1$^{\text{er}}$ janv. 1960), avant le « raz-de-marée » de 1960 (Sénégal, Mali, Côte d'Ivoire, etc. en août-septembre 1960). Il le fait \textbf{dans le cadre de la Communauté française}, conservant des accords de coopération militaire, monétaire (Franc CFA), technique et culturelle très étroits avec la France.
\end{propriete}

La Constitution du 21 février 1960 instaure une \cle{République} à régime \cle{présidentiel} : le Président de la République (Ahidjo, élu par l'Assemblée nationale) est chef de l'État, chef du gouvernement et commandant en chef des armées. Le Parlement est monocaméral (Assemblée nationale). Ce texte, bien que démocratique en apparence, concentre les pouvoirs entre les mains du Président, préfigurant le régime autoritaire qui s'installe dès 1962 (parti unique) et 1966 (fusion UC-KNDP -> UNC).

\courssub{La question du British Cameroons : le plébiscite de février 1961}

L'indépendance du Cameroun français laisse en suspens le sort des \cle{British Cameroons} (Nord et Sud), toujours sous tutelle britannique et administrés comme partie intégrante du Nigeria (provinces du Nord et de l'Est). L'ONU impose un \cle{plébiscite} (résolution 1352 XIV, oct. 1959) pour permettre aux populations de choisir leur avenir : indépendance par réunion au Cameroun indépendant ou indépendance par intégration au Nigeria (indépendant depuis le 1$^{\text{er}}$ oct. 1960).

\begin{methode}[Justifier le vote du Southern Cameroons pour la réunification (1961)]
Pour analyser le résultat (70\% pour la réunification), croiser quatre facteurs explicatifs :
\begin{enumerate}
\item \textbf{Peur de l'absorption par le Nigeria :} Démographie écrasante du Nigeria (35 millions vs 1 million), domination économique (ports, chemins de fer, administration), risque de marginalisation politique des minorités anglophones dans une fédération nigériane à trois régions (Nord, Ouest, Est).
\item \textbf{Affinités culturelles et historiques transfrontalières :} Peuples Bamiléké, Bamoun, Tikar, Bassa, Douala coupés par la frontière de 1916. Réseaux commerciaux, parentés, chefferies traditionnelles (Fon, Sultan) traversent la ligne de partage.
\item \textbf{Promesses fédérales d'autonomie (Foumban anticipé) :} Le KNDP (John Ngu Foncha) négocie et promet un statut fédéral garantissant : Parlement propre (House of Assembly), Premier ministre propre, judiciaire propre (Common Law), système éducatif propre. Ahidjo accepte ce cadre pour obtenir le « Oui ».
\item \textbf{Leadership de Foncha et du KNDP :} Foncha (Premier ministre du Southern Cameroons depuis 1959) mène une campagne efficace (« \emph{One Cameroon, One People, One Destiny} »), marginalisant l'option « Nigeria » (KNC - Kamerun National Congress, soutenu par le Nord) et l'option « Indépendance pure » (marginale).
\end{enumerate}
\end{methode}

\begin{center}
\begin{tabularx}{\linewidth}{|l|X|X|}
\hline
\rowcolor{popblueL}
\textbf{Territoire} & \textbf{Question posée} & \textbf{Résultat (11 fév. 1961)} \\
\hline
\textbf{Northern Cameroons} & Intégration au Nigeria OU Réunification au Cameroun & \textbf{60\% pour le Nigeria} (Intégration effective 1$^{\text{er}}$ juin 1961) \\
\hline
\textbf{Southern Cameroons} & Intégration au Nigeria OU Réunification au Cameroun & \textbf{70\% pour la Réunification} (Indépendance 1$^{\text{er}}$ oct. 1961) \\
\hline
\end{tabularx}
\end{center}

\courssub{La Conférence de Foumban (juillet 1961) : négocier la Constitution fédérale}

Le vote du Sud impose la tenue d'une conférence constitutionnelle pour définir les modalités de l'union. Elle se tient à \cle{Foumban} (chef-lieu du département du Noun, région de l'Ouest, terre bamoun, symbolique pour le Sultan Njoya et l'Histoire) du 17 au 21 juillet 1961. Deux délégations s'affrontent/négocient :
\begin{itemize}
\item \textbf{Délégation de la République du Cameroun (East Cameroon) :} Menée par Ahidjo, forte de 30 délégués. Position : État unitaire décentralisé, présidence forte, intégration rapide.
\item \textbf{Délégation du Southern Cameroons (West Cameroon) :} Menée par Foncha, 20 délégués. Position : Fédération forte, deux États autonomes (législatif, exécutif, judiciaire), présidence tournante ou vice-présidence forte pour le VP (Foncha), maintien du Common Law et de l'anglais.
\end{itemize}

\begin{definition}[Réunification (1961)]
Processus politique et juridique par lequel le Southern Cameroons britannique, après un plébiscite organisé par l'ONU (11 fév. 1961), rejoint la République du Cameroun (indépendante depuis le 1$^{\text{er}}$ janv. 1960) pour former la \textbf{République Fédérale du Cameroun} (1$^{\text{er}}$ oct. 1961). Elle résulte de la Conférence de Foumban (juil. 1961) qui négocie la Constitution fédérale.
\end{definition}

\begin{propriete}[Unicité du modèle camerounais]
La réunification camerounaise est un \textbf{cas unique en Afrique} : deux territoires colonisés par des puissances différentes (France et Royaume-Uni), placés sous tutelle de l'ONU, dont l'un accède à l'indépendance avant l'autre, et qui choisissent \textbf{librement} (par plébiscite) de s'unir dans une structure fédérale, avant de fusionner ultérieurement (1972).
\end{propriete}

Le compromis de Foumban accouche de la Constitution fédérale promulguée le 1$^{\text{er}}$ septembre 1961, entrée en vigueur le 1$^{\text{er}}$ octobre 1961 :
\begin{itemize}
\item \textbf{État fédéral :} République Fédérale du Cameroun.
\item \textbf{Deux États fédérés :} \cle{East Cameroon} (ex-Cameroun français, 6 régions) et \cle{West Cameroon} (ex-Southern Cameroons, 2 divisions : Bamenda, Kumba). Chaque État a : un Parlement (Assemblée législative), un Premier ministre, un pouvoir judiciaire (Common Law à l'Ouest, Droit civil à l'Est).
\item \textbf{Exécutif fédéral :} Président de la République (Ahidjo) — élu au suffrage universel direct — chef de l'État, du gouvernement, des armées. Vice-Président de la République (Foncha) — président de l'Assemblée législative du West Cameroon — assure l'intérim.
\item \textbf{Législatif fédéral :} Assemblée nationale fédérale (députés élus au suffrage universel direct, proportionnelle aux populations).
\item \textbf{Partage des compétences :} Fédéral (défense, diplomatie, monnaie, justice suprême, éducation supérieure, transports internationaux) ; États (éducation primaire/secondaire, santé, administration locale, justice courante, culture, développement rural).
\end{itemize}

\begin{savaistu}[Le drapeau fédéral (1961-1975)]
Le drapeau de la République Fédérale du Cameroun (loi n°61-24 du 1$^{\text{er}}$ nov. 1961) reprenait le tricolore \textbf{vert-rouge-jaune} (vertical) avec \textbf{deux étoiles d'or} à cinq branches placées dans la bande verte (côté hampe). Ces deux étoiles symbolisaient les \textbf{deux États fédérés} : East Cameroon et West Cameroon. En 1975 (loi n°75-9), après le référendum du 20 mai 1972 instituant l'État unitaire, la deuxième étoile est supprimée, remplacée par une unique étoile d'or au centre de la bande rouge, symbolisant l'unité nationale.
\end{savaistu}

\courssub{Bilan : une construction à deux vitesses (1960-1961)}

Le cheminement institutionnel révèle une asymétrie fondamentale :
\begin{enumerate}
\item \textbf{Le Cameroun français} accède à l'indépendance par un \textbf{transfert de souveraineté négocié avec la France} (loi-cadre -> élections -> gouvernement Ahidjo -> proclamation 1$^{\text{er}}$ janv. 1960). La légitimité repose sur le suffrage universel (élections 1957, 1960) et la reconnaissance internationale (ONU).
\item \textbf{Le Southern Cameroons} accède à l'indépendance \textbf{par la réunification} (plébiscite ONU 11 fév. 1961 -> Conférence de Foumban -> Constitution fédérale 1$^{\text{er}}$ oct. 1961). Sa légitimité repose sur l'autodétermination (plébiscite) et le contrat fédéral (Foumban).
\end{enumerate}

Cette dualité d'origine — indépendance \emph{par le haut} (transfert colonial français) vs indépendance \emph{par le bas} (choix plébiscitaire britannique) — structure les tensions politiques futures : centralisation vs fédéralisme, droit civil vs Common Law, francophonie vs anglophonie, qui resurgiront violemment à partir de 2016 (crise anglophone).

\begin{exoresolu}[Commentaire de document : Extrait du discours d'Ahidjo à Foumban, 19 juillet 1961]
\textbf{Document :} « \emph{Nous ne voulons pas d'une fédération qui soit une simple juxtaposition de deux États souverains. Nous voulons une fédération forte, au service de l'unité nationale. Le Président fédéral doit avoir les moyens de sa politique. Le Vice-Président sera le garant des intérêts de l'Ouest au sein de l'exécutif fédéral.} »

\textbf{Question :} Ce extrait révèle-t-il la conception unitaire ou fédérale d'Ahidjo ? Justifiez par deux arguments tirés du texte et du contexte.

\textbf{Solution détaillée :}
\begin{enumerate}
\item \textbf{Conception unitaire (centralisatrice) :} « \emph{Nous ne voulons pas d'une fédération qui soit une simple juxtaposition de deux États souverains} » -> Rejet de la fédération confédérale (État fort, États fédérés faibles). « \emph{Le Président fédéral doit avoir les moyens de sa politique} » -> Concentration de l'exécutif fédéral (Ahidjo) au détriment de l'autonomie des États (West Cameroon).
\item \textbf{Concession fédérale formelle :} « \emph{Le Vice-Président sera le garant des intérêts de l'Ouest} » -> Acceptation du principe de la vice-présidence (Foncha) et de la représentation de l'Ouest à la tête de l'État, marqueur fédéral.
\item \textbf{Contexte :} Ahidjo dispose de la majorité démographique (East Cameroon ~ 4,5 millions vs West Cameroon ~ 700 000) et du contrôle de l'armée/administration. La Constitution de 1961 lui donne les pleins pouvoirs (décrets-lois, état d'urgence). La fédération est dès l'origine « à géométrie variable » : fédérale sur le papier, unitaire dans la pratique (fusion partis 1966, fin fédéralisme 1972).
\end{enumerate}
\end{exoresolu}

\begin{attention}[Piège Bac : Les dates commémoratives]
Ne pas confondre les trois dates majeures :
\begin{itemize}
\item \textbf{1$^{\text{er}}$ janvier 1960} : Indépendance du Cameroun français (proclamation Ahidjo). \textbf{Pas la fête nationale.}
\item \textbf{1$^{\text{er}}$ octobre 1961} : Réunification / Indépendance du Southern Cameroons / Naissance de la Fédération. \textbf{Pas la fête nationale.}
\item \textbf{20 mai 1972} : Référendum sur l'État unitaire (Oui massif). \textbf{C'est la date de la Fête nationale actuelle (20 mai).}
\end{itemize}
L'examinateur peut demander : « Quelle événement commémore le 20 mai ? » -> Réponse : Le référendum de 1972 (passage à l'État unitaire), \emph{pas} l'indépendance ni la réunification.
\end{attention}

\coursec{Héritages, mémoires et enjeux historiographiques (Complément Bac / Éducation à la citoyenneté)}

Nous venons de voir, dans la section précédente, comment le Cameroun est passé de la tutelle de l'ONU à l'indépendance (1960), puis à la réunification (1961). Mais cette marche institutionnelle ne s'est pas faite sans violence ni sans débats. La répression de l'UPC, la guerre du maquis, l'assassinat des leaders nationalistes et l'abandon du fédéralisme en 1972 ont laissé des traces profondes dans la mémoire collective camerounaise. Cette dernière section étudie comment cette histoire douloureuse a été racontée, cachée, puis redécouverte, et pourquoi elle reste un enjeu citoyen majeur aujourd'hui. C'est un complément précieux pour le Baccalauréat technique, car il entraîne à la démarche de l'historien : croiser les sources, comparer les versions, justifier ses conclusions.

\courssub{La « Guerre cachée » (1955-1971) : un bilan humain lourd}

\begin{definition}[Guerre cachée]
On appelle \cle{guerre cachée} la guerre de répression menée entre 1955 et 1971, d'abord par l'administration coloniale française puis par le gouvernement camerounais d'Ahmadou Ahidjo (avec l'appui militaire français), contre les militants et sympathisants de l'\cle{UPC} (Union des populations du Cameroun), interdite en juillet 1955. Elle est dite « cachée » parce qu'elle a été longtemps passée sous silence dans les manuels scolaires et le discours officiel.
\end{definition}

Après les émeutes de mai 1955 et l'interdiction de l'UPC, une partie des militants upécistes passe dans la clandestinité et organise le \cle{maquis}, c'est-à-dire une résistance armée dans les forêts et les montagnes, surtout dans deux régions : le \cle{pays Bassa} (département du Sanaga-Maritime, autour d'Édéa et de Douala) et le \cle{pays Bamiléké} (région de l'Ouest). Le conflit dure bien après l'indépendance : il ne prend fin qu'avec la mort des derniers chefs maquisards, notamment Ernest Ouandié, arrêté en 1970 et exécuté le 15 janvier 1971 à Bafoussam.

Le bilan humain est très difficile à établir précisément, car les archives sont longtemps restées fermées. Les historiens proposent des estimations qui varient de plusieurs dizaines de milliers à plus de \num{100000} morts, selon les sources et les méthodes de calcul. Cette fourchette très large est elle-même un enseignement : quand un État ne compte pas ses victimes, l'historien doit travailler avec des témoignages, des registres administratifs partiels et des recoupements.

La répression a aussi provoqué d'importants \cle{déplacements de populations}. Pour isoler les maquisards de leurs soutiens, les autorités ont créé des \cle{villages de regroupement} : les habitants des campagnes étaient déplacés de force et concentrés dans des villages surveillés, le long des axes contrôlés par l'armée. Des centaines de milliers de personnes, notamment en pays Bamiléké et Bassa, ont ainsi été déracinées, perdant leurs champs, leurs cases et leurs cimetières familiaux.

Ces violences ont produit des \cle{traumatismes intergénérationnels} : les familles des victimes ont transmis le souvenir de la répression à leurs enfants et petits-enfants, souvent dans le silence ou la peur, car parler de l'UPC était longtemps dangereux. Les régions Bamiléké et Bassa gardent ainsi une mémoire vive et douloureuse de cette période, qui explique en partie la sensibilité particulière de ces populations aux questions de justice et de reconnaissance.

\begin{attention}[Sourcer ses affirmations]
Le programme exige de « rédiger et justifier une démarche ». Au Baccalauréat, toute affirmation sur les violences doit être \cle{sourcée} : écrivez par exemple « selon les rapports de la mission de visite de l'ONU de 1958 », « selon les archives militaires françaises partiellement déclassifiées », ou « selon les témoignages recueillis par les historiens ». Une affirmation non sourcée sur un sujet sensible est une faute de méthode qui coûte des points.
\end{attention}

\courssub{La construction de l'historiographie : trois versions d'une même histoire}

\begin{definition}[Historiographie]
L'\cle{historiographie} est l'étude de la manière dont l'histoire a été écrite : par qui, dans quel but, avec quelles sources. Elle ne raconte pas les faits eux-mêmes, mais compare les différents récits des faits. Ici, on oppose l'\cle{histoire officielle} (écrite par le pouvoir) à l'histoire « par le bas » (oralité des familles, archives de l'UPC, mémoires des anciens maquisards).
\end{definition}

Sur la guerre du maquis, on peut distinguer trois grandes façons d'écrire l'histoire, qui se sont succédé et affrontées.

\textbf{1. La version officielle (Ahidjo, puis Biya).} Du point de vue du pouvoir, les maquisards de l'UPC étaient des « bandits », des « rebelles » ou des « terroristes » qu'il fallait « pacifier ». Le mot même de \cle{pacification}, hérité du vocabulaire colonial, présente la répression comme une opération de maintien de l'ordre, et non comme une guerre. Dans cette version, l'indépendance de 1960 est un aboutissement pacifique obtenu par la négociation, et l'UPC armée est un obstacle à la construction nationale. Cette version a longtemps dominé les manuels scolaires et les discours officiels.

\textbf{2. La mémoire UPC.} Du point de vue des militants et des familles de victimes, il s'agit d'une \cle{guerre de libération nationale} : l'UPC, premier parti à réclamer l'indépendance immédiate et la réunification, aurait été écrasée par la France pour installer au pouvoir des dirigeants « dociles ». Ses leaders assassinés — Ruben Um Nyobé (tué par l'armée française le 13 septembre 1958), Félix Moumié (empoisonné à Genève en novembre 1960), Ernest Ouandié (exécuté en 1971) — sont présentés comme des \cle{martyrs} de l'indépendance. Cette mémoire s'est transmise par l'oralité, les chansons, les veillées familiales et les réseaux de la diaspora.

\textbf{3. L'approche universitaire récente.} Depuis les années 1990-2000, des historiens camerounais et français ont pu travailler sur de nouvelles sources : \cle{archives françaises partiellement déclassifiées} (militaires, diplomatiques), rapports des missions de l'ONU, témoignages oraux recueillis dans les villages. Ces travaux confirment l'ampleur de la répression (exécutions sommaires, villages de regroupement, bombardements) et montrent le rôle direct de l'armée française après 1960. L'historien ne choisit pas un « camp » : il établit les faits à partir des sources croisées.

\begin{methode}[Traiter un sujet sensible en dissertation]
Pour traiter un sujet comme la guerre du maquis au Baccalauréat, suivez quatre étapes :
\begin{enumerate}
\item \cle{Contextualiser} : situer le conflit dans la Guerre froide et la décolonisation (la France craint un Cameroun « communiste » et veut garder son influence).
\item \cle{Présenter les thèses antagonistes sans parti pris} : exposer la version officielle (« pacification ») et la mémoire UPC (« guerre de libération ») avec neutralité.
\item \cle{Croiser les sources} : archives administratives, témoignages oraux, rapports de l'ONU, travaux universitaires.
\item \cle{Conclure sur l'apaisement mémoriel} : montrer que la reconnaissance des faits est une condition de la réconciliation nationale.
\end{enumerate}
\end{methode}

\courssub{Justice et réparation : une reconnaissance tardive et inachevée}

La reconnaissance officielle des crimes de la guerre cachée a été lente et partielle.

\textbf{Les demandes de reconnaissance.} Des associations de familles de victimes et des historiens réclament la vérité sur l'assassinat de Félix Moumié (empoisonné au thallium à Genève en 1960, crime attribué aux services secrets français), sur la mort de Ruben Um Nyobé (tué par l'armée française en septembre 1958, son corps exposé pour décourager les partisans), et plus largement sur les crimes de guerre commis pendant la « pacification ». Une partie de ces demandes vise la France, dont l'armée et les services ont joué un rôle central.

\textbf{La restitution des archives.} Une partie des archives militaires et diplomatiques françaises sur le Cameroun reste classée « Secret Défense ». Les historiens et les associations demandent leur déclassification complète et leur restitution, car sans archives, impossible d'établir précisément les responsabilités et le nombre des victimes.

\textbf{La réhabilitation des maquisards.} La loi n°91/022 du 16 décembre 1991, relative aux mérites nationaux, a permis de réhabiliter officiellement Ruben Um Nyobé, Félix Moumié et Ernest Ouandié comme \cle{héros nationaux}. Cette réhabilitation, plus de trente ans après les faits, marque une évolution : les « bandits » d'hier sont reconnus comme des acteurs de l'indépendance. Mais beaucoup estiment qu'elle reste incomplète : les maquisards ordinaires, les familles déplacées et les villages détruits n'ont reçu ni réparation ni reconnaissance équivalente.

\begin{savaistu}[Une réhabilitation tardive]
La loi n°91/022 du 16 décembre 1991 réhabilite Ruben Um Nyobé, Félix Moumié et Ernest Ouandié comme « héros nationaux ». Pourtant, plus de trente ans après, une partie des archives militaires françaises sur la « pacification » du Cameroun demeure classée « Secret Défense », ce qui freine le travail des historiens et les demandes de justice des familles de victimes.
\end{savaistu}

\courssub{Un enjeu citoyen actuel : diversité, fédéralisme et crise anglophone}

L'histoire de la marche vers l'indépendance n'est pas seulement un sujet de mémoire : elle éclaire directement les débats politiques actuels du Cameroun.

\textbf{Fédéralisme ou État unitaire ?} La réunification de 1961 s'était faite sur la base d'un État \cle{fédéral} : la République fédérale du Cameroun comprenait deux États fédérés (Cameroun oriental, francophone ; Cameroun occidental, anglophone), chacun avec ses institutions. Le référendum du 20 mai 1972 a mis fin au fédéralisme et instauré la République unie du Cameroun, État unitaire et centralisé. Pour beaucoup d'anglophones, cette évolution a été vécue comme une absorption et une marginalisation de leur spécificité (langue, système juridique de Common Law, système éducatif).

\textbf{La crise anglophone (depuis 2016).} En 2016, des avocats et des enseignants des régions du Nord-Ouest et du Sud-Ouest (NOSO) protestent contre la « francophonisation » des tribunaux et des écoles. La répression des manifestations transforme la revendication en crise politique majeure : grèves, « villes mortes », puis affrontements armés entre forces de sécurité et groupes séparatistes. De nombreux analystes y voient une \cle{résurgence des frustrations nées de Foumban} : la conférence de Foumban (juillet 1961), qui devait fixer les termes du fédéralisme, est jugée par les anglophones comme déséquilibrée, et le non-respect ultérieur de l'esprit fédéral alimente la rancœur. Le Grand Dialogue National de 2019 a abouti à l'octroi d'un \cle{statut spécial} aux régions du Nord-Ouest et du Sud-Ouest, dans le cadre de la décentralisation prévue par la Constitution de 1996.

\textbf{La place de la diaspora.} La diaspora camerounaise (Europe, Amérique du Nord, Afrique) joue un rôle croissant : elle finance des associations mémorielles, fait circuler les travaux historiques sur l'UPC et la guerre cachée, et pèse dans les débats politiques par les réseaux numériques. Elle est à la fois un relais de la mémoire et un acteur de la vie citoyenne.

\begin{popfigure}
\begin{center}
\begin{tikzpicture}[scale=0.72, every node/.style={font=\small}]
\draw[thick, -{Stealth}] (0,0) -- (11.5,0) node[right]{Années};
\foreach \x/\lab in {0/1960, 2/1972, 4/1982, 6/1996, 8/2016, 10/2019}
  \draw (\x,0.12) -- (\x,-0.12) node[below]{\lab};
\node[align=center, fill=popgreenL, rounded corners, inner sep=2pt] at (1,1.1) {Indépendance 1960\\Réunification 1961\\État fédéral};
\draw[popgreen, thick] (1,0.15) -- (1,0.7);
\node[align=center, fill=poporangeL, rounded corners, inner sep=2pt] at (3,1.7) {Référendum 1972\\Fin du fédéralisme\\État unitaire};
\draw[poporange, thick] (3,0.15) -- (3,1.25);
\node[align=center, fill=popblueL, rounded corners, inner sep=2pt] at (5,1.1) {Ère Biya (1982)\\Constitution 1996\\Décentralisation};
\draw[popblue, thick] (5,0.15) -- (5,0.7);
\node[align=center, fill=poppinkL, rounded corners, inner sep=2pt] at (7,1.7) {Crise anglophone 2016\\Villes mortes\\Revendications NOSO};
\draw[poppink, thick] (7,0.15) -- (7,1.25);
\node[align=center, fill=poppurpleL, rounded corners, inner sep=2pt] at (9.5,1.1) {Grand Dialogue\\National 2019\\Statut spécial NOSO};
\draw[poppurple, thick] (9.5,0.15) -- (9.5,0.7);
\end{tikzpicture}
\end{center}
\legende{Frise chronologique : de l'État fédéral (1961) à la crise anglophone et au statut spécial (2019). Les jalons montrent le passage du fédéralisme à l'État unitaire (1972), puis les tentatives de décentralisation (1996, 2019) face aux revendications des régions du Nord-Ouest et du Sud-Ouest.}
\end{popfigure}

\courssub{Application à la vie courante : lire la mémoire dans l'espace public}

\begin{exemplebox}[Les noms de rues comme marqueurs mémoriels]
À Douala et à Yaoundé, l'espace public porte la trace des mémoires concurrentes : la \cle{rue Um Nyobé} honore le leader de l'UPC, le \cle{boulevard Félix Moumié} rappelle le martyr de Genève, tandis que des lieux portent le nom d'Ahmadou Ahidjo, artisan de l'indépendance de 1960 mais aussi chef de l'État pendant la répression. Ces noms de rues sont des \cle{marqueurs mémoriels} : ils montrent qu'une même ville abrite plusieurs versions de l'histoire nationale. Observer la toponymie (les noms de lieux) de sa commune est donc un exercice d'éducation à la citoyenneté : qui est honoré ? Qui est absent ? Pourquoi ?
\end{exemplebox}

\begin{exoresolu}[Sujet type Bac : « La guerre du maquis au Cameroun : une guerre cachée ? »]
\textbf{Énoncé.} À partir de vos connaissances et des documents étudiés, montrez que la guerre du maquis (1955-1971) peut être qualifiée de « guerre cachée » et expliquez pourquoi sa mémoire reste débattue au Cameroun.
\tcblower
\textbf{Corrigé détaillé (plan et éléments de rédaction).}

\emph{Introduction.} Après l'interdiction de l'UPC en 1955, une guerre oppose jusqu'en 1971 les maquisards upécistes aux forces françaises puis camerounaises. Longtemps absente des manuels et du discours officiel, elle est qualifiée de « guerre cachée ». Problématique : pourquoi cette guerre a-t-elle été occultée, et pourquoi sa mémoire divise-t-elle encore ?

\emph{I. Une guerre réelle et meurtrière.} Contexte : Guerre froide et décolonisation ; la France veut écarter l'UPC, jugée proche du communisme. Faits : maquis en pays Bassa et Bamiléké ; répression (exécutions, villages de regroupement, déplacements massifs) ; bilan estimé entre plusieurs dizaines de milliers et plus de \num{100000} morts selon les sources ; mort des leaders Um Nyobé (1958), Moumié (1960), Ouandié (1971).

\emph{II. Une guerre occultée.} Le pouvoir parle de « pacification » et de « bandits » ; les manuels scolaires passent l'UPC sous silence ; les familles transmettent la mémoire par l'oralité, dans la peur ; les archives françaises restent partiellement classées « Secret Défense ».

\emph{III. Une mémoire disputée mais en voie d'apaisement.} Trois historiographies s'affrontent : version officielle, mémoire UPC (« guerre de libération », « martyrs »), approche universitaire fondée sur les archives déclassifiées et les témoignages. La loi de 1991 réhabilite Um Nyobé, Moumié et Ouandié comme héros nationaux ; les noms de rues (rue Um Nyobé, boulevard Moumié) inscrivent cette mémoire dans l'espace public.

\emph{Conclusion.} La guerre du maquis fut « cachée » parce qu'elle gênait le récit officiel d'une indépendance pacifique. Son étude, fondée sur le croisement des sources, est aujourd'hui une condition de la réconciliation nationale et un exercice de citoyenneté.
\end{exoresolu}

\begin{aretenir}[L'essentiel de la section]
\begin{itemize}
\item La \cle{guerre cachée} (1955-1971) : répression de l'UPC, de plusieurs dizaines de milliers à plus de \num{100000} morts selon les estimations, villages de regroupement, traumatismes transmis aux générations suivantes (pays Bassa et Bamiléké).
\item Trois \cle{historiographies} : version officielle (« pacification »), mémoire UPC (« guerre de libération »), approche universitaire (archives déclassifiées, témoignages oraux).
\item Justice et réparation : réhabilitation des leaders par la loi de 1991, mais archives françaises encore partiellement secrètes.
\item Enjeu actuel : la crise anglophone (depuis 2016) renoue avec les frustrations nées de l'abandon du fédéralisme (1972) ; réponse par la décentralisation et le statut spécial des régions NOSO (2019).
\item Vie courante : les noms de rues et de lieux sont des marqueurs mémoriels que le citoyen doit savoir lire et questionner.
\end{itemize}
\end{aretenir}

\begin{exercice}[Pour aller plus loin]
\begin{enumerate}
\item Définis les termes suivants : historiographie, maquis, village de regroupement, marqueur mémoriel.
\item Explique en dix lignes pourquoi les estimations du nombre de morts de la guerre cachée varient autant d'une source à l'autre.
\item Enquête citoyenne : relève dans ta commune trois noms de rues ou de lieux publics, identifie les personnes ou événements honorés, et rédige un court texte (quinze lignes) justifiant ta démarche et tes sources.
\end{enumerate}
\end{exercice}

\newpage

\coursec{Activité d'intégration}

\courssub{Objectifs de l'activité}

Cette activité d'intégration clôt la leçon sur la marche du Cameroun vers l'indépendance. Elle prend la forme d'un atelier de simulation historique intitulé « Pétition à l'ONU ». À travers la rédaction d'une pétition adressée au Conseil de tutelle des Nations unies en mars 1958, l'élève est amené à mobiliser l'ensemble des savoirs de la leçon et à les appliquer à une situation concrète, comme le demande le programme officiel.

À l'issue de l'atelier, l'élève doit être capable de :
\begin{itemize}
\item expliquer le fonctionnement du régime international de tutelle et le rôle du Conseil de tutelle de l'ONU dans l'évolution politique du Cameroun ;
\item caractériser les grandes tendances du nationalisme camerounais après 1945 (UPC, UC, KNDP) et les stratégies de leurs leaders (Ruben Um Nyobè, Ahmadou Ahidjo, John Ngu Foncha) ;
\item analyser des documents historiques variés (rapports de l'ONU, tracts, discours, cartes électorales, photographies) pour en extraire des arguments ;
\item rédiger et justifier une démarche argumentée dans une forme administrative contrainte (la pétition) ;
\item confronter son point de vue à celui des autres dans un débat contradictoire réglé, compétence utile dans la vie citoyenne courante.
\end{itemize}

\begin{aretenir}[Compétence visée]
L'activité évalue la capacité à \cle{appliquer les notions de l'unité à une situation concrète} et à \cle{rédiger et justifier une réponse argumentée}, deux savoir-faire exigibles au Probatoire et au Baccalauréat technique.
\end{aretenir}

\courssub{Situation}

\textbf{Nous sommes en mars 1958.} Le Conseil de tutelle des Nations unies siège à New York pour examiner le rapport annuel sur les deux Cameroun sous tutelle : le Cameroun oriental, administré par la France (environ \num{3 200 000} habitants), et le Cameroun occidental (Southern Cameroons), administré par le Royaume-Uni comme partie du Nigeria (environ \num{800 000} habitants). Conformément à l'article 87 de la Charte des Nations unies, le Conseil peut recevoir et examiner des \cle{pétitions} émanant des habitants des territoires sous tutelle.

La situation sur le terrain est tendue. Depuis l'interdiction de l'Union des populations du Cameroun (UPC) par l'administration française en juillet 1955, une partie des nationalistes est passée dans la clandestinité et le maquis. L'administration française a engagé une politique de \cle{regroupement des villages} dans les zones de Sanaga-Maritime et du Mungo : des dizaines de milliers de paysans ont été déplacés de force vers des villages contrôlés par l'armée. En février 1958, les élections à l'Assemblée législative du Cameroun (ALCAM) ont porté au pouvoir l'Union camerounaise (UC) d'Ahmadou Ahidjo, favorable à une indépendance négociée avec la France. De son côté, Ruben Um Nyobè, secrétaire général de l'UPC réfugié dans la forêt du Bassa, continue d'exiger l'indépendance immédiate et la réunification des deux Cameroun. Dans le Cameroun britannique, John Ngu Foncha et son parti, le Kamerun National Democratic Party (KNDP), réclament la séparation d'avec le Nigeria et posent la question de la réunification avec le Cameroun français.

Chaque groupe de 4 à 5 élèves incarne l'une des cinq délégations suivantes, invitées à déposer une pétition devant le Conseil :
\begin{enumerate}
\item la délégation de l'\textbf{UPC clandestine} (représentant Ruben Um Nyobè) ;
\item la délégation de l'\textbf{Union camerounaise} (gouvernement d'Ahmadou Ahidjo) ;
\item la délégation du \textbf{KNDP} (British Southern Cameroons, John Ngu Foncha) ;
\item la délégation de l'\textbf{administration française}, puissance administrante du Cameroun oriental ;
\item la délégation de l'\textbf{administration britannique}, puissance administrante du Cameroun occidental.
\end{enumerate}

\textbf{Mission :} chaque délégation rédige une pétition formelle (deux pages maximum) exigeant du Conseil de tutelle une action précise : indépendance immédiate, organisation d'élections libres supervisées par l'ONU, plébiscite dans le Cameroun britannique, maintien de l'ordre et poursuite de la tutelle, ou report de l'échéance.

\textbf{Contraintes obligatoires :}
\begin{itemize}
\item employer le vocabulaire juridique de l'ONU : \emph{Charte des Nations unies, chapitre XII}, \emph{article 76 b} (conduire les populations vers l'auto-gouvernement ou l'indépendance), \emph{Conseil de tutelle}, \emph{puissance administrante}, \emph{résolution} ;
\item citer au moins deux chiffres ou faits précis du dossier documentaire ;
\item proposer une solution concrète et datée ;
\item respecter la forme administrative : en-tête, formule d'appel, exposé des motifs, demande finale, signatures.
\end{itemize}

La restitution se fait oralement (3 minutes par délégation), suivie d'un débat contradictoire dirigé par le professeur, qui incarne le Président du Conseil de tutelle. L'évaluation repose sur une grille critériée : exactitude historique, qualité de l'argumentation, respect de la forme administrative, esprit de synthèse.

\courssub{Consignes}

\textbf{Phase 1 — Analyse du dossier documentaire (30 min).}
\begin{enumerate}
\item À partir de l'extrait du rapport de l'ONU de 1958, relevez deux informations chiffrées qui peuvent servir votre cause et notez-les avec leur source.
\item Sur la carte des résultats électoraux de février 1958, identifiez les régions où l'UC est majoritaire et celles où l'UPC clandestine conserve une influence. Quel argument chaque délégation peut-elle en tirer ?
\item Observez la photographie d'un village regroupé : quels éléments visuels permettent de dénoncer — ou au contraire de justifier — la politique de regroupement ?
\end{enumerate}

\textbf{Phase 2 — Rédaction de la pétition (45 min).}
\begin{enumerate}
\setcounter{enumi}{3}
\item Définissez en une phrase la demande principale de votre délégation (l'« action précise » exigée du Conseil).
\item Construisez trois arguments historiques fondés sur les savoirs de la leçon : le cadre de la tutelle, le bilan de l'administration, les aspirations nationalistes.
\item Rédigez la pétition en respectant la forme administrative et les contraintes de vocabulaire juridique.
\end{enumerate}

\textbf{Phase 3 — Restitution et débat (45 min).}
\begin{enumerate}
\setcounter{enumi}{6}
\item Présentez oralement votre pétition en 3 minutes : un élève lit la demande, un autre défend les arguments.
\item Pendant le débat, répondez aux objections des délégations adverses en citant des faits précis (dates, chiffres, textes).
\item En conclusion individuelle (5 lignes), indiquez quelle délégation vous a paru la plus convaincante et justifiez votre choix.
\end{enumerate}

\courssub{Ressources à mobiliser}

\begin{aretenir}[Rappels de cours utiles]
\begin{itemize}
\item \textbf{La tutelle} : instituée par la résolution de l'Assemblée générale de l'ONU du 13 décembre 1946, elle succède au mandat de la Société des Nations. L'article 76 b de la Charte fixe le but : conduire les populations vers l'auto-gouvernement ou l'indépendance. Le Conseil de tutelle contrôle les rapports annuels des puissances administrantes et reçoit les pétitions.
\item \textbf{Le nationalisme} : né dans les syndicats (USCC de Ruben Um Nyobè, 1944) et les associations, il s'organise en partis après 1945 : UPC (1948), qui exige indépendance immédiate et réunification ; UC d'Ahidjo, favorable à une indépendance négociée ; KNDP de Foncha (1955) dans le Cameroun britannique.
\item \textbf{Les dates clés} : 1946 (tutelle), 1948 (création de l'UPC), 1955 (émeutes de mai, interdiction de l'UPC en juillet), 1956 (loi-cadre Defferre), 1957 (statut interne, ALCAM), 1958 (élections, mort d'Um Nyobè le 13 septembre), 1959 (résolution de l'ONU), 1\textsuperscript{er} janvier 1960 (indépendance du Cameroun français), 1961 (plébiscite et réunification).
\item \textbf{Méthode du document} : identifier la nature, l'auteur, la date, le destinataire ; dégager l'idée principale ; croiser avec ses connaissances.
\end{itemize}
\end{aretenir}

\courssub{Démarche et corrigé commenté}

\begin{exoresolu}[Exemple de pétition attendue : délégation de l'UPC clandestine]
Énoncé. Rédiger, pour la délégation UPC, une pétition exigeant l'indépendance immédiate et la réunification, en respectant les contraintes de l'atelier.

\tcblower

\textbf{Étape 1 — Définition de la demande.} La délégation UPC demande au Conseil de tutelle d'adopter une résolution fixant l'indépendance du Cameroun au plus tard en 1960, l'organisation d'élections générales supervisées par l'ONU et un plébiscite de réunification avec le Cameroun britannique.

\textbf{Étape 2 — Choix des arguments.} Trois arguments fondés sur les savoirs de la leçon :
\begin{itemize}
\item \emph{Argument juridique} : l'article 76 b de la Charte engage la puissance administrante à conduire le territoire vers l'indépendance ; douze ans après 1946, ce but n'est pas atteint.
\item \emph{Argument démocratique} : les élections de février 1958 se sont déroulées alors que l'UPC, principal parti nationaliste créé en 1948, est interdite depuis juillet 1955 ; elles ne reflètent donc pas la volonté du peuple.
\item \emph{Argument humanitaire} : la politique de regroupement des villages a déplacé de force des dizaines de milliers de paysans en Sanaga-Maritime, ce que le rapport de l'ONU de 1958 mentionne ; la tutelle doit protéger les populations, non les opprimer.
\end{itemize}

\textbf{Étape 3 — Rédaction (modèle abrégé).}

\begin{quote}
\textit{Au Président et aux membres du Conseil de tutelle des Nations unies, New York.}

\textit{Pétition de l'Union des populations du Cameroun, représentée par son secrétaire général Ruben Um Nyobè.}

\textit{Exposé des motifs.} Depuis la résolution du 13 décembre 1946, le Cameroun est placé sous tutelle internationale conformément au chapitre XII de la Charte. Or, l'article 76 b de cette Charte assigne à la tutelle le but de conduire les peuples vers l'indépendance. Douze ans après, la France, puissance administrante, maintient notre pays dans la dépendance : elle a interdit notre parti en juillet 1955, organisé en février 1958 des élections sans le principal mouvement nationaliste, et déplacé de force des dizaines de milliers de paysans dans des villages regroupés, comme le constate votre propre rapport de 1958.

\textit{Demande.} Nous demandons au Conseil de tutelle d'adopter une résolution : 1) fixant l'indépendance du Cameroun au plus tard en 1960 ; 2) ordonnant des élections générales libres, supervisées par l'ONU et ouvertes à tous les partis ; 3) organisant un plébiscite sur la réunification des deux Cameroun.

\textit{Fait dans la clandestinité, mars 1958. Pour l'UPC, le secrétaire général.}
\end{quote}

\textbf{Commentaire.} Cette pétition respecte la forme administrative (en-tête, motifs, demande numérotée, signature), cite le vocabulaire juridique exigé (Charte, chapitre XII, article 76 b, résolution, puissance administrante) et s'appuie sur des faits datés et chiffrés du dossier. Les autres délégations suivent la même démarche avec des demandes opposées : l'UC valorise la légalité des élections de 1958 et demande une indépendance négociée ; le KNDP réclame un plébiscite pour trancher entre rattachement au Nigeria et réunification ; les administrations française et britannique plaident pour un report au nom du maintien de l'ordre.
\end{exoresolu}

\begin{attention}[Erreurs fréquentes à éviter]
Ne pas confondre mandat de la SDN (1922) et tutelle de l'ONU (1946) ; ne pas faire siéger l'UPC à l'ALCAM de 1958 (elle est interdite) ; ne pas écrire que la réunification a eu lieu en 1960 (plébiscite de février 1961, réunification le 1\textsuperscript{er} octobre 1961) ; éviter les anachronismes de vocabulaire (pas de « République du Cameroun » avant 1960).
\end{attention}

\courssub{Bilan de l'activité}

Cet atelier a d'abord permis de \cle{concrétiser le fonctionnement de la tutelle} : les élèves ont compris que le Conseil de tutelle n'était pas une abstraction, mais une instance où les peuples colonisés pouvaient faire entendre leur voix — Um Nyobè s'y est d'ailleurs rendu à plusieurs reprises entre 1952 et 1954. La simulation montre ensuite la \cle{pluralité des nationalismes camerounais} : indépendance immédiate pour l'UPC, voie négociée pour l'UC, question spécifique de la réunification pour le KNDP. Le débat contradictoire a mis en lumière que l'indépendance de 1960 et la réunification de 1961 sont le résultat d'un rapport de forces complexe entre nationalistes, puissances administrantes et ONU, et non d'un simple don de la France. Enfin, en rédigeant un texte administratif argumenté et en défendant oralement une position, les élèves ont exercé des compétences directement transposables à la vie citoyenne : pétitionner, argumenter, débattre dans le respect des règles — compétences au cœur du savoir-faire exigé par le programme et évalué au Baccalauréat technique.

\newpage

\coursec{Méthodes et exercices résolus}

\coursec{Chapitre 6 : La marche vers l’indépendance}

\courssub{Méthodes et exercices résolus}

\begin{methode}[Analyser et commenter un document historique (texte, discours, rapport)]
\textbf{Objectif :} Maîtriser l'exercice type du Bac : commentaire de document (texte ou iconographie).\par
\textbf{Démarche en 5 étapes :}
\begin{enumerate}
    \item \textbf{Identification (Nature, Auteur, Date, Contexte, Source) :} Préciser le type de document (discours, traité, article de presse, rapport ONU), son auteur (personnalité, institution), la date exacte, le contexte historique immédiat et la source (archives, manuel, journal).
    \item \textbf{Analyse du contenu (Vocabulaire, Idées forces, Structure) :} Relever les mots-clés (\cle{mandat}, \cle{tutelle}, \cle{indépendance}, \cle{réunification}, \cle{UPC}, \cle{Assemblée législative}). Dégager 2 à 3 idées principales et la thèse défendue.
    \item \textbf{Explication / Mise en perspective (Connaissances) :} Mobiliser le cours pour expliciter les allusions, vérifier la fiabilité (intentionnalité de l'auteur), confronter à d'autres points de vue. C'est l'apport personnel noté.
    \item \textbf{Synthèse (Portée et Limites) :} Résumer l'apport du document à la compréhension du thème (ex : la tutelle, le nationalisme). Signaler ce qu'il ne dit pas (angles morts, parti pris).
    \item \textbf{Conclusion argumentée :} Répondre explicitement à la problématique implicite ou posée, en une phrase forte.
\end{enumerate}
\textbf{Reflexes Bac :} Toujours citer le document (\guillemotleft ...\guillemotright) en intégrant la citation dans sa phrase. Utiliser les connecteurs logiques (\emph{donc, cependant, par ailleurs, en conséquence}). Soigner l'orthographe des noms propres (Um Nyobé, Ahidjo, Foncha, Endeley).
\end{methode}

\begin{exoresolu}[Commentaire d'un extrait du discours de Ruben Um Nyobé à l'ONU (décembre 1952)]
\textbf{Énoncé :} \emph{« Le peuple du Cameroun, par ma voix, demande à l'Organisation des Nations Unies de prendre les mesures nécessaires pour que le Cameroun accède à l'indépendance immédiate et totale, dans l'unité de son territoire et sous un gouvernement démocratique élu au suffrage universel. »} (Extrait de la pétition de l'UPC au Comité de tutelle de l'ONU, New York, décembre 1952).\par
\textbf{Consignes :} Présentez le document, en expliquez le sens et la portée, en le situant dans le contexte de la tutelle de l'ONU et de l'essor du nationalisme camerounais.

\tcblower
\textbf{Solution détaillée :}

\textbf{1. Identification}
\begin{itemize}
    \item \textbf{Nature :} Extrait de pétition / discours officiel prononcé devant une instance internationale.
    \item \textbf{Auteur :} \cle{Ruben Um Nyobé} (1913-1958), secrétaire général de l'\cle{UPC} (Union des Populations du Cameroun), figure majeure du nationalisme radical.
    \item \textbf{Date :} Décembre 1952. Contexte : 4ème voyage d'Um Nyobé à l'ONU (premier en 1952). La France refuse l'indépendance immédiate ; la loi-cadre Defferre n'existe pas encore (1956).
    \item \textbf{Source :} Archives de l'ONU, Comité de tutelle. Document public, propagande politique assumée.
\end{itemize}

\textbf{2. Analyse du contenu}
\emph{Idée force 1 : La revendication d'indépendance immédiate et totale.} Rupture avec les évolutions progressives (loi-cadre 1956). Le terme \guillemotleft immédiate \guillemotright rejette l'autonomie interne.
\emph{Idée force 2 : L'unité territoriale.} Revendication de la réunification du Cameroun français et des deux zones du Cameroun britannique (Nord et Sud) \emph{avant} toute élection. Opposition au risque de balkanisation.
\emph{Idée force 3 : La légitimité démocratique.} \guillemotleft Gouvernement démocratique élu au suffrage universel \guillemotright : critique de l'administration coloniale et des chefferies traditionnelles (indirect rule) jugées non représentatives. Revendication du droit des peuples à disposer d'eux-mêmes (Charte ONU, Art. 76 b).

\textbf{3. Explication / Mise en perspective (Connaissances)}
\begin{itemize}
    \item \textbf{Le cadre de la tutelle :} Depuis 1946, le Cameroun est un \cle{territoire sous tutelle} de l'ONU (ex-mandat français). La France (puissance administrante) a l'obligation de le mener à l'autonomie/indépendance (Art. 76 Charte). L'UPC utilise cette tribune internationale pour contourner l'interdiction française.
    \item \textbf{Stratégie de l'UPC :} \cle{Internationalisation} du conflit. Um Nyobé joue sur la Guerre froide et la décolonisation asiatique (Bandung 1955 à venir) pour faire pression sur la France.
    \item \textbf{Réaction française :} Refus de l'indépendance immédiate. Création de l'Assemblée territoriale (1946) puis législative (1952) au suffrage restreint puis élargi, mais contrôle du scrutin. Interdiction de l'UPC en 1955 (loi-cadre Defferre + émeutes).
    \item \textbf{Contexte britannique :} Au même moment, le Kamerun National Congress (KNC) et le KPP au British Cameroons demandent aussi l'indépendance/réunification, mais la stratégie diffère (intégration au Nigeria vs réunification).
\end{itemize}

\textbf{4. Synthèse (Portée et Limites)}
\textbf{Portée :} Document fondateur du \cle{nationalisme radical camerounais}. Pose les 3 piliers : Indépendance, Unité, Démocratie. Montre l'usage habile de la tribune onusienne.
\textbf{Limites :} Vision partisane (UPC = \guillemotleft le peuple \guillemotright, occultant les autres forces : ANC, KNDP, chefferies, modérés). Ne mentionne pas la guerre froide (accusation de communisme par la France). Silencieux sur la faisabilité immédiate (cadres, économie).

\textbf{5. Conclusion}
Ce document illustre la \cle{stratégie diplomatique de l'UPC} pour transformer la tutelle onusienne en levier d'indépendance immédiate, posant dès 1952 les termes non négociables (unité, suffrage universel) qui structureront le conflit jusqu'à la réunification de 1961, au prix d'une guerre civile occultée par l'histoire officielle.
\end{exoresolu}

\begin{methode}[Rédiger une réponse argumentée structurée (type dissertation / question de synthèse)]
\textbf{Objectif :} Répondre à une question de cours ou de réflexion (souvent au Bac : \guillemotleft Expliquez... \guillemotright, \guillemotleft Montrez que... \guillemotright, \guillemotleft Quels furent les facteurs... \guillemotright).\par
\textbf{Démarche en 4 étapes :}
\begin{enumerate}
    \item \textbf{Analyse du sujet :} Souligner les mots-clés (dates, acteurs, concepts : \cle{nationalisme}, \cle{tutelle}, \cle{réunification}). Définir les bornes chronologiques. Formuler la \cle{problématique} (question centrale à laquelle le plan répondra).
    \item \textbf{Construction du plan (2 ou 3 parties équilibrées) :} Chaque partie = un argument majeur. Chaque sous-partie = un fait précis (date, acteur, événement, chiffre) + explication. Utiliser \textbf{I, II, (III)} puis \textbf{A, B} puis \textbf{1, 2}.
    \item \textbf{Rédaction :} \emph{Introduction} (accroche, définition, problématique, annonce du plan). \emph{Développement} (un paragraphe par sous-partie, phrase d'intro, faits, analyse, phrase de transition). \emph{Conclusion} (réponse synthétique à la problématique, ouverture).
    \item \textbf{Relecture :} Vérifier les dates, les noms propres, la chronologie, la syntaxe, l'orthographe historique.
\end{enumerate}
\textbf{Formule à retenir :} \textbf{1 Paragraphe = 1 Idée force + 1 Exemple précis (Date/Acteur/Chiffre) + 1 Lien avec la problématique.}
\end{methode}

\begin{exoresolu}[Dissertation : Les facteurs de l'essor du nationalisme camerounais (1945-1955)]
\textbf{Énoncé :} \guillemotleft Expliquez les facteurs qui ont favorisé l'essor du nationalisme camerounais entre 1945 et 1955. \guillemotright

\tcblower
\textbf{Solution détaillée :}

\textbf{Introduction}
\textbf{[Accroche]} 1945 marque un tournant mondial : la fin de la Seconde Guerre mondiale fragilise les empires coloniaux. \textbf{[Définition]} Le nationalisme camerounais désigne la prise de conscience collective d'appartenir à une nation camerounaise distincte de la puissance coloniale, et la volonté d'accéder à la souveraineté. \textbf{[Bornes]} 1945 (fin guerre, création ONU) -- 1955 (interdiction UPC, basculement dans la clandestinité/maquis). \textbf{[Problématique]} Comment la conjonction de facteurs externes (contexte international, tutelle ONU) et internes (exploitation coloniale, émergence d'élites, structuration syndicale/partisane) a-t-elle permis l'éclosion d'un nationalisme de masse au Cameroun ? \textbf{[Annonce]} Nous analyserons d'abord le cadre international favorable (I), puis les dynamiques internes syndicales et politiques (II), enfin la radicalisation face au blocage colonial (III).

\textbf{Développement}

\textbf{I. Un contexte international et institutionnel porteur (1945-1950)}
\textbf{A. L'impact de la Seconde Guerre mondiale et la Charte de l'ONU}
\begin{itemize}
    \item Participation des tirailleurs camerounais : découverte de la métropole, remise en cause du mythe de l'invincibilité blanche, attente de récompense (citoyenneté, droits).
    \item Charte de l'ONU (1945), Art. 76 : principe de l'autodétermination des peuples. Création du \cle{système de tutelle} (1946) pour le Cameroun français : obligation pour la France de rendre des comptes à l'ONU (visites périodiques, pétitions). Cela offre une \cle{tribune internationale} aux nationalistes.
\end{itemize}
\textbf{B. La dynamique décolonisation et panafricanisme}
\begin{itemize}
    \item Indépendance de l'Inde (1947), conférence de Bandung (1955) : modèles de libération.
    \item Congrès de Manchester (1945) : panafricanisme (Nkrumah, Kenyatta) influence les élites camerounaises (Um Nyobé, Moumié).
\end{itemize}

\textbf{II. Les ferment internes : de la revendication sociale à la revendication politique (1945-1952)}
\textbf{A. L'école syndicale : l'USCC et la grève de 1945-1946}
\begin{itemize}
    \item Création de l'\cle{USCC} (Union des Syndicats Confédérés du Cameroun, 1945) par Gaston Donnat / Um Nyobé.
    \item Grève générale de 1945-1946 (Douala, Édéa) : revendications salariales \rightarrow conscience de classe \rightarrow conscience nationale. Répression (arrestations) forge des martyrs.
\end{itemize}
\textbf{B. L'émergence des partis politiques : de l'ARCA à l'UPC}
\begin{itemize}
    \item \cle{ARCA} (Assemblée Républicaine du Cameroun, 1946) : modérés (Louis-Paul Aujoulat), évolution dans le cadre français.
    \item \cle{UPC} (Union des Populations du Cameroun, 10 avril 1948) : \cle{Ruben Um Nyobé}, \cle{Félix-Roland Moumié}, \cle{Ernest Ouandié}. Programme radical : indépendance immédiate, réunification, suffrage universel. Recrutement large (syndicalistes, planteurs bassa, bamiléké, intellectuels).
    \item \cle{Partis modérés / ethno-régionaux} : Bloc Démocratique Camerounais (BDC, Ahidjo, Nord), Paysans Indépendants (Sud).
\end{itemize}

\textbf{III. La radicalisation et l'affrontement (1952-1955)}
\textbf{A. L'internationalisation du combat (ONU)}
\begin{itemize}
    \item Pétitions d'Um Nyobé (1952, 1953, 1954). Dénonciation du \guillemotleft pacte colonial \guillemotright. Succès médiatique : l'UPC devient l'interlocuteur privilégié de l'ONU.
\end{itemize}
\textbf{B. Le blocage institutionnel français et la scission}
\begin{itemize}
    \item Élections à l'Assemblée législative (déc. 1952) : suffrage inégal (double collège), fraudes. UPC boycotte ou perd. Ahidjo élu.
    \item \cle{Loi-cadre Defferre (juin 1956)} : autonomie interne, mais maintien des prérogatives françaises (diplomatie, défense, monnaie). Rejetée par l'UPC (indépendance \emph{avant} autonomie).
\end{itemize}
\textbf{C. Basculement dans la violence (mai 1955)}
\begin{itemize}
    \item Manifestations UPC (mai 1955) réprimées (dizaines de morts). Interdiction de l'UPC (13 juillet 1955). Direction exilée (Khoumba, Moumié au Ghana/Guinée/Égypte) \rightarrow \cle{Maquis} (guerre de libération nationale, 1955-1971).
\end{itemize}

\textbf{Conclusion}
\textbf{[Bilan]} L'essor du nationalisme (1945-1955) résulte de la rencontre entre une \cle{opportunité internationale} (tutelle ONU, droit des peuples) et une \cle{maturité interne} (syndicalisme, parti de masse UPC, élites formées). \textbf{[Réponse]} Il se structure autour de deux pôles : un nationalisme \cle{radical et unitaire} (UPC) refusant le partage, et un nationalisme \cle{modéré et fédéraliste} (Ahidjo, Foncha) acceptant l'autonomie progressive. \textbf{[Ouverture]} Cette dualité conditionnera l'échec de la réunification immédiate (plébiscite 1961) et la guerre civile post-indépendance.
\end{exoresolu}

\begin{methode}[Comparer des figures historiques et leurs stratégies politiques (Dossier 3)]
\textbf{Objectif :} Traiter le \guillemotleft Dossier 3 \guillemotright du programme : présenter, comparer et évaluer l'action des grandes figures (Um Nyobé, Moumié, Ouandié, Ahidjo, Foncha, Endeley, Ndeh Ntumazah).\par
\textbf{Démarche en 4 étapes :}
\begin{enumerate}
    \item \textbf{Fiche d'identité pour chaque figure :} Nom, dates, origine ethnique/régionale, formation, parti/organisation, rôle précis (fondateur, stratège, diplomate, chef militaire, homme d'État).
    \item \textbf{Classement par courant idéologique / stratégie :}
        \begin{itemize}
            \item \cle{Radical / Indépendance immédiate / Réunification / Maquis} : Um Nyobé, Moumié, Ouandié (UPC).
            \item \cle{Modéré / Autonomie progressive / Fédéralisme / Négociation} : Ahidjo (Nord, BDC/UC), Foncha (Sud britannique, KNDP), Endeley (Nord britannique, KNC/KPP).
            \item \cle{Intellectuel / Diplomate / Fédéraliste} : Ndeh Ntumazah (UPC puis KNDP/ONU).
        \end{itemize}
    \item \textbf{Tableau comparatif (Critères communs) :} Objectif final, Moyens (légal/illégal, ONU/armes), Base sociale, Rapport à la France/UK, Sort personnel, Héritage mémoriel.
    \item \textbf{Synthèse argumentée :} Expliquer \emph{pourquoi} telle stratégie a prévalu (contexte Guerre froide, rapport de force militaire, soutien extérieur, structure étatique).
\end{enumerate}
\textbf{Reflexe Bac :} Ne pas faire de \guillemotleft portraits \guillemotright juxtaposés. Toujours \textbf{comparer} : \guillemotleft Contrairement à Um Nyobé qui..., Ahidjo choisit de... \guillemotright. Utiliser le vocabulaire précis : \cle{légalisme}, \cle{clandestinité}, \cle{maquis}, \cle{fédéralisme}, \cle{centralisation}.
\end{methode}

\begin{exoresolu}[Portraits croisés : Ruben Um Nyobé vs Ahmadou Ahidjo (Stratégies pour l'indépendance)]
\textbf{Énoncé :} \guillemotleft Comparez les stratégies politiques de Ruben Um Nyobé et d'Ahmadou Ahidjo pour l'accession du Cameroun à l'indépendance (1948-1960). \guillemotright

\tcblower
\textbf{Solution détaillée :}

\textbf{Tableau comparatif synthétique}

\begin{center}
\begin{tabularx}{\linewidth}{|l|X|X|}
\hline
\textbf{Critères} & \textbf{Ruben Um Nyobé (UPC)} & \textbf{Ahmadou Ahidjo (BDC / UC / UNC)} \\
\hline
\textbf{Origine / Base} & Bassa (Sanaga-Maritime), syndicaliste, autodidacte. Base : planteurs Bassa, Bamiléké, syndicalistes, jeunesse urbaine. & Peul/Fulani (Garoua), fils de chef, école des fils de chefs, administration. Base : Nord (chefferies, lamidats), élites administratives, modérés du Sud. \\
\hline
\textbf{Objectif proclamé} & \cle{Indépendance immédiate et totale} + \cle{Réunification} (Cameroun français + britannique) + \cle{Démocratie sociale} (réformes agraires). & \cle{Autonomie interne progressive} (loi-cadre 1956) \rightarrow \cle{Indépendance} (1960) + \cle{Fédéralisme} (réunification 1961) + \cle{Stabilité / Ordre}. \\
\hline
\textbf{Moyens d'action} & \cle{Légal (1948-1955)} : Pétitions ONU, grèves, meetings, presse (\emph{Écho du Cameroun}). \textbf{Illégal (1955-1958)} : Clandestinité, \cle{Maquis} (guérilla), diplomatie itinérante (Pékin, Moscou, Accra, Le Caire). & \cle{Légalisme strict} : Participation aux institutions coloniales (Assemblée territoriale 1946, Assemblée législative 1952, Assemblée de l'Union française 1953). Négociation avec Paris (Loi-cadre, Accords de Yaoundé 1958). Alliances politiciennes (ralliement modérés Sud). \\
\hline
\textbf{Rapport à la France} & \cle{Rupture totale}. Dénonciation du \guillemotleft colonialisme \guillemotright, refus de la Communauté française. Soutien bloc de l'Est / pays non-alignés. & \cle{Négociation / Continuité}. Acceptation de la \cle{Communauté française} (1958). Accords de coopération (défense, monnaie, matières premières). Soutien militaire français contre maquis UPC. \\
\hline
\textbf{Rapport au British Cameroons} & \cle{Réunification immédiate et unitaire}. Soutien au KNDP (Foncha) contre l'intégration au Nigeria. & \cle{Fédéralisme} (2 États, 1 Fédération). Négociation avec Foncha (Conférence de Foumban 1961). Acceptation du plébiscite ONU (1961). \\
\hline
\textbf{Sort / Héritage} & Assassiné par l'armée française (13 sept 1958). \cle{Martyr} du nationalisme radical. \guillemotleft Damnatio memoriae \guillemotright (interdit de citation) jusqu'aux années 1990. Réhabilitation tardive (héros national 1991). & 1er PM (1958), 1er Président (1960-1982). \cle{Constructeur de l'État} (centralisation 1966, parti unique 1966). Héritage controversé : stabilité vs autoritarisme, marginalisation anglophone. \\
\hline
\end{tabularx}
\end{center}

\textbf{Analyse comparative structurée (Rédaction type Bac)}

\textbf{1. Deux conceptions antagonistes de la souveraineté}
Um Nyobé incarne la \cle{rupture révolutionnaire} : la souveraineté se conquiert \emph{par le bas} (peuple en armes, suffrage universel) et \emph{par le droit international} (ONU). L'indépendance est \emph{indivisible} (politique, économique, culturelle). Ahidjo incarne la \cle{transition négociée} : la souveraineté se construit \emph{par le haut} (élites, administration, armée) et \emph{par le droit constitutionnel français} (loi-cadre, Communauté). L'indépendance est \emph{partageable} (coopération, francophonie).

\textbf{2. Le facteur décisif : la violence d'État et la Guerre froide}
La stratégie d'Um Nyobé échoue militairement (maquis isolé, manque d'armes lourdes, répression française massive \emph{Opération Mouette}, quadrillage) et diplomatiquement (veto français à l'ONU, contexte Guerre froide : UPC perçue comme pro-communiste). La stratégie d'Ahidjo réussit car elle \cle{s'aligne sur la Realpolitik française} : la France préfère transférer le pouvoir à un allié fiable (Ahidjo) qui garantit ses intérêts (accords de coopération) plutôt qu'à un mouvement révolutionnaire imprévisible (UPC).

\textbf{3. La réunification : Unité unitaire vs Unité fédérale}
Um Nyobé veut un \cle{État unitaire} fort (condition de la survie nationale). Ahidjo impose le \cle{fédéralisme} (Foumban 1961) pour rassurer les Anglais du Sud (Foncha) et conserver le contrôle du centre (Yaoundé). \emph{Paradoxe :} Ahidjo réalisera l'unité (1961) que réclamait Um Nyobé, mais sous une forme (fédérale puis unitaire 1972) et avec des alliés (modérés) que ce dernier combattait.

\textbf{Conclusion}
La comparaison révèle que l'\cle{indépendance du Cameroun (1960) et la réunification (1961) sont le fruit de la stratégie légaliste d'Ahidjo}, soutenue par la France, mais qu'elles portent en germe les contradictions (marginalisation des vainqueurs de la décolonisation radicale, contentieux anglophone) issues de l'élimination physique et symbolique de la stratégie d'Um Nyobé.
\end{exoresolu}

\begin{methode}[Reconstituer et expliquer un processus institutionnel chronologique (Échéancier / Frise) ]
\textbf{Objectif :} Maîtriser la chronologie complexe 1956-1961 (Loi-cadre \rightarrow Indépendance \rightarrow Plébiscite \rightarrow Réunification) et expliquer la logique causale entre les étapes.\par
\textbf{Démarche en 3 étapes :}
\begin{enumerate}
    \item \textbf{Établir la frise chronologique annotée (Repères obligatoires Bac) :} Placer les 10-12 dates clés avec l'événement \textbf{ET} sa signification (cause/conséquence).
    \item \textbf{Identifier les tournants (Ruptures) :} 1956 (Loi-cadre / Interdiction UPC), 1958 (Ahidjo PM / Accords Yaoundé), 1960 (Indépendance Est), 1961 (Plébiscite / Réunification).
    \item \textbf{Expliquer la logique d'acteurs :} Pour chaque étape, nommer l'acteur principal (France, ONU, Ahidjo, Foncha, Um Nyobé/Moumié) et son intérêt.
\end{enumerate}
\textbf{Repères chronologiques incontournables (à connaître par cœur) :}
\begin{itemize}
    \item \textbf{23 juin 1956 :} Loi-cadre Defferre (autonomie interne) + Interdiction UPC (13 juil).
    \item \textbf{1957 :} 1er gouvernement africain (André-Marie Mbida PM, puis Ahidjo VP).
    \item \textbf{Fév 1958 :} Ahmadou Ahidjo devient Premier Ministre (chute Mbida).
    \item \textbf{14 oct 1958 :} \cle{Accords de Yaoundé} (France / Ahidjo) : Indépendance fixée au 1er janv 1960, coopération étroite.
    \item \textbf{13 mars 1959 :} \cle{Référendum constitutionnel} (OUI massif) \rightarrow Constitution du 16 avril 1959.
    \item \textbf{1er janv 1960 :} \cle{Indépendance du Cameroun oriental} (République du Cameroun). Ahidjo Président.
    \item \textbf{11 fév 1961 :} \cle{Plébiscite ONU} au British Cameroons. Nord \rightarrow Nigeria (60\%), Sud \rightarrow Réunification (70\%).
    \item \textbf{1er oct 1961 :} \cle{Réunification} / Indépendance du Sud. Naissance de la \cle{République Fédérale du Cameroun} (Ahidjo PR, Foncha VP).
    \item \textbf{Août 1961 :} Conférence de \cle{Foumban} (Constitution fédérale).
\end{itemize}
\end{methode}

\begin{exoresolu}[Ordonner et expliquer : Les étapes institutionnelles de 1956 à 1961]
\textbf{Énoncé :} \guillemotleft Classez les événements suivants dans l'ordre chronologique et expliquez en deux lignes la portée de chacun pour l'accession à la souveraineté : Accords de Yaoundé ; Loi-cadre Defferre ; Plébiscite du 11 février 1961 ; Indépendance du 1er janvier 1960 ; Conférence de Foumban ; Référendum constitutionnel de mars 1959. \guillemotright

\tcblower
\textbf{Solution détaillée :}

\textbf{Ordre chronologique et Portée}

\begin{enumerate}
    \item \textbf{23 juin 1956 -- Loi-cadre Defferre (Loi n°56-619).}
    \emph{Portée :} Instauré l'\cle{autonomie interne} (gouvernement responsable devant l'Assemblée, suffrage universel). Tournant majeur : fin de l'administration directe, début du pouvoir africain. \textbf{Mais} : la France garde la diplomatie, la défense, la monnaie, la justice. Provoque la scission : UPC refuse (indépendance d'abord) \rightarrow interdiction (juil. 1955) \rightarrow maquis.
    \item \textbf{14 octobre 1958 -- Accords de Yaoundé (France / Ahidjo).}
    \emph{Portée :} \cle{Négociation bilatérale} actant l'indépendance pour le 1er janv 1960. Ahidjo obtient la souveraineté politique \emph{contre} la signature d'accords de coopération (défense, monnaie CFA, matières premières, conseillers techniques). Scelle l'alliance Ahidjo-France et marginalise l'UPC (exclue des négociations).
    \item \textbf{13 mars 1959 -- Référendum constitutionnel.}
    \emph{Portée :} Approbation par le peuple (suffrage universel) du projet de Constitution préparant l'indépendance. Légitimité démocratique interne pour le régime Ahidjo. Score officiel : 99,9\% OUI (vote encadré, UPC interdite/clandestine).
    \item \textbf{1er janvier 1960 -- Indépendance du Cameroun oriental (République du Cameroun).}
    \emph{Portée :} \cle{Souvenir international} (adhésion ONU sept 1960). Ahidjo 1er Président. Fin de la tutelle française. \textbf{Limite :} Ne concerne que l'ex-Cameroun français. La question de la réunification avec le British Cameroons reste posée (résolution ONU 1352 XIV déc 1959 organisant le plébiscite).
    \item \textbf{11 février 1961 -- Plébiscite ONU au British Cameroons.}
    \emph{Portée :} \cle{Choix démocratique} imposé par l'ONU (Rés. 1352). Deux questions distinctes :
    \begin{itemize}
        \item Nord : Intégration au Nigeria (60\% OUI) \rightarrow Perte pour la réunification.
        \item Sud : Réunification avec la République du Cameroun (70\% OUI) \rightarrow Victoire de l'option Foncha/KNDP (soutenue par Ahidjo/ONU) contre Endeley/KNC (intégration Nigeria).
    \end{itemize}
    \item \textbf{Juillet-Août 1961 -- Conférence de Foumban (Rédaction Constitution fédérale).}
    \emph{Portée :} Négociation des \cle{modalités institutionnelles} de la réunification (1er oct 1961). Compromis : \cle{Fédéralisme} (2 États fédérés : East/West, gouvernement fédéral fort). Ahidjo impose un fédéralisme \guillemotleft centralisé \guillemotright (contrôle budget, armée, parti unique naissant). Foncha obtient des garanties (bicamérisme, VP anglophone).
\end{enumerate}

\textbf{Synthèse causale :} La \cle{Loi-cadre (1956)} ouvre la voie institutionnelle ; les \cle{Accords de Yaoundé (1958)} scellent le deal France-Ahidjo ; le \cle{Référendum (1959)} légitime le pouvoir ; l'\cle{Indépendance (1960)} acte la souveraineté de l'Est ; le \cle{Plébiscite (1961)} règle le sort de l'Ouest (perte Nord / gain Sud) ; \cle{Foumban (1961)} institutionnalise la réunification fédérale.
\end{exoresolu}

\begin{methode}[Exploiter un croquis ou une carte schématique pour illustrer une dynamique territoriale]
\textbf{Objectif :} Réaliser ou commenter une carte/croquis schématique (ex: partition 1919, tutelle 1946, plébiscite 1961, réunification 1961). Compétence transversale (Histoire/Géo).\par
\textbf{Démarche en 4 étapes :}
\begin{enumerate}
    \item \textbf{Lecture de la consigne / du sujet :} Identifier le thème (ex: \guillemotleft Le Cameroun sous tutelle \guillemotright, \guillemotleft Le plébiscite de 1961 \guillemotright) et l'échelle (pays, région).
    \item \textbf{Choix des informations pertinentes (Légende) :} Sélectionner 4 à 6 items max (Zones, Villes, Flèches, Symboles). Hiérarchiser : Fond (zones colorées) > Forme (villes, frontières) > Flux (flèches).
    \item \textbf{Réalisation graphique (Rigueur) :}
        \begin{itemize}
            \item Contour simplifié (polygone reconnaissable).
            \item \textbf{Code couleur strict :} Zone française (bleu), Zone britannique Nord (rouge), Zone britannique Sud (vert), Frontière internationale (trait fin), Frontière inter-zones (trait pointillé épais).
            \item Villes : Points noirs + noms (Yaoundé, Douala, Garoua, Buea, Bamenda).
            \item Flèches : Résultats plébiscite (vers Nigeria / vers Yaoundé), Maquis (zones de guérilla).
        \end{itemize}
    \item \textbf{Légende organisée et Titre :} Titre informatif (Thème + Date). Légende structurée (Fonds, Signes, Flux). \textbf{Échelle approximative} et \textbf{Flèche Nord}.
\end{enumerate}
\textbf{Attention :} Pas de coloriage artistique. Traits à la règle. Écriture lisible (minuscules pour légende, majuscules pour titre). Crayons de couleur (pas feutres qui traversent).
\end{methode}

\begin{exoresolu}[Croquis commenté : Le plébiscite du 11 février 1961 et la réunification]
\textbf{Énoncé :} \guillemotleft Réalisez un croquis schématique du Cameroun au moment du plébiscite de février 1961. La légende doit faire apparaître les deux zones de tutelle britannique, le résultat du vote et la capitale fédérale future. \guillemotright

\tcblower
\textbf{Solution détaillée : Description du croquis attendu (Code TikZ indicatif pour l'enseignant / Description textuelle pour l'élève)}

\textbf{1. Structure du croquis (Description textuelle pour réalisation manuelle)}

\begin{popfigure}
\begin{tikzpicture}[scale=0.8]
% Contour très simplifié du Cameroun (polygone approximatif)
\draw[popdark, thick] (0,0) -- (6,0) -- (7,2) -- (6.5,5) -- (4,6) -- (1,5.5) -- (0,3) -- cycle;
% Frontière Est/Ouest (approximative)
\draw[popdark, dashed, thick] (3.5,0) -- (3.5,1.5) -- (4,2.5) -- (3.8,4) -- (3.5,5.5);
% Frontière Nord/Sud zone britannique
\draw[popdark, dotted, thick] (4.5,2.5) -- (6.5,2.5);

% Zones colorées (Fonds)
\fill[popblueL!50] (0,0) -- (6,0) -- (7,2) -- (6.5,5) -- (4,6) -- (1,5.5) -- (0,3) -- cycle; % Base tout le pays
% Zone Française (Est) - Bleu
\fill[popblue!30] (0,0) -- (3.5,0) -- (3.5,1.5) -- (4,2.5) -- (3.8,4) -- (3.5,5.5) -- (1,5.5) -- (0,3) -- cycle;
% Zone Britannique Nord - Rouge
\fill[poporange!30] (3.5,0) -- (6,0) -- (7,2) -- (6.5,2.5) -- (4.5,2.5) -- (3.5,1.5) -- cycle;
% Zone Britannique Sud - Vert
\fill[popgreen!30] (4.5,2.5) -- (6.5,2.5) -- (6.5,5) -- (4,6) -- (3.8,4) -- (4,2.5) -- cycle;

% Villes
\node[circle,fill=popdark,inner sep=1.5pt] (Yde) at (2.5,3.5) {};
\node[right] at (Yde) {\small Yaoundé (Capitale Fédérale)};
\node[circle,fill=popdark,inner sep=1.5pt] (Dla) at (1.5,1.5) {};
\node[below] at (Dla) {\small Douala};
\node[circle,fill=popdark,inner sep=1.5pt] (Gar) at (4.5,0.8) {};
\node[below] at (Gar) {\small Garoua};
\node[circle,fill=popdark,inner sep=1.5pt] (Buea) at (5.5,3.5) {};
\node[right] at (Buea) {\small Buea (Cap. Ouest)};
\node[circle,fill=popdark,inner sep=1.5pt] (Bam) at (5.5,4.5) {};
\node[right] at (Bam) {\small Bamenda};

% Flèches Résultats Plébiscite
\draw[->, poporange, thick] (5.2,1.2) -- (7.5,1.2) node[right] {\small Nord : Intégration Nigeria (60\%)};
\draw[->, popgreen, thick] (5.5,4) -- (7.5,4) node[right] {\small Sud : Réunification (70\%)};

% Flèche Nord
\node[draw, circle, fill=popgold, inner sep=2pt] at (7,5.5) {N};
\draw[->, popdark] (7,5.3) -- (7,4.8);

% Légende (simulée ici par des nodes)
\begin{scope}[xshift=8.5cm, yshift=-1cm]
\node[anchor=west] at (0,2.5) {\textbf{LÉGENDE}};
\fill[popblue!30] (0,2) rectangle (0.5,2.3); \node[anchor=west] at (0.6,2.15) {Cameroun français (Indép. 1960)};
\fill[poporange!30] (0,1.5) rectangle (0.5,1.8); \node[anchor=west] at (0.6,1.65) {British Northern Cameroons};
\fill[popgreen!30] (0,1) rectangle (0.5,1.3); \node[anchor=west] at (0.6,1.15) {British Southern Cameroons};
\draw[popdark, dashed, thick] (0,0.5) -- (0.5,0.5); \node[anchor=west] at (0.6,0.5) {Frontière inter-zones (1916/1946)};
\draw[popdark, dotted, thick] (0,0) -- (0.5,0); \node[anchor=west] at (0.6,0) {Frontière Nord/Sud (Brit.)};
\draw[->, poporange, thick] (0,-0.5) -- (0.5,-0.5); \node[anchor=west] at (0.6,-0.5) {Vote : Intégration Nigeria};
\draw[->, popgreen, thick] (0,-1) -- (0.5,-1); \node[anchor=west] at (0.6,-1) {Vote : Réunification};
\node[circle,fill=popdark,inner sep=1.5pt] at (0.25,-1.6) {}; \node[anchor=west] at (0.6,-1.6) {Ville / Capitale};
\end{scope}
\end{tikzpicture}
\legende{Croquis schématique : Le Cameroun au plébiscite de février 1961. Partition ONU (1946), résultats du vote (Nord/Nigeria, Sud/Réunification), capitales.}
\end{popfigure}

\textbf{2. Commentaire type (Rédaction attendue)}
\guillemotleft Ce croquis illustre la \cle{partition du Kamerun} issue du traité de Versailles (1919) confirmée par la tutelle ONU (1946). La \cle{zone française (Est, bleu)} accède à l'indépendance le 1er janv 1960. La \cle{zone britannique} est divisée en deux unités administratives distinctes : le \cle{Nord (orange)} et le \cle{Sud (vert)}. Le plébiscite du 11 fév 1961, organisé par l'ONU (Rés. 1352), scelle le sort de ces deux territoires : le Nord choisit l'intégration à la Fédération du Nigeria (flèche orange), le Sud choisit la réunification avec la République du Cameroun (flèche verte). La capitale fédérale, \cle{Yaoundé}, située en zone Est, devient le centre de décision de la nouvelle République Fédérale du Cameroun (1er oct 1961). Ce croquis met en évidence l'\cle{échec partiel de l'idéal réunificationniste} (perte du Nord) et la construction d'un État fédéral asymétrique. \guillemotright
\end{exoresolu}

\begin{methode}[Mobiliser l'histoire pour une réflexion citoyenne (Éducation à la citoyenneté / Héritages)]
\textbf{Objectif :} Répondre aux questions de l'axe \guillemotleft Héritages, mémoires et enjeux historiographiques \guillemotright (Complément Bac). Lier passé et présent sans anachronisme.\par
\textbf{Démarche en 3 étapes :}
\begin{enumerate}
    \item \textbf{Identifier le fait historique fondateur :} Ex: Guerre de l'UPC (1955-1971), Réunification (1961), Fédéralisme (1961-1972), Parti unique (1966).
    \item \textbf{Analyser la construction mémorielle :} Qui commémore ? Qui oublie ? Quels sont les \cle{lieux de mémoire} (stèles, noms de rues, programmes scolaires, jours fériés) ? Ex: \cle{Journée de la Réunification (1er mai)} vs \cle{Silence sur la guerre du maquis} (jusqu'aux années 90). \cle{Réhabilitation d'Um Nyobé (1991)}.
    \item \textbf{Formuler l'enjeu citoyen actuel :} En quoi ce passé éclaire le présent ? (Ex: Crise anglophone 2016-... \leftarrow non-respect engagements Foumban / centralisation 1972 ; Gouvernance \leftarrow héritage parti unique / clientélisme ; Identité nationale \leftarrow pluralisme mémoriel).
\end{enumerate}
\textbf{Formule de rédaction :} \guillemotleft L'héritage de [Événement] se manifeste aujourd'hui par [Réalité actuelle], car [Mécanisme historique : continuité institutionnelle, trauma non résolu, construction identitaire]. La citoyenneté exige [Reconnaissance / Réparation / Dialogue / Réforme]. \guillemotright
\end{methode}

\begin{exoresolu}[Héritages et citoyenneté : La réunification de 1961 face à la crise anglophone actuelle]
\textbf{Énoncé :} \guillemotleft En vous appuyant sur les accords de Foumban (1961) et la loi de 1972, expliquez en quoi l'héritage institutionnel de la réunification explique une partie des revendications de la crise anglophone (depuis 2016). \guillemotright

\tcblower
\textbf{Solution détaillée :}

\textbf{1. L'engagement de Foumban (Août 1961) : Le contrat fédéral}
\begin{itemize}
    \item \textbf{Contexte :} Négociation entre Ahidjo (Rép. du Cameroun, 5,5M hab) et Foncha (Rép. Fédérale du Sud, 1M hab). Rapport de force défavorable au Sud.
    \item \textbf{Contenu clé (Constitution fédérale 1er sept 1961) :}
        \begin{itemize}
            \item \cle{Fédéralisme} : 2 États fédérés (East Cameroon / West Cameroon) avec assemblées, gouvernements, judicature (Common Law), services publics propres.
            \item \cle{Bicamérisme} : Assemblée nationale (députés) + Sénat (représentation égale des États ? Non, non appliqué initialement).
            \item \cle{Présidence tournante ?} Non. Ahidjo Président Fédéral, Foncha Vice-Président.
            \item \cle{Partis politiques} : Maintien des partis d'État (UNC à l'Est, KNDP à l'Ouest) \rightarrow Fusion forcée 1966 (CNU).
        \end{itemize}
\end{itemize}

\textbf{2. La dérive unitaire (1962-1972) : Rupture du contrat}
\begin{itemize}
    \item \textbf{1962 :} Code de la nationalité (uniformisation).
    \item \textbf{1966 :} Création du \cle{CNU} (Parti unique). Dissolution KNDP. Perte de l'autonomie politique de l'Ouest.
    \item \textbf{1968 / 1972 :} Remplacement de Foncha (VP) par Muna (Sud) puis Bello (Nord). Affaiblissement de la vice-présidence.
    \item \textbf{20 mai 1972 :} \cle{Référendum} (contesté, OUI massif officiel) \rightarrow \cle{République Unie du Cameroun}. \textbf{Abrogation unilatérale du fédéralisme.} Suppression des institutions de l'Ouest (Assemblée, PM, Judicature Common Law absorbée). Centralisation totale à Yaoundé.
    \item \textbf{1984 :} Décret Biya \rightarrow \cle{République du Cameroun} (nom de l'ex-Est). Symbole effacement de l'identité fédérale/ouest.
\end{itemize}

\textbf{3. Lien avec la crise anglophone (2016-...) : Mécanismes historiques}
\begin{itemize}
    \item \textbf{Grief 1 : Non-respect du \guillemotleft contrat de Foumban \guillemotright.} Les anglophones perçoivent 1972 comme une \cle{annexion} / violation du traité international (accord entre deux États souverains). Revendication : \cle{Retour au fédéralisme} (2 ou 10 États) ou \cle{Sécession} (Ambazonie).
    \item \textbf{Grief 2 : Effacement du \cle{Common Law} et de l'\cle{éducation} anglo-saxonne.} Harmonisation forcée (droit civil français, système francophone). Perte de l'identité professionnelle (avocats, enseignants) \rightarrow Grèves 2016.
    \item \textbf{Grief 3 : \cle{Marginalisation politique et économique}.} Centralisation des recettes (pétrole Sud-Ouest / Nord-Ouest) \rightarrow redistribution opaques à Yaoundé. Sous-représentation aux postes clés (ministères régaliens, armée).
    \item \textbf{Grief 4 : \cle{Mémoire conflictuelle}.} Héroïsation d'Ahidjo/Biya (bâtisseurs) vs Oubli/Stigmatisation des leaders ouest (Foncha, Muna, Endeley) et surtout des \cle{martyrs UPC} (Um Nyobé) qui portaient aussi l'unité. La \cle{réhabilitation d'Um Nyobé (1991)} a rouvert la boîte de Pandore mémorielle.
\end{itemize}

\textbf{Conclusion citoyenne}
La crise anglophone n'est pas un accident conjoncturel mais la \cle{résurgence d'un contentieux constitutionnel non résolu} (1961-1972). La citoyenneté camerounaise actuelle exige une \cle{relecture critique de l'histoire} (sortir du roman national unique), la \cle{reconnaissance du traumatisme} de la guerre du maquis et de la perte du fédéralisme, et l'ouverture d'un \cle{dialogue inclusif} sur la forme de l'État (décentralisation effective vs fédéralisme) pour transformer l'héritage conflictuel en contrat social renouvelé.
\end{exoresolu}

\begin{methode}[Calculer et interpréter des données statistiques historiques (Élections, Référendum, Démographie)]
\textbf{Objectif :} Exploiter des tableaux chiffrés (résultats électoraux, référendums, population) pour appuyer une démonstration historique. Compétence quantitative.\par
\textbf{Démarche en 4 étapes :}
\begin{enumerate}
    \item \textbf{Lecture critique du tableau :} Source (officielle ? contestée ?), Date, Population concernée, Mode de scrutin (suffrage universel ? restreint ? boycott ?).
    \item \textbf{Calculs de base (obligatoires) :}
        \begin{itemize}
            \item \cle{Taux de participation} = (Votants / Inscrits) $\times$ 100.
            \item \cle{Pourcentage des voix} = (Voix liste / Votants exprimés) $\times$ 100.
            \item \cle{Taux de variation} (évolution entre deux dates) = $\frac{Valeur_{final} - Valeur_{initial}}{Valeur_{initial}} \times 100$.
        \end{itemize}
    \item \textbf{Interprétation historique :} Ne pas laisser le chiffre nu. Expliquer \emph{pourquoi} ce chiffre (contexte, fraude, boycott, propagande, enjeu). Comparer aux attentes.
    \item \textbf{Intégration dans l'argumentation :} \guillemotleft Le score de 99,9\% au référendum de 1959 traduit l'absence de pluralisme (UPC interdite) plus qu'un plébiscite populaire spontané. \guillemotright
\end{enumerate}
\textbf{Attention :} Pourcentages avec \textbf{virgule} (français). Arrondir au dixième (ex: 70,5\%). Unités : \%, millions d'habitants.
\end{methode}

\begin{exoresolu}[Analyse statistique : Les résultats du plébiscite de février 1961]
\textbf{Énoncé :} \guillemotleft À partir des données ci-dessous, calculez les pourcentages de votes \guillemotleft Pour la réunification \guillemotright (Sud) et \guillemotleft Pour l'intégration au Nigeria \guillemotright (Nord). Interprétez les résultats au regard de la géographie humaine et des enjeux politiques. \guillemotright
\begin{center}
\begin{tabular}{|l|c|c|c|}\hline
\textbf{Zone} & \textbf{Inscrits} & \textbf{Votants} & \textbf{Voix \guillemotleft Réunification \guillemotright / \guillemotleft Intégration \guillemotright} \\ \hline
Nord (Intégration Nigeria) & 175 000 & 158 000 & 97 000 (Intégration) / 61 000 (Réunification) \\ \hline
Sud (Réunification) & 210 000 & 195 000 & 138 000 (Réunification) / 57 000 (Intégration) \\ \hline
\end{tabular}
\end{center}

\tcblower
\textbf{Solution détaillée :}

\textbf{1. Calculs}

\textbf{Zone Nord (British Northern Cameroons) :}
\begin{itemize}
    \item Taux de participation : $\frac{158\,000}{175\,000} \times 100 = \mathbf{90,3\%}$ (Forte mobilisation).
    \item \% Intégration (Nigeria) : $\frac{97\,000}{158\,000} \times 100 = \mathbf{61,4\%}$.
    \item \% Réunification (Cameroun) : $\frac{61\,000}{158\,000} \times 100 = \mathbf{38,6\%}$.
    \item \textbf{Résultat :} Majorité pour l'\cle{intégration au Nigeria}.
\end{itemize}

\textbf{Zone Sud (British Southern Cameroons) :}
\begin{itemize}
    \item Taux de participation : $\frac{195\,000}{210\,000} \times 100 = \mathbf{92,9\%}$ (Très forte mobilisation).
    \item \% Réunification (Cameroun) : $\frac{138\,000}{195\,000} \times 100 = \mathbf{70,8\%}$.
    \item \% Intégration (Nigeria) : $\frac{57\,000}{195\,000} \times 100 = \mathbf{29,2\%}$.
    \item \textbf{Résultat :} Majorité nette pour la \cle{réunification avec la République du Cameroun}.
\end{itemize}

\textbf{2. Interprétation historique}

\textbf{A. Une fracture géographique, ethnique et religieuse structurante}
\begin{itemize}
    \item \textbf{Nord (Peuls / Musulmans / Éleveurs) :} Liens séculaires avec le Nord-Nigeria (Sokoto, Bornou), commerce transsaharien, langue haoussa/fulfulde. L'intégration au Nigeria (État du Nord, musulman) assure la continuité culturelle et économique. Le KNC (Endeley) puis NPPP mobilisent sur ce registre.
    \item \textbf{Sud (Bamiléké, Bamoun, Anglophones côtiers / Chrétiens / Planteurs) :} Histoire partagée avec l'Ouest camerounais français (révoltes 1950, UPC, maquis). Liens économiques vers Douala/Yaoundé (chemin de fer, port). Le KNDP (Foncha) joue la carte de la \cle{réunification des frères} et de l'autonomie fédérale promise par Ahidjo.
\end{itemize}

\textbf{B. Le rôle des élites et de la campagne}
\begin{itemize}
    \item \textbf{Nord :} \cle{Ahmadou Ahidjo} (originaire du Nord, PM du Cameroun Est) intervient personnellement... \emph{pour l'intégration au Nigeria} ? \textbf{Non.} Paradoxe : Ahidjo, pour assurer sa majorité fédérale future, a \cle{laissé faire} (voire soutenu discrètement) l'option Nigeria au Nord (peuplé, mais hostile à l'UPC/maquis), tout en faisant campagne activement pour la réunification au Sud (Foncha allié). Calcul politicien : Mieux vaut un Sud fédéral loyal (1M hab) qu'un Nord fédéral hostile (0,5M hab) qui aurait renforcé l'opposition UPC/Nord.
    \item \textbf{Sud :} \cle{John Ngu Foncha} (KNDP) mène une campagne efficace : \guillemotleft Votez pour vos frères \guillemotright, promesse de fédéralisme fort (Common Law, éducation).
\end{itemize}

\textbf{C. Conséquences immédiates et à long terme}
\begin{itemize}
    \item \textbf{1er juin 1961 :} Indépendance du Nord \rightarrow Intégration Nigeria (État de Sardauna).
    \item \textbf{1er oct 1961 :} Indépendance du Sud \rightarrow Réunification (État fédéré de l'Ouest).
    \item \textbf{Héritage :} Le Cameroun indépendant \cle{perd le Nord} (ressources, cohérence territoriale) mais \cle{gagne le Sud} (côte, plantations, main-d'œuvre qualifiée, port de Victoria/Limbé). La frontière Nord/Sud (ligne de partage des eaux Bénoué/Cross River) devient frontière internationale Cameroun/Nigeria (source de contentieux Bakassi plus tard).
\end{itemize}

\textbf{Conclusion chiffrée :} Le plébiscite valide le \cle{partage du British Cameroons} (61\% Nord pour Nigeria, 71\% Sud pour Cameroun). Il reflète des identités profondes (islam vs christianité, sahélien vs forestier) instrumentalisées par les élites (Ahidjo/Foncha) pour construire un État fédéral viable \emph{mais amputé} de sa partie septentrionale.
\end{exoresolu}

\begin{attention}[Les 5 erreurs qui coûtent des points au Bac (Histoire Terminale Tech)]
\begin{enumerate}
    \item \textbf{Confondre \cle{Mandat} (SDN, 1919-1946) et \cle{Tutelle} (ONU, 1946-1960).} Le mandat impliquait une \guillemotleft mission civilisatrice \guillemotright sans contrôle réel ; la tutelle impose l'\cle{obligation d'indépendance} (Art. 76) et un \cle{contrôle onusien} (visites, pétitions). \emph{Piège :} Dire \guillemotleft La France a géré le Cameroun comme une colonie sous mandat après 1945 \guillemotright $\rightarrow$ \textbf{Faux}.
    \item \textbf{Réduire le nationalisme à l'UPC seul.} Occulter les \cle{nationalismes modérés / fédéralistes} (Ahidjo, Foncha, Endeley, Aujoulat) qui ont \emph{réussi} la transition institutionnelle. \emph{Piège :} Traiter \guillemotleft Le nationalisme camerounais \guillemotright comme un bloc monolithique radical. \textbf{Nuance :} Il y a \cle{deux nationalismes} concurrents (Radical vs Modéré) aux stratégies opposées.
    \item \textbf{Inverser la chronologie 1958-1961.} Placer l'indépendance (1er janv 1960) \emph{après} le plébiscite (fév 1961) ou la réunification (1er oct 1961). \emph{Ordre impératif :} Loi-cadre (56) $\rightarrow$ Accords Yaoundé (58) $\rightarrow$ Indépendance Est (1 janv 60) $\rightarrow$ Plébiscite (11 fév 61) $\rightarrow$ Réunification (1er oct 61) $\rightarrow$ Foumban (août 61).
    \item \textbf{Oublier le sort du \cle{Nord} au plébiscite de 1961.} Écrire \guillemotleft Le British Cameroons a choisi la réunification \guillemotright. \textbf{Faux.} Le Nord a choisi le Nigeria (60\%). Seule la minorité (Sud, 70\%) a choisi la réunification. C'est un \cle{partage}, pas une réunification totale.
    \item \textbf{Traiter la \cle{guerre du maquis (1955-1971)} comme un simple \guillemotleft trouble à l'ordre public \guillemotright ou l'ignorer.} C'est un \cle{conflit armé majeur} (dizaines de milliers de morts, déplacements, torture, villages regroupés) qui structure la mémoire et la politique (réhabilitation 1991, crise anglophone). Le Bac attend une \cle{analyse distanciée} : causes (interdiction UPC, répression), acteurs (maquisards, armée française, armée camerounaise post-60), conséquences (hégémonie Ahidjo, silence mémoriel).
\end{enumerate}
\textbf{Conseil final :} \cle{Chronologie + Acteurs + Concepts (Tutelle, Fédéralisme, Maquis, Réunification) + Nuance} = Succès au Probatoire/Bac.
\end{attention}

\newpage

\coursec{Exercices}

\courssub{Je vérifie mes connaissances}

\begin{exercice}[QCM : révisions rapides]
Pour chaque question, une seule réponse est exacte. Entoure la bonne lettre.
\begin{enumerate}
\item La différence essentielle entre le \emph{mandat} de la SDN (1922) et la \emph{tutelle} de l'ONU (1946) est :
\begin{enumerate}
\item la tutelle vise explicitement à conduire les peuples vers l'autonomie et l'indépendance ;
\item la tutelle supprime toute présence française au Cameroun ;
\item le mandat prévoyait déjà l'indépendance du Cameroun ;
\item il n'y a aucune différence entre les deux statuts.
\end{enumerate}
\item L'Union des Populations du Cameroun (UPC) est créée en :
\begin{enumerate}
\item 1945 \quad b) 1948 \quad c) 1955 \quad d) 1958
\end{enumerate}
\item Le premier Premier ministre du Cameroun, nommé en 1957, est :
\begin{enumerate}
\item Ruben Um Nyobé \quad b) Ahmadou Ahidjo \quad c) André-Marie Mbida \quad d) John Ngu Foncha
\end{enumerate}
\item Lors des plébiscites de février 1961, le Cameroun septentrional (Nord) choisit :
\begin{enumerate}
\item le rattachement à la République du Cameroun ;
\item le rattachement au Nigeria ;
\item l'indépendance séparée ;
\item le maintien sous tutelle britannique.
\end{enumerate}
\item La Conférence de Foumban, qui adopte la Constitution de la République fédérale du Cameroun, se tient en :
\begin{enumerate}
\item juillet 1961 \quad b) octobre 1961 \quad c) janvier 1960 \quad d) février 1961
\end{enumerate}
\item Les deux étoiles du drapeau camerounais adopté en 1961 symbolisent :
\begin{enumerate}
\item les deux religions principales du pays ;
\item le Cameroun oriental et le Cameroun occidental réunifiés ;
\item les deux langues officielles ;
\item le Cameroun et le Nigeria.
\end{enumerate}
\end{enumerate}
\end{exercice}

\begin{exercice}[Vrai ou faux justifié]
Indique si chaque affirmation est vraie ou fausse, puis justifie ta réponse en une ou deux phrases.
\begin{enumerate}
\item Le Cameroun devient un territoire sous tutelle de l'ONU le 13 décembre 1946.
\item L'UPC est interdite par l'administration française en juillet 1955.
\item Ruben Um Nyobé plaide la cause de l'indépendance du Cameroun à la tribune de l'ONU à New York.
\item Le Cameroun méridional (Sud) vote son rattachement au Nigeria lors du plébiscite de 1961.
\item La réunification entre le Cameroun oriental et le Cameroun occidental est effective le 1\ier{} octobre 1961.
\end{enumerate}
\end{exercice}

\begin{exercice}[Définitions]
Définis en deux ou trois phrases chacune des notions suivantes, en donnant pour chacune un exemple tiré du cas camerounais :
\begin{enumerate}
\item la \cle{tutelle} internationale ;
\item le \cle{nationalisme} ;
\item le \cle{plébiscite} ;
\item la \cle{réunification}.
\end{enumerate}
\end{exercice}

\begin{exercice}[Dates et acteurs]
Complète le tableau suivant en associant chaque date à l'événement et à l'acteur correspondants.
\begin{center}
\begin{tabularx}{\linewidth}{|l|X|X|}
\hline
\rowcolor{popblueL} \textbf{Date} & \textbf{Événement} & \textbf{Acteur(s)} \\
\hline
1948 & \ldots & \ldots \\
\hline
13 décembre 1952 & \ldots & \ldots \\
\hline
13 février 1958 & \ldots & \ldots \\
\hline
1\ier{} janvier 1960 & \ldots & \ldots \\
\hline
11-12 février 1961 & \ldots & \ldots \\
\hline
\end{tabularx}
\end{center}
\end{exercice}

\courssub{Je m'entraîne}

\begin{exercice}[Le cadre de la tutelle]
À partir de tes connaissances sur l'accord de tutelle de 1946 :
\begin{enumerate}
\item Rappelle les deux puissances administrantes du Cameroun et les parties du territoire qu'elles administrent.
\item Explique en cinq à huit lignes en quoi l'article 76 de la Charte de l'ONU crée un cadre « contraignant » pour la puissance administrante.
\item Cite deux obligations de la France en tant que puissance administrante.
\end{enumerate}
\end{exercice}

\begin{exercice}[Des syndicats aux partis]
\begin{enumerate}
\item Présente en quelques lignes le rôle des syndicats (notamment l'USCC) dans l'éveil politique camerounais après 1945.
\item Explique pourquoi l'UPC est considérée comme le premier grand parti nationaliste camerounais : programme, implantation, méthodes.
\item Compare brièvement les stratégies de l'UPC et de l'UC (Union camerounaise) d'Ahmadou Ahidjo.
\end{enumerate}
\end{exercice}

\begin{exercice}[Chronologie raisonnée]
Replace sur une frise chronologique les huit événements suivants, puis choisis-en trois et explique pour chacun en quoi il constitue une étape décisive vers l'indépendance : création de l'UPC ; discours d'Um Nyobé à l'ONU ; interdiction de l'UPC ; loi-cadre de 1956 ; assassinat d'Um Nyobé ; indépendance du Cameroun oriental ; plébiscites de 1961 ; Conférence de Foumban.
\end{exercice}

\begin{exercice}[Portraits croisés]
Choisis deux grandes figures du nationalisme camerounais parmi Ruben Um Nyobé, Ahmadou Ahidjo, André-Marie Mbida et John Ngu Foncha. Pour chacune :
\begin{enumerate}
\item présente son parcours et son parti en trois lignes ;
\item décris sa stratégie pour accéder à l'indépendance ;
\item dégage en conclusion un point de convergence et un point de divergence entre les deux figures.
\end{enumerate}
\end{exercice}

\courssub{Je me prépare au Bac}

\begin{exercice}[Dissertation — sujet type Baccalauréat]
\textbf{Sujet :} « En vous appuyant sur l'étude du Dossier 3 et sur vos connaissances, montrez que l'accession du Cameroun à la souveraineté internationale (1960-1961) résulte de la convergence entre une dynamique internationale (ONU) et des stratégies d'acteurs nationaux aux objectifs divergents. »

\textbf{Consignes :}
\begin{enumerate}
\item Rédige une introduction comportant une accroche, une présentation du sujet, une problématique et l'annonce du plan.
\item Développe ton analyse en deux ou trois parties équilibrées (par exemple : le rôle de l'ONU et du cadre de la tutelle ; les stratégies divergentes des nationalistes camerounais ; la convergence aboutissant à l'indépendance et à la réunification).
\item Rédige une conclusion qui répond à la problématique et propose une ouverture.
\end{enumerate}
\textbf{Repères à mobiliser :} accord de tutelle (1946), article 76 de la Charte, UPC (1948), tribune de l'ONU (1952-1953), loi-cadre (1956), indépendance (1\ier{} janvier 1960), plébiscites (février 1961), réunification (1\ier{} octobre 1961).
\end{exercice}

\begin{exercice}[Analyse de document — type Probatoire]
\textbf{Document :} extrait du discours de Ruben Um Nyobé, secrétaire général de l'UPC, devant la Quatrième Commission de l'Assemblée générale de l'ONU, New York, décembre 1952 : « Nous demandons l'unification des deux Cameroun et l'indépendance du pays après un délai de dix ans. [\ldots] L'article 76 de la Charte des Nations unies impose à la puissance administrante l'obligation de développer les institutions politiques des territoires sous tutelle et de les conduire progressivement vers l'autonomie. »

\textbf{Questions :}
\begin{enumerate}
\item Identifie la nature, la date, l'auteur et le contexte de ce document. (4 points)
\item Explique la référence faite à l'article 76 de la Charte de l'ONU : quelles obligations impose-t-il à la puissance administrante ? (4 points)
\item Quelles sont les deux revendications principales exprimées par Um Nyobé ? (2 points)
\item Montre comment ce document illustre la stratégie « dualiste » de l'UPC : lutte intérieure au Cameroun et recours à la tribune internationale de l'ONU. (6 points)
\item Ce discours a-t-il atteint ses objectifs ? Justifie ta réponse à l'aide de la suite des événements (1955-1961). (4 points)
\end{enumerate}
\end{exercice}

\begin{exercice}[Croquis historique noté — travaux pratiques]
\textbf{Sujet :} « Le Cameroun sous tutelle (1946-1961) : espaces de conflictualité et de négociation. »

Réalise un croquis historique sur une demi-feuille, en respectant les consignes suivantes :
\begin{enumerate}
\item \textbf{Fond de carte :} trace les frontières des territoires sous tutelle française et britannique ; place les capitales (Yaoundé, Buea) et les principales zones d'implantation de l'UPC (pays bassa, pays bamiléké, Sanaga-Maritime) et de l'UC (Nord-Cameroun).
\item \textbf{Éléments figurés :} hachure les zones de maquis upéciste ; marque par des ronds les villes étapes de la négociation internationale (New York, Foumban) ; dessine des flèches figurant la réunification de 1961 (Cameroun méridional vers la République du Cameroun, Cameroun septentrional vers le Nigeria).
\item \textbf{Légende organisée} en quatre rubriques : Contexte, Acteurs, Événements, Résultats.
\item \textbf{Titre} en haut du croquis, \textbf{orientation} (flèche du nord) et échelle indicative.
\end{enumerate}
\textbf{Barème indicatif :} fond de carte exact (5 points), figurés correctement placés (8 points), légende organisée (5 points), propreté et présentation (2 points).
\end{exercice}

\newpage

\coursec{Corrigés des exercices}

\coursec{Leçon 08 : La marche vers l’indépendance (1945-1961)}
\courssub{Chapitre 6 : La marche vers l’indépendance}

\begin{corrige}[Exercice 1 — QCM : révisions rapides]
\begin{enumerate}
\item \textbf{Réponse a.} La \cle{tutelle de l'ONU} (article 76 de la Charte) fixe explicitement comme objectif le développement progressif des institutions politiques et la marche des territoires vers l'autonomie ou l'indépendance, alors que le mandat de la SDN ne prévoyait pas d'échéance vers l'indépendance.
\item \textbf{Réponse b.} L'\cle{UPC} est créée en 1948, à Douala, à partir notamment de syndicats comme l'\cle{USCC}.
\item \textbf{Réponse c.} \cle{André-Marie Mbida} devient le premier Premier ministre du Cameroun en mai 1957, à la tête du premier gouvernement autonome.
\item \textbf{Réponse b.} Le \cle{Cameroun septentrional} (Nord) choisit le rattachement au Nigeria lors du plébiscite du 11 février 1961 (effectif le 1\ier{} juin 1961).
\item \textbf{Réponse a.} La \cle{Conférence de Foumban} se tient du 17 au 21 juillet 1961 et adopte la Constitution de la République fédérale du Cameroun.
\item \textbf{Réponse b.} Les deux étoiles symbolisent les deux entités réunifiées : le Cameroun oriental (francophone) et le Cameroun occidental (anglophone, Sud).
\end{enumerate}
\textbf{Points de vigilance :} ne pas confondre les dates des plébiscites (février 1961) avec celle de la réunification (1\ier{} octobre 1961) ; ne pas confondre Mbida (premier Premier ministre, 1957) et Ahidjo (Premier ministre à partir de février 1958, puis président).
\end{corrige}

\begin{corrige}[Exercice 2 — Vrai ou faux justifié]
\begin{enumerate}
\item \textbf{Vrai.} L'Assemblée générale de l'ONU approuve les accords de tutelle le 13 décembre 1946 : le Cameroun passe du statut de mandat de la SDN à celui de territoire sous tutelle, administré par la France et le Royaume-Uni.
\item \textbf{Vrai.} L'UPC est dissoute et interdite par décret de l'administration française le 13 juillet 1955, après les émeutes de mai 1955 ; ses dirigeants passent dans la clandestinité ou en exil.
\item \textbf{Vrai.} \cle{Ruben Um Nyobé} se rend à plusieurs reprises à New York (notamment en décembre 1952 et 1953) pour plaider devant la Quatrième Commission de l'ONU l'unification des deux Cameroun et l'indépendance.
\item \textbf{Faux.} C'est le Cameroun septentrional (Nord) qui choisit le rattachement au Nigeria ; le Cameroun méridional (Sud) vote pour le rattachement à la République du Cameroun.
\item \textbf{Vrai.} Le 1\ier{} octobre 1961, le Cameroun méridional rejoint la République du Cameroun : naît la République fédérale du Cameroun, concrétisant la \cle{réunification}.
\end{enumerate}
\textbf{Points de vigilance :} une justification est exigée pour chaque item, même quand l'affirmation est vraie ; pour les affirmations fausses, corriger l'erreur ne suffit pas sans donner le fait exact.
\end{corrige}

\begin{corrige}[Exercice 3 — Définitions]
\begin{enumerate}
\item \textbf{La tutelle internationale} : régime institué par la Charte de l'ONU (1945) plaçant un territoire non autonome sous l'administration d'une puissance dite « administrante », chargée de le conduire vers l'autonomie ou l'indépendance, sous le contrôle de l'ONU. \emph{Exemple :} le Cameroun, placé sous tutelle française et britannique par les accords du 13 décembre 1946.
\item \textbf{Le nationalisme} : courant de pensée et mouvement politique affirmant la volonté d'un peuple de constituer une nation souveraine et de s'affranchir de la domination étrangère. \emph{Exemple :} le nationalisme camerounais incarné par l'UPC, qui revendique l'indépendance et l'unification des deux Cameroun.
\item \textbf{Le plébiscite} : consultation par laquelle une population se prononce directement par vote sur une question précise, ici son destin national. \emph{Exemple :} les plébiscites des 11 et 12 février 1961 organisés sous l'égide de l'ONU dans les Cameroun britanniques.
\item \textbf{La réunification} : réunion en un seul État de territoires d'une même entité historique précédemment séparés. \emph{Exemple :} le rattachement du Cameroun méridional à la République du Cameroun le 1\ier{} octobre 1961, formant la République fédérale du Cameroun.
\end{enumerate}
\textbf{Points de vigilance :} une définition complète comporte l'idée générale, les éléments essentiels et un exemple ; l'exemple camerounais est ici explicitement demandé, son absence coûte des points.
\end{corrige}

\begin{corrige}[Exercice 4 — Dates et acteurs]
\begin{center}
\begin{tabularx}{\linewidth}{|l|X|X|}
\hline
\rowcolor{popblueL} \textbf{Date} & \textbf{Événement} & \textbf{Acteur(s)} \\
\hline
1948 & Création de l'UPC à Douala & Ruben Um Nyobé, Félix-Roland Moumié, Ernest Ouandié \\
\hline
13 décembre 1952 & Discours à la tribune de l'ONU à New York : revendication de l'unification et de l'indépendance & Ruben Um Nyobé \\
\hline
13 février 1958 & Démission d'André-Marie Mbida ; Ahmadou Ahidjo devient Premier ministre & André-Marie Mbida, Ahmadou Ahidjo \\
\hline
1\ier{} janvier 1960 & Proclamation de l'indépendance du Cameroun oriental (République du Cameroun) & Ahmadou Ahidjo \\
\hline
11-12 février 1961 & Plébiscites sous l'égide de l'ONU dans les Cameroun britanniques & Populations des Cameroun septentrional et méridional, John Ngu Foncha \\
\hline
\end{tabularx}
\end{center}
\textbf{Points de vigilance :} pour le 13 février 1958, l'événement clé est le passage de relais Mbida-Ahidjo ; pour 1961, bien préciser qu'il s'agit des plébiscites et non de la réunification elle-même (1\ier{} octobre 1961).
\end{corrige}

\begin{corrige}[Exercice 5 — Le cadre de la tutelle]
\begin{enumerate}
\item Les deux puissances administrantes sont la \textbf{France}, qui administre la plus grande partie du territoire (le Cameroun oriental, capitale Yaoundé), et le \textbf{Royaume-Uni}, qui administre deux bandes frontalières occidentales (les Cameroun britanniques : Cameroun septentrional et méridional, capitale Buea), rattachées administrativement au Nigeria.
\item L'article 76 de la Charte crée un cadre contraignant car il fixe des \emph{objectifs obligatoires} à la puissance administrante : favoriser le progrès politique, économique et social des populations, développer leurs institutions politiques et les conduire progressivement vers l'\textbf{autonomie} ou l'\textbf{indépendance}. La puissance administrante n'est donc plus souveraine sur le territoire : elle exerce une administration sous \emph{contrôle international}, avec obligation de rendre compte chaque année à l'Assemblée générale de l'ONU, qui peut entendre des pétitions et envoyer des missions d'enquête. Le Cameroun n'est plus une possession coloniale ordinaire : son statut international engage juridiquement la France et le Royaume-Uni devant la communauté internationale.
\item Deux obligations de la France : présenter des \textbf{rapports annuels} à l'ONU sur l'administration du territoire ; développer les \textbf{institutions politiques} locales (assemblées représentatives, participation des Camerounais) en vue de l'autonomie. On peut aussi citer : recevoir les pétitions des populations et les missions de visite de l'ONU.
\end{enumerate}
\textbf{Points de vigilance :} la question 2 exige un développement rédigé (cinq à huit lignes), pas une simple liste ; citer explicitement l'article 76 et le contrôle de l'ONU.
\end{corrige}

\begin{corrige}[Exercice 6 — Des syndicaux aux partis]
\begin{enumerate}
\item Après 1945, les syndicats jouent un rôle d'\emph{école politique} : créé en 1944, l'\cle{USCC} (Union des syndicats confédérés du Cameroun), dirigée par Ruben Um Nyobé, défend les travailleurs (salaires, conditions de travail) mais dépasse vite le cadre revendicatif pour dénoncer la domination coloniale. Les syndicats forment des cadres, diffusent une culture de l'organisation et de la lutte, et servent de tremplin à la création des partis politiques nationalistes.
\item L'\cle{UPC} (1948) est le premier grand parti nationaliste par son \textbf{programme} : indépendance immédiate et unification des deux Cameroun ; par son \textbf{implantation} : enracinement populaire massif (pays bassa, pays bamiléké, Sanaga-Maritime), dizaines de milliers de militants ; par ses \textbf{méthodes} : pétitions à l'ONU, tribune internationale (Um Nyobé à New York), presse, meetings, puis boycott et lutte armée après son interdiction en juillet 1955.
\item L'\textbf{UPC} adopte une stratégie de rupture : indépendance immédiate, contestation frontale de la présence française, recours à l'ONU puis lutte armée clandestine après 1955. L'\textbf{UC} d'Ahmadou Ahidjo (1958) choisit une stratégie légaliste et graduelle : participation aux institutions issues de la loi-cadre, négociation avec la France, indépendance obtenue par étapes dans le cadre de la coopération. Les deux partis visent la souveraineté, mais divergent sur le rythme, les moyens et les relations avec l'ancienne puissance administrante.
\end{enumerate}
\textbf{Points de vigilance :} ne pas réduire l'UPC à la seule lutte armée : son action à l'ONU (1948-1955) est antérieure et essentielle ; la comparaison doit porter sur des critères identiques (objectif, méthodes, rapport à la France).
\end{corrige}

\begin{corrige}[Exercice 7 — Chronologie raisonnée]
Ordre chronologique des huit événements :
\begin{enumerate}
\item Création de l'UPC — 1948 ;
\item Discours d'Um Nyobé à l'ONU — décembre 1952 ;
\item Interdiction de l'UPC — juillet 1955 ;
\item Loi-cadre — 23 juin 1956 ;
\item Assassinat d'Um Nyobé — 13 septembre 1958 ;
\item Indépendance du Cameroun oriental — 1\ier{} janvier 1960 ;
\item Plébiscites — 11-12 février 1961 ;
\item Conférence de Foumban — juillet 1961.
\end{enumerate}
Trois étapes décisives (exemples attendus) :
\begin{itemize}
\item \textbf{Création de l'UPC (1948)} : pour la première fois, un parti de masse formule clairement les revendications d'indépendance et d'unification ; il fait entrer la question camerounaise à l'ONU et impose le nationalisme comme force politique incontournable.
\item \textbf{Loi-cadre de 1956} : elle accorde l'autonomie interne aux territoires sous tutelle française ; les élections de 1956-1957 et la création d'un gouvernement camerounais (Mbida puis Ahidjo) constituent le passage institutionnel décisif vers la souveraineté.
\item \textbf{Plébiscites de février 1961} : sous contrôle de l'ONU, ils règlent démocratiquement le sort des Cameroun britanniques et rendent possible la réunification du 1\ier{} octobre 1961.
\end{itemize}
\textbf{Points de vigilance :} la frise doit respecter l'ordre et l'échelle du temps ; chaque explication doit montrer le \emph{lien causal} avec l'indépendance, pas seulement décrire l'événement.
\end{corrige}

\begin{corrige}[Exercice 8 — Portraits croisés]
Exemple de correction avec Ruben Um Nyobé et Ahmadou Ahidjo (tout autre couple correctement traité est recevable).
\begin{enumerate}
\item \textbf{Ruben Um Nyobé} (1913-1958) : syndicaliste bassa, secrétaire général de l'USCC puis de l'UPC ; il incarne l'aile radicale du nationalisme. \textbf{Ahmadou Ahidjo} (1924-1989) : originaire du Nord, ancien fonctionnaire des Postes, député puis fondateur de l'UC (1958), Premier ministre en 1958 puis premier président de la République.
\item \textbf{Stratégie d'Um Nyobé} : indépendance immédiate et unification des deux Cameroun, revendiquées par les pétitions et la tribune de l'ONU (1952-1953), puis, après l'interdiction de l'UPC (1955), par la lutte clandestine et le maquis ; il est tué en 1958. \textbf{Stratégie d'Ahidjo} : voie légaliste et négociée — participation aux institutions de la loi-cadre, gouvernement autonome, demande d'indépendance acceptée par la France et l'ONU, aboutissant au 1\ier{} janvier 1960, puis négociation de la réunification avec Foncha (Foumban, juillet 1961).
\item \textbf{Convergence} : les deux hommes œuvrent à la souveraineté nationale du Cameroun et affirment l'existence d'une nation camerounaise. \textbf{Divergence} : Um Nyobé privilégie la rupture avec la France et l'indépendance immédiate (voie de la contestation, y compris armée), tandis qu'Ahidjo privilégie la négociation graduelle et la coopération avec l'ancienne puissance administrante (voie légaliste).
\end{enumerate}
\textbf{Points de vigilance :} respecter le format demandé (parcours en trois lignes, stratégie, puis convergence/divergence en conclusion) ; les dates et les partis doivent être exacts (UPC 1948, UC 1958).
\end{corrige}

\begin{corrige}[Exercice 9 — Dissertation : éléments de corrigé et barème]
\textbf{Introduction attendue.} \emph{Accroche} : au lendemain de la Deuxième Guerre mondiale, le mouvement de décolonisation s'accélère sous l'égide de l'ONU. \emph{Présentation du sujet} : le Cameroun, placé sous tutelle franco-britannique depuis 1946, accède à l'indépendance le 1\ier{} janvier 1960 (partie orientale) puis se réunifie partiellement le 1\ier{} octobre 1961. \emph{Problématique} : comment la combinaison d'un cadre international contraignant et de stratégies nationalistes divergentes a-t-elle conduit le Cameroun à la souveraineté ? \emph{Annonce du plan} : I. Le cadre de la tutelle : une dynamique internationale vers l'indépendance ; II. Des stratégies nationales divergentes ; III. La convergence aboutissant à l'indépendance et à la réunification.

\textbf{Développement attendu.}
\begin{itemize}
\item \textbf{I.} L'accord de tutelle de 1946 et l'article 76 de la Charte obligent les puissances administrantes à conduire le territoire vers l'autonomie ; rapports annuels, pétitions, missions de l'ONU ; la tribune de New York offre aux nationalistes une caisse de résonance (Um Nyobé, 1952-1953) ; la loi-cadre de 1956 accorde l'autonomie interne ; les résolutions de l'ONU fixent l'indépendance (1959) et organisent les plébiscites de 1961.
\item \textbf{II.} Divergences des acteurs : l'UPC (Um Nyobé, Moumié, Ouandié) revendique l'indépendance immédiate et l'unification, par la voie internationale puis armée après son interdiction (1955) ; Mbida (1957-1958) refuse d'abord de demander l'indépendance immédiate ; Ahidjo (UC) choisit la négociation graduelle avec la France ; Foncha (KNDP) fait campagne pour le rattachement du Cameroun méridional à la République du Cameroun.
\item \textbf{III.} Convergence : la pression de l'ONU et celle du nationalisme se renforcent mutuellement ; Ahidjo obtient l'indépendance du Cameroun oriental (1\ier{} janvier 1960) ; les plébiscites de février 1961 et la Conférence de Foumban (juillet 1961) aboutissent à la République fédérale (1\ier{} octobre 1961) — réalisant, par d'autres voies, une partie du programme upéciste.
\end{itemize}
\textbf{Conclusion attendue} : réponse claire à la problématique (l'indépendance résulte bien de la convergence ONU/acteurs nationaux) ; ouverture possible : les limites de la réunification (le Nord rejoint le Nigeria) ou la mémoire conflictuelle de l'UPC.

\textbf{Barème indicatif (sur 20)} : introduction complète (4 pts) ; cohérence du plan et de la démonstration (4 pts) ; exactitude et mobilisation des connaissances (dates, acteurs, repères du sujet) (8 pts) ; conclusion avec ouverture (2 pts) ; qualité de l'expression (2 pts).

\textbf{Points de vigilance :} la problématique doit être une question explicite ; chaque partie doit mobiliser au moins deux repères datés ; éviter la simple chronologie narrative sans analyse de la « convergence » ; ne pas écrire que l'UPC a obtenu l'indépendance : c'est le gouvernement Ahidjo qui la proclame, l'UPC y ayant contribué par sa pression.
\end{corrige}

\begin{corrige}[Exercice 10 — Analyse de document : corrigé détaillé]
\begin{enumerate}
\item \textbf{Présentation (4 pts)} : il s'agit d'un \emph{discours politique} (extrait d'une allocution officielle), prononcé en \emph{décembre 1952} à New York par \emph{Ruben Um Nyobé}, secrétaire général de l'UPC, devant la Quatrième Commission de l'Assemblée générale de l'ONU (chargée des questions de tutelle). \emph{Contexte} : le Cameroun est territoire sous tutelle franco-britannique depuis 1946 ; l'UPC, créée en 1948, multiplie les pétitions et obtient d'être entendue par l'ONU.
\item \textbf{Article 76 (4 pts)} : cet article de la Charte impose à la puissance administrante de favoriser le progrès politique, économique et social des populations, de développer leurs institutions politiques et de les conduire progressivement vers l'autonomie ou l'indépendance, compte tenu de leurs aspirations. Um Nyobé s'appuie donc sur le droit international pour rappeler la France à ses obligations.
\item \textbf{Revendications (2 pts)} : l'\textbf{unification} des deux Cameroun (français et britannique) et l'\textbf{indépendance} du pays après un délai de dix ans.
\item \textbf{Stratégie « dualiste » (6 pts)} : à l'intérieur du Cameroun, l'UPC mène une lutte de masse — syndicats, meetings, presse, pétitions collectives — pour mobiliser la population ; à l'extérieur, elle utilise la tribune de l'ONU pour internationaliser la question camerounaise, faire pression sur la France et obtenir l'appui de l'opinion mondiale. Le document illustre ce double registre : Um Nyobé parle au nom du peuple camerounais (ancrage intérieur) tout en invoquant l'article 76 devant une instance internationale (levier extérieur).
\item \textbf{Bilan (4 pts)} : partiellement. À court terme, non : la France refuse l'indépendance immédiate, l'UPC est interdite en 1955 et Um Nyobé est tué en 1958. Mais à moyen terme, ses objectifs sont en partie réalisés : le Cameroun oriental accède à l'indépendance le 1\ier{} janvier 1960 et la réunification avec le Cameroun méridional intervient le 1\ier{} octobre 1961 — soit dans le délai de dix ans annoncé — même si le Cameroun septentrional rejoint le Nigeria et si l'indépendance est obtenue par d'autres acteurs (Ahidjo).
\end{enumerate}
\textbf{Points de vigilance :} en présentation de document, les quatre éléments (nature, date, auteur, contexte) sont chacun notés ; la question 4 exige de \emph{croiser} le document avec les connaissances ; la question 5 demande un jugement nuancé et justifié par des faits datés, pas une opinion.
\end{corrige}

\begin{corrige}[Exercice 11 — Croquis historique : corrigé-type et barème]
\begin{popfigure}
\begin{tikzpicture}[scale=0.75, every node/.style={font=\small}]
% --- Fond de carte : contour simplifié du Cameroun ---
\draw[popdark, thick] (0,0) -- (8,0) -- (7.5,6) -- (5.5,8) -- (2,8) -- (0.5,5) -- cycle;

% --- Limite Est/Ouest (Tutelle française / britannique) ---
\draw[popblue, dashed, thick] (2.5,0) -- (2.5,2.5) -- (3,3.5) -- (3,5.5) -- (3.5,7);

% --- Limite Nord/Sud (Cameroun britanniques) ---
\draw[poporange, dashed, thick] (0.5,2.5) -- (2.5,2.5);

% --- Zones UPC (hachures) ---
\pattern[pattern=north east lines, pattern color=popgreen!60] (1,0.5) rectangle (3,2.5); % Sanaga-Maritime / Bassa
\pattern[pattern=north west lines, pattern color=popgreen!60] (1.5,3) rectangle (3.5,5); % Bamiléké

% --- Zone UC (Nord) ---
\fill[poporange!20] (3,5.5) -- (7.5,6) -- (5.5,8) -- (3,7) -- cycle;
\node[poporange!80!black, align=center] at (4.5,6.5) {Zone d'implantation\\ de l'\cle{UC}\\(Nord)};

% --- Capitales ---
\node[circle, fill=popblue, inner sep=2pt, label=above:\textbf{Yaoundé}] (Ya) at (3.5,3) {};
\node[circle, fill=popblue, inner sep=2pt, label=left:\textbf{Buea}] (Bu) at (1.2,1.8) {};
\node[circle, fill=popdark, inner sep=1.5pt, label=right:\textbf{Douala}] at (1.5,0.8) {};
\node[circle, fill=popdark, inner sep=1.5pt, label=above:\textbf{Bafoussam}] at (2.2,3.8) {};
\node[circle, fill=popdark, inner sep=1.5pt, label=right:\textbf{Garoua}] at (5,5.5) {};
\node[circle, fill=popdark, inner sep=1.5pt, label=right:\textbf{Maroua}] at (5.5,6.5) {};

% --- Zones de maquis (symboles) ---
\foreach \x/\y in {1.5/1.5, 2/2, 2.5/3.5, 3/4} {
    \draw[popgreen!70!black, decorate, decoration={snake, amplitude=1pt, segment length=4pt}] (\x,\y) -- ++(0.3,0.3) node[right, font=\tiny] {Maquis};
}

% --- Lieux de négociation ---
\node[star, star points=5, star point ratio=2.25, fill=popgold, draw=popdark, inner sep=2pt, label=above:\textbf{Foumban}] at (2.5,4.2) {};
% New York (encart)
\begin{scope}[shift={(9,6)}, scale=0.5]
\draw[popdark] (0,0) rectangle (2,1);
\node at (1,0.5) {\textbf{New York (ONU)}};
\draw[->, popblue, thick] (-1,0.5) -- (0,0.5);
\end{scope}

% --- Flèches 1961 ---
% Sud -> Cameroun (Réunification)
\draw[->, popblue, very thick, >=Stealth] (1.2,1.5) -- (3,2.5) node[midway, above, sloped, font=\tiny] {Rattachement (1\ier{} oct. 1961)};
% Nord -> Nigeria
\draw[->, poporange, very thick, >=Stealth] (1.5,4.5) -- (0,5.5) node[midway, above, sloped, font=\tiny] {Rattachement au Nigeria (1\ier{} juin 1961)};

% --- Légende intégrée (simplifiée pour la figure) ---
\begin{scope}[shift={(-4,-0.5)}]
\matrix[draw=popline, fill=popcream, inner sep=3pt, row sep=2pt, column sep=5pt] {
\node[pattern=north east lines, pattern color=popgreen!60, minimum size=6pt] {}; & \node[anchor=west] {Implantation UPC (Bassa, Sanaga-Maritime, Bamiléké)}; \\
\node[pattern=north west lines, pattern color=popgreen!60, minimum size=6pt] {}; & \node[anchor=west] {Zones de maquis UPC (1955-1960)}; \\
\fill[poporange!20] (0,0) rectangle (6pt,6pt); & \node[anchor=west] {Implantation UC (Nord-Cameroun)}; \\
\draw[popblue, dashed, thick] (0,3pt) -- (6pt,3pt); & \node[anchor=west] {Limite tutelle française / britannique}; \\
\draw[poporange, dashed, thick] (0,3pt) -- (6pt,3pt); & \node[anchor=west] {Limite Cameroun septentrional / méridional}; \\
\node[circle, fill=popblue, inner sep=2pt] {}; & \node[anchor=west] {Capitales administratives (Yaoundé, Buea)}; \\
\node[star, star points=5, fill=popgold, draw=popdark, inner sep=2pt] {}; & \node[anchor=west] {Lieu de négociation (Foumban 1961)}; \\
\draw[->, popblue, very thick] (0,3pt) -- (6pt,3pt); & \node[anchor=west] {Flux réunification (Sud $\to$ Cameroun)}; \\
\draw[->, poporange, very thick] (0,3pt) -- (6pt,3pt); & \node[anchor=west] {Flux rattachement Nord $\to$ Nigeria}; \\
};
\end{scope}

% --- Nord, Échelle, Titre ---
\node[anchor=south west, font=\Large\bfseries] at (-4,8.5) {Le Cameroun sous tutelle (1946-1961) : espaces de conflictualité et de négociation};
\node[anchor=south east] at (8,-0.5) {\textbf{N}};
\draw[->, thick] (7.5,-0.5) -- (7.5,0.5);
\node[anchor=north west] at (-4,-0.8) {Échelle indicative : $\approx$ 200 km};
\end{tikzpicture}
\legende{Croquis historique : le Cameroun sous tutelle (1946-1961). La légende organisée en rubriques (Contexte, Acteurs, Événements, Résultats) figure dans la matrice en bas à gauche. Les hachures vertes indiquent l'implantation de l'UPC et les zones de maquis ; l'orange la zone UC. Les flèches bleues et oranges montrent les résultats des plébiscites de 1961.}
\end{popfigure}

\textbf{Fond de carte attendu (5 pts)} : contour du Cameroun ; ligne de partage entre le territoire sous tutelle française (Est, grande partie) et les Cameroun britanniques (deux bandes occidentales le long du Nigeria) ; capitales Yaoundé (Est) et Buea (Ouest) placées ; flèche du nord et échelle indicative.

\textbf{Figurés attendus (8 pts)} :
\begin{itemize}
\item Zones d'implantation de l'UPC : pays bassa et Sanaga-Maritime (littoral, autour de Douala-Edéa), pays bamiléké (Ouest, autour de Bafoussam) ; zone de l'UC : Nord-Cameroun (Garoua, Maroua).
\item Hachures sur les zones de maquis upéciste (Sanaga-Maritime, pays bamiléké, pays bassa).
\item Ronds pour les villes de négociation : New York (indiquée en encart ou par une flèche « vers l'ONU ») et Foumban.
\item Flèches de 1961 : du Cameroun méridional (Sud) vers la République du Cameroun ; du Cameroun septentrional (Nord) vers le Nigeria.
\end{itemize}

\textbf{Légende organisée (5 pts)} en quatre rubriques :
\begin{itemize}
\item \emph{Contexte} : territoire sous tutelle française ; territoires sous tutelle britannique (Nord et Sud).
\item \emph{Acteurs} : implantation UPC ; implantation UC ; capitales administratives.
\item \emph{Événements} : zones de maquis (1955-1960) ; lieux de négociation (ONU/New York, Foumban 1961).
\item \emph{Résultats} : indépendance du Cameroun oriental (1\ier{} janvier 1960) ; rattachement du Sud au Cameroun (1\ier{} octobre 1961) ; rattachement du Nord au Nigeria (1\ier{} juin 1961).
\end{itemize}

\textbf{Propreté et présentation (2 pts)} : titre en haut (« Le Cameroun sous tutelle, 1946-1961 : espaces de conflictualité et de négociation »), tracés nets, légende claire et non surchargée.

\textbf{Points de vigilance :} tout figuré utilisé sur la carte doit figurer dans la légende (et inversement) ; ne pas confondre les deux flèches de 1961 (Sud vers le Cameroun, Nord vers le Nigeria) ; la légende doit être \emph{organisée en rubriques}, exigence spécifique du croquis historique au Baccalauréat.
\end{corrige}

\newpage

\coursec{Fiche bilan}

\begin{popfigure}
\begin{center}\tcbox[colback=popink,colframe=popink,coltext=white,arc=9pt,boxsep=4pt,left=6pt,right=6pt]{\bfseries\small La marche vers l'indépendance du Cameroun (1945--1961)}\end{center}
\vspace{-2pt}
\begin{tcbraster}[raster columns=3,raster equal height=rows,raster column skip=6pt,raster row skip=6pt,enhanced,blankest]
\begin{tcolorbox}[enhanced,colback=white,colframe=popblue,boxrule=0.9pt,arc=3pt,title={Tutelle de l'ONU},fonttitle=\bfseries\scriptsize,coltitle=white,colbacktitle=popblue,left=4pt,right=4pt,top=2pt,bottom=2pt,boxsep=1pt,toptitle=1.5pt,bottomtitle=1.5pt,halign=left]
\begin{itemize}[nosep,leftmargin=*,itemsep=1pt,font=\scriptsize]
\item Accords de tutelle 1946
\item France et Grande-Bretagne administratrices
\end{itemize}
\end{tcolorbox}
\begin{tcolorbox}[enhanced,colback=white,colframe=poporange,boxrule=0.9pt,arc=3pt,title={Essor du nationalisme},fonttitle=\bfseries\scriptsize,coltitle=white,colbacktitle=poporange,left=4pt,right=4pt,top=2pt,bottom=2pt,boxsep=1pt,toptitle=1.5pt,bottomtitle=1.5pt,halign=left]
\begin{itemize}[nosep,leftmargin=*,itemsep=1pt,font=\scriptsize]
\item Syndicats et élites
\item UPC fondée en 1948
\end{itemize}
\end{tcolorbox}
\begin{tcolorbox}[enhanced,colback=white,colframe=poppink,boxrule=0.9pt,arc=3pt,title={Répression},fonttitle=\bfseries\scriptsize,coltitle=white,colbacktitle=poppink,left=4pt,right=4pt,top=2pt,bottom=2pt,boxsep=1pt,toptitle=1.5pt,bottomtitle=1.5pt,halign=left]
\begin{itemize}[nosep,leftmargin=*,itemsep=1pt,font=\scriptsize]
\item Émeutes de mai 1955
\item UPC interdite (juillet 1955)
\end{itemize}
\end{tcolorbox}
\begin{tcolorbox}[enhanced,colback=white,colframe=popgreen,boxrule=0.9pt,arc=3pt,title={Grandes figures},fonttitle=\bfseries\scriptsize,coltitle=white,colbacktitle=popgreen,left=4pt,right=4pt,top=2pt,bottom=2pt,boxsep=1pt,toptitle=1.5pt,bottomtitle=1.5pt,halign=left]
\begin{itemize}[nosep,leftmargin=*,itemsep=1pt,font=\scriptsize]
\item Um Nyobè
\item Ahidjo, Mbida, Foncha
\end{itemize}
\end{tcolorbox}
\begin{tcolorbox}[enhanced,colback=white,colframe=poppurple,boxrule=0.9pt,arc=3pt,title={Étapes institutionnelles},fonttitle=\bfseries\scriptsize,coltitle=white,colbacktitle=poppurple,left=4pt,right=4pt,top=2pt,bottom=2pt,boxsep=1pt,toptitle=1.5pt,bottomtitle=1.5pt,halign=left]
\begin{itemize}[nosep,leftmargin=*,itemsep=1pt,font=\scriptsize]
\item Loi-cadre 1956
\item Indépendances 1960
\end{itemize}
\end{tcolorbox}
\begin{tcolorbox}[enhanced,colback=white,colframe=popgold,boxrule=0.9pt,arc=3pt,title={Réunification},fonttitle=\bfseries\scriptsize,coltitle=white,colbacktitle=popgold,left=4pt,right=4pt,top=2pt,bottom=2pt,boxsep=1pt,toptitle=1.5pt,bottomtitle=1.5pt,halign=left]
\begin{itemize}[nosep,leftmargin=*,itemsep=1pt,font=\scriptsize]
\item Plébiscites 1961
\item Fédération 1\textsuperscript{er} oct. 1961
\end{itemize}
\end{tcolorbox}
\end{tcbraster}

\legende{Carte mentale de la leçon : six branches essentielles à maîtriser pour le Bac technique.}
\end{popfigure}

\begin{aretenir}[L'essentiel de la leçon : définitions clés]
\begin{itemize}
  \item \cle{Tutelle} : régime international institué par l'ONU en 1946, confiant l'administration du Cameroun à la France (Est) et à la Grande-Bretagne (Ouest), avec obligation de conduire le territoire vers l'autonomie puis l'indépendance.
  \item \cle{Nationalisme} : volonté politique d'un peuple colonisé d'accéder à la souveraineté ; au Cameroun, il passe des syndicats et associations aux partis politiques après 1945.
  \item \cle{UPC} (Union des populations du Cameroun) : parti créé en 1948, revendiquant l'indépendance immédiate et la réunification ; interdit en juillet 1955.
  \item \cle{Loi-cadre} (23 juin 1956) : loi française instaurant le suffrage universel et des institutions autonomes dans les territoires d'outre-mer et sous tutelle.
  \item \cle{Plébiscite} : consultation populaire organisée par l'ONU les 11 et 12 février 1961 dans les deux Cameroun britanniques.
\end{itemize}
\end{aretenir}

\begin{aretenir}[Dates à connaître par cœur]
\begin{center}
\begin{tabularx}{\linewidth}{|l|X|}\hline
\rowcolor{popblueL} \textbf{Date} & \textbf{Événement} \\ \hline
13 décembre 1946 & Accords de tutelle ONU (France et Grande-Bretagne) \\ \hline
10 avril 1948 & Création de l'UPC \\ \hline
Mai 1955 & Émeutes réprimées à Douala, Yaoundé, Ngambé \\ \hline
13 juillet 1955 & Interdiction de l'UPC par la France \\ \hline
23 juin 1956 & Loi-cadre ; 1957 : Assemblée et gouvernement autonomes (Mbida) \\ \hline
3 septembre 1958 & Assassinat de Ruben Um Nyobè \\ \hline
1\textsuperscript{er} janvier 1960 & Indépendance du Cameroun français (Ahidjo) \\ \hline
11--12 février 1961 & Plébiscites dans les Cameroun britanniques \\ \hline
1\textsuperscript{er} juin 1961 & Le Nord britannique rejoint le Nigeria \\ \hline
1\textsuperscript{er} octobre 1961 & Réunification : République fédérale du Cameroun \\ \hline
\end{tabularx}
\end{center}
\end{aretenir}

\begin{methode}[Méthodes express pour le Bac]
\begin{enumerate}
  \item Commentaire de document : identifier nature, auteur, date, contexte ; dégager l'idée principale ; confronter au cours.
  \item Sujet de synthèse : problématique, plan en deux ou trois parties chronologiques, dates précises et acteurs nommés.
  \item Toujours distinguer les deux espaces : Cameroun français (indépendance 1960) et Cameroun britannique (plébiscite 1961).
\end{enumerate}
\end{methode}

\begin{aretenir}[Checklist avant le Bac]
\begin{itemize}
  \item $\square$ Je sais définir la tutelle de l'ONU et citer les deux puissances administratrices.
  \item $\square$ Je connais la date des accords de tutelle : 13 décembre 1946.
  \item $\square$ Je sais situer la création de l'UPC (1948) et ses deux revendications : indépendance immédiate et réunification.
  \item $\square$ Je connais les grandes figures : Ruben Um Nyobè, Félix-Roland Moumié, Ahmadou Ahidjo, André-Marie Mbida, John Ngu Foncha.
  \item $\square$ Je sais expliquer la répression de mai 1955 et l'interdiction de l'UPC (13 juillet 1955).
  \item $\square$ Je connais le rôle de la loi-cadre du 23 juin 1956 dans le cheminement institutionnel.
  \item $\square$ Je sais que le Cameroun français devient indépendant le 1\textsuperscript{er} janvier 1960 avec Ahidjo.
  \item $\square$ Je connais les plébiscites des 11 et 12 février 1961 et leurs résultats différenciés (Nord vers le Nigeria, Sud vers le Cameroun).
  \item $\square$ Je sais dater la réunification : 1\textsuperscript{er} octobre 1961, naissance de la République fédérale.
  \item $\square$ Je sais construire un commentaire de document en quatre étapes et rédiger une conclusion justifiée.
  \item $\square$ Je sais relier la lutte nationaliste camerounaise au contexte mondial de décolonisation après 1945.
  \item $\square$ Je vérifie l'orthographe des noms propres et la précision des dates dans ma copie.
\end{itemize}
\end{aretenir}

\findecours

\end{document}
